\chapter[Graded unitarity through the BRST coset]{Graded unitarity through the BRST coset: from affine $\mathfrak{sl}_2$ at level $-4/3$ to Lee--Yang}\label{ch:L28}
\begingroup\setlength{\emergencystretch}{3em}\sloppy

\providecommand{\ltwentyeightsl}{\mathfrak{sl}_2}
\providecommand{\ltwentyeightvac}{\mathbf 1}
\providecommand{\ltwentyeightgr}{\operatorname{gr}}
\providecommand{\ltwentyeightGU}{\mathrm{GU}}
\providecommand{\ltwentyeightQ}{\mathbf Q}
\providecommand{\ltwentyeightLY}{L(-\tfrac{22}{5},0)}
\providecommand{\ltwentyeightA}{L_{-4/3}}
\providecommand{\ltwentyeightinv}{\operatorname{inv}}

This chapter was written against an earlier state of the LeanVOA library, and every Lean quotation below was
copied from the source files of that state. For this edition, the book's automatic check confirms that every Lean quotation below occurs verbatim (up to whitespace and comments) in the source files of this release. The author
of this chapter did not run \texttt{lake build}; the aggregate-build claims are records of the project's earlier builds, and the
\texttt{\#print axioms} results that were re-run are listed in Section~\ref{l28:sec:verification}.

The chapter covers the Lean files that the book assigns to it (about 225 files and 53\,700 lines):
the files \leanpath{VOA/\allowbreak Examples/\allowbreak Sl2SemiInfinite*},
\leanpath{Sl2Coset*}, \leanpath{GradedUnitaryCohomology*}, \leanpath{BGPackage*}, \leanpath{BRST*},
\leanpath{GhostConj*}, \leanpath{HarmonicMetric*}, \leanpath{ConcreteGood*}, \leanpath{ConcreteHodge*},
\leanpath{KoszulConcrete*}, \leanpath{KoszulResidual*}, \leanpath{LeadingKoszul*}, \leanpath{InvariantSymbol*},
\leanpath{CyclicEval*}, \leanpath{SymbolModel*}, \leanpath{JetFermionic*}, \leanpath{JetProduct*},
\leanpath{JetStrict*}, \leanpath{Adamovic*}, \leanpath{LYGram*}, \leanpath{RRLower*}, \leanpath{LevelOneHeis*},
\leanpath{LevelOneUnitary*},
\leanpath{LeeYangCharacter*}, \leanpath{LeeYangCosetTransfer*},
\leanpath{LeeYangShortcut*} and \leanpath{LeeYangTransferAssembly*}. We present about forty declarations one by one,
those on which the headline theorem rests or which are results in their own right; the rest
(intermediate lemmas, bookkeeping of gradings on tensor products, mode computations) are described collectively in
Section~\ref{l28:sec:infra}.

\paragraph{Conventions for quotations and status boxes.}
\begin{itemize}
\item \emph{Omitted proofs.} A Lean quotation shows a statement, not its proof. A block that ends in \lean{:=},
  \lean{:= by} or \lean{where} with nothing after it is a statement or definition whose proof or body is
  omitted; it is not a truncated quotation. Documentation comments (\lean{/-- ... -/}) inside quoted
  structures are omitted; the verbatim check ignores comments.
\item \emph{``Axioms standard''.} This means that \texttt{\#print axioms} reports exactly
  \lean{[propext, Classical.choice, Quot.sound]}, the three standard axioms of Lean's logic; see
  Chapter~\ref{ch:01}, Section~\ref{ch01:sec:trust}. It does \emph{not} mean ``unconditional'':
  \texttt{\#print axioms} does not list the explicit hypotheses of a theorem (Section~\ref{ch01:sec:uncond}). Every
  status box names the hypotheses separately. The PBW theorem \lean{VOA.pbw} of
  Chapter~\ref{ch:L02} is now a theorem, not an axiom (it was a project axiom in earlier versions of the library), so the library has no project axiom; some project notes written
  before that change list \lean{VOA.pbw} among the axioms of the headline below, which is no longer the case.
\item \emph{Where a status claim comes from.} ``Re-checked'' means that \texttt{\#print axioms} was re-run on a scratch file
  against the library's build products; the release check of this edition re-ran it for every declaration so marked
  (Section~\ref{l28:sec:verification}). Aggregate-build claims are records of the project's earlier builds, also listed there.
\end{itemize}

\paragraph{Reading the Lean.} Generic Lean idioms (structures and their fields, \lean{Submodule},
\lean{Module.finrank}, coercions $\uparrow$, quotients \lean{⧸}, the order of \lean{∀} and \lean{∃}, \lean{def}
versus \lean{theorem}, why a \lean{Prop}-valued \lean{def} is a \emph{statement} and not a claim) are explained in
Chapter~\ref{ch:01}, Section~\ref{ch01:sec:lean}. One idiom recurs so often in this chapter that we recall it here:
a declaration \lean{def P : Prop := ...} only \emph{names} a proposition; a later \lean{theorem p : P} proves it; and a
theorem that takes \lean{(h : P)} as an argument is \emph{conditional} on \lean{P}. The library uses this pattern to
record its open inputs, and many of them were later proved; when that happened the definition's docstring often
still says ``OPEN''. Section~\ref{l28:sec:defects} lists those stale docstrings.

\paragraph{Prerequisites.} The following are used here and developed elsewhere in the book.
(1) The graded-unitarity predicate of Arabi Ardehali, Beem, Lemos and Rastelli (\lean{GradedUnitary},
\lean{GradedUnitaryAt}, \lean{LayerNondegenerate}, \lean{LayerPositive}), sesquilinear forms, \lean{IsPosOn},
\lean{signedForm} and real structures (\lean{RealStructure}, \lean{gradedSesq}): Chapter~\ref{ch:L19},
Sections~\ref{l19:sec:sesq}--\ref{l19:sec:gu}, Definitions~\ref{l19:def:real}, \ref{l19:def:gradedsesq},
\ref{l19:def:gu}. (2) Graded unitarity of the Virasoro VOAs with the length flag (\lean{VirGradedUnitary},
\lean{virGradedFlag}), and the Lee--Yang model: Definitions~\ref{l19:def:lengthflag} and \ref{l19:def:virgu},
Section~\ref{l19:sec:leeyang}. The \emph{graded isometric embeddings} of the transfer theorems are
Definition~\ref{l19:def:gie}. (3) The universal and simple affine VOAs $V^k(\ltwentyeightsl)$ and
$L_k(\ltwentyeightsl)$, their currents, Sugawara vector and canonical invariant form: Chapter~\ref{ch:L08}
(Definition~\ref{l08:def:Vk}) and Chapter~\ref{ch:L09} (Definition~\ref{l09:def:Lk}). (4) Super vertex algebras,
tensor products and BRST complexes (\lean{SuperVertexAlgebra}, \lean{TensorVA.tensorSVA}, \lean{BRST.BRSTComplex} and
its cohomology vertex algebra): Chapter~\ref{ch:L04}, Sections~\ref{l04:sec:tensor} and \ref{l04:sec:brst}, and
Chapter~\ref{ch:L17} for the $bc$ ghost system (Definition~\ref{l17:def:bc}). (5) Graded flags and good filtrations:
Chapter~\ref{ch:L11}, Definition~\ref{l11:def:flag}. (6) The universal and simple Virasoro VOAs $V(c,0)$, $L(c,0)$,
creation words and the canonical form \lean{Bsimp}: Chapter~\ref{ch:L06}. (7) The vacuum character of
$L_{-4/3}(\ltwentyeightsl)$ (Arakawa--Dai): Chapter~\ref{ch:L20}, Theorem~\ref{l20:thm:character}.

\paragraph{What this chapter is about, in one paragraph.} The Lee--Yang model $M(2,5)$, whose vacuum VOA is the
simple Virasoro VOA $\ltwentyeightLY$, is not unitary in the usual sense: already $\|L_{-2}\ltwentyeightvac\|^2=c/2=-11/5$.
Arabi Ardehali, Beem, Lemos and Rastelli (ABLR) conjectured that it is \emph{graded unitary}: with respect to the
filtration by the number of stress-tensor modes, the sign of the norm alternates with the filtration degree
(Chapter~\ref{ch:L19}). This conjecture is open. The library proves an \emph{implication}: graded unitarity of the
simple affine VOA $L_{-4/3}(\ltwentyeightsl)$ (for its canonical form, its compact conjugation and its ``corrected''
filtration, in which the currents and the Sugawara vector all have degree one) implies graded unitarity of
Lee--Yang. The bridge is the Goddard--Kent--Olive coset $L_{-4/3}\otimes L_1\supset L_{-1/3}\otimes(\text{Virasoro})$,
whose Virasoro part has $c=-22/5$, realized as the relative semi-infinite (BRST) cohomology of
$L_{-4/3}\otimes L_1\otimes V^{-11/3}$ with three pairs of ghosts, together with a theorem of Beem and Garner
that BRST cohomology inherits graded unitarity, and a dimension count that identifies the cohomology with
Lee--Yang. The headline is \lean{VOA.AdamovicSymbol.leeYang_gradedUnitary_of_affine}. Its only hypothesis is the
affine graded unitarity, which is \textbf{open}. The headline is therefore a proved implication with the explicit hypothesis
\lean{hA}, \texttt{\#print axioms}-clean conditional on \lean{hA}; nothing in this chapter proves that Lee--Yang is graded unitary.

\paragraph{Why formalize this.} The paper proof of the implication cites four substantial results: a coset
identification (Arakawa--Creutzig--Linshaw), the Beem--Garner theorem, an ``invariant-symbol'' dimension formula
for $\ltwentyeightgr L_{-4/3}$, and the Rogers--Ramanujan form of the Lee--Yang character. In the Lean development
every one of them is either proved or avoided. The coset identification is avoided altogether
(Section~\ref{l28:sec:coset}); the Beem--Garner positivity argument is proved as finite-dimensional linear algebra
and instantiated on the actual complex (Sections~\ref{l28:sec:hodge} and \ref{l28:sec:package}); the
invariant-symbol bound is proved through a jet-ring model and a $q$-series identity
(Sections~\ref{l28:sec:koszul} and \ref{l28:sec:jet}); and the Lee--Yang character is proved from the Kac
determinant (Section~\ref{l28:sec:character}). Several of these are results of independent interest: the Lee--Yang
vacuum character, unitarity of $L_1(\ltwentyeightsl)$, graded unitarity of $V^k(\ltwentyeightsl)$ for every real
$k<-2$, and a fermionic $q$-series identity for a jet ring.

\begin{remark}[The logical arc and the ledger]\label{l28:rem:arc}
Write $A=L_{-4/3}(\ltwentyeightsl)$, $B=L_1(\ltwentyeightsl)$, $U=V^{-11/3}(\ltwentyeightsl)$ and $L=\ltwentyeightLY$.
Let $C_N$ be the relative BRST complex of $A\otimes B\otimes U\otimes(\text{ghosts})$ at conformal weight $N$,
$H^0_N$ its cohomology at ghost number zero, $F_rH^0_N$ the classes with a representative of $R$-degree $\le r$,
and $F_rL_N$ the span of Virasoro words of length $\le r$ in $L_N$. Fix $N$ and $r$. The proof is the chain
\begin{align*}
\dim F_rL_N
&\ \le\ \dim\bigl(\text{image of }F_rV(-\tfrac{22}{5},0)_N\text{ in }H^0_N\bigr)
  &&\text{(a) Virasoro image, Section~\ref{l28:sec:coset}}\\
&\ \le\ \dim F_rH^0_N
  &&\text{(b) the image is filtered, Theorem~\ref{l28:thm:filtered}}\\
&\ \le\ \textstyle\sum_{s\le r}\dim H^0(Q^-)_{s,N}
  &&\text{(c) harmonic bound, Section~\ref{l28:sec:hodge}}\\
&\ =\ \textstyle\sum_{s\le r}\dim\bigl(\ltwentyeightgr_sA\otimes B\bigr)^{\ltwentyeightsl[t]}_N
  &&\text{(d) Koszul step, Section~\ref{l28:sec:koszul}}\\
&\ \le\ \textstyle\sum_{s\le r}G_s(N)
  &&\text{(e) invariant-symbol bound, Sections~\ref{l28:sec:koszul}--\ref{l28:sec:jet}}\\
&\ =\ \dim F_rL_N
  &&\text{(f) Lee--Yang PBW character, Theorem~\ref{l28:thm:pbwchar}.}
\end{align*}
Here $G_s(N)$ is the number of partitions of $N$ into $s$ parts, each $\ge2$, with consecutive parts differing by at
least $2$. Since the two ends agree, every inequality is an equality. That forces the Virasoro image to be a
\emph{filtered isomorphism} onto $H^0_N$ (Section~\ref{l28:sec:squeeze}); by Beem--Garner, $H^0$ carries a
nondegenerate flag with sign-alternating layers (Section~\ref{l28:sec:hodge}), and the isomorphism carries this to
the length flag of $L$. The hypothesis $\ltwentyeightGU(A)$ is the only hypothesis. Its main use is to make the modified Hermitian form
on the complex positive definite, which is what (c) and the sign statement need; it is also used to define the $R$-grading of the complex
(Definition~\ref{l28:def:good}), so every concrete statement about that grading carries it.

The ledger. ``Proved'' means proved in Lean with the three standard axioms and no hypothesis beyond those displayed.
\begin{center}\small
\begin{longtable}{p{0.36\textwidth}p{0.32\textwidth}p{0.22\textwidth}}
\toprule
step & Lean & status\\
\midrule\endhead
$\ltwentyeightGU(A)$ (the hypothesis) & \lean{hA : GradedUnitary affineForm affineRealStructure affineFlag} & \textbf{open}; numerically checked through weight 7\\
$\ltwentyeightGU(U)$, $U=V^{-11/3}$ & \lean{auxiliaryGradedUnitary} & proved\\
$B=L_1$ is unitary & \lean{LevelOneUnitary.bUnitary} & proved\\
$Q^2=0$ at total level $-4$ & \lean{semiInfiniteQ_sq_zero_iff}, \lean{leeYangBRSTComplex} & proved\\
coset VOA of $c=-22/5$ & \lean{Sl2CosetConcrete.cosetVOA} & proved\\
Beem--Garner positivity, abstract & \lean{KahlerComplex}, \lean{GUData.layer_signed_pos} & proved\\
Beem--Garner package, concrete (goodness, Kähler identities, ghost conjugation, harmonic metric) & \lean{relativeDifferentialGood_of_affine}, \lean{ghostConjRhoD_proved}, \lean{ghostConjBE_proved}, \lean{rep*_proved} & proved from $\ltwentyeightGU(A)$ and $\ltwentyeightGU(U)$\\
(a) Virasoro image & \lean{virasoroToH0_ker_le_maxIdeal} & proved\\
(b) the image is filtered & \lean{endpointKahler_filtered} & proved from $\ltwentyeightGU(A)$\\
(c) harmonic bound & \lean{finrank_cohomFilt_row_le_lead} & proved\\
(d) Koszul step & \lean{leadCohomEquivInvariants}, \lean{kExact_endpoint} & proved\\
(e) invariant-symbol bound & \lean{finrank_le_gapTwoCoefficient_of_cyclicEvaluation}, \lean{Desc.wordSymbolDescends}, \lean{jetFermionicIdentity} & proved\\
(f) Lee--Yang character & \lean{finrank_pieceIn_le}, \lean{leeYangRRLowerBound_holds}, \lean{leeYangPBWCharacter_}\allowbreak\lean{iff_lowerBound} & proved\\
squeeze and transfer & \lean{gradedUnitaryAt_of_class_family} & proved\\
headline & \lean{AdamovicSymbol.leeYang_gradedUnitary_of_affine} & proved, \textbf{conditional on} $\ltwentyeightGU(A)$\\
\bottomrule
\end{longtable}
\end{center}
\end{remark}

\paragraph{The findings of this chapter, in brief.}
\begin{enumerate}
\item \textbf{The headline is a conditional theorem with exactly one hypothesis,} and that hypothesis is the intended
  one: graded unitarity of the concrete simple affine VOA $L_{-4/3}(\ltwentyeightsl)$ with its normalized canonical
  form, its compact conjugation and its corrected flag, all three \emph{constructed}, none of them a free parameter
  (Theorem~\ref{l28:thm:headline}). The conclusion is also pinned: the Lee--Yang VOA, its canonical form, the
  partition-basis real structure and the length flag. \texttt{\#print axioms} is the standard three (re-checked).
\item \textbf{No coset identification is used.} The library constructs the coset VOA and its central charge, but the
  proof never identifies it with Lee--Yang. Instead it maps the universal Virasoro VOA into the cohomology through the
  coset stress tensor, and the dimension count does the rest (Section~\ref{l28:sec:coset}).
\item \textbf{Proved along the way, without hypotheses:} the Lee--Yang vacuum character (Rocha-Caridi form and
  Rogers--Ramanujan form, Theorem~\ref{l28:thm:lychar}); unitarity of $L_1(\ltwentyeightsl)$
  (Theorem~\ref{l28:thm:bunitary}); $Q^2=0$ exactly at total level $-4$, given the structural hypotheses displayed there
(Theorem~\ref{l28:thm:qsq}); the fermionic
  jet identity $\sum_rq^r\sum_k\left[\begin{smallmatrix}2r\\k\end{smallmatrix}\right]_q/(q)_r=((q^3;q^3)_\infty/(q;q)_\infty)^3$
  (Theorem~\ref{l28:thm:jetfermionic}).
\item \textbf{A name clash.} Two different theorems are called \lean{leeYang_gradedUnitary_of_affine}. The one in
  namespace \lean{VOA.LeeYangTransferAssembly} is a superseded conditional assembly with eight further named inputs,
  among them the open coset identification; it is not the headline (Warning~\ref{l28:warn:clash}).
\item \textbf{A reduction ladder.} About fifty theorems named \lean{leeYang_gradedUnitary_of_affine_*} record the
  successive stages of the proof, each discharging inputs of the previous one. They are described collectively in
  Section~\ref{l28:sec:ladder}; a few rest on hypotheses that are probably unsatisfiable, and are dead ends
  (Warning~\ref{l28:warn:vacuous}).
\item \textbf{One direction only.} The library proves the implication from $L_{-4/3}(\ltwentyeightsl)$ to Lee--Yang; this
  edition states and claims no converse implication.
\item \textbf{Stale docstrings.} About twenty \lean{Prop}-valued definitions on the proof path still carry the
  docstring ``OPEN'' although they are proved (Section~\ref{l28:sec:defects}). No statement-level defect was found.
\end{enumerate}

\paragraph{Notation used in this chapter.}
\begin{center}\small
\begin{longtable}{p{0.24\textwidth}p{0.68\textwidth}}
\toprule
symbol & meaning\\
\midrule\endhead
$e,f,h$; colours $0,1,2$ & the Chevalley basis of $\ltwentyeightsl$; in Lean the colours \lean{0, 1, 2} of \lean{Fin 3} are $e,f,h$.
  The normalized form is $\kappa(e,f)=1$, $\kappa(h,h)=2$ (Lean \lean{kappa}); \lean{fabc a b c} are the structure constants
  $[x_a,x_b]=\sum_cf_{ab}{}^cx_c$.\\
$J^a_n$, $a_{(n)}$ & the mode $n$ of the current $a$; for a weight-one state $a$, $a_{(n)}=J^a_n$. In Lean
  \lean{mode (V.Y u) n} is $u_{(n)}$ (Chapter~\ref{ch:01}).\\
$k$; $A$, $B$, $U$ & the level; $A=L_{-4/3}(\ltwentyeightsl)$ (Lean \lean{AVOA}), $B=L_1(\ltwentyeightsl)$ (\lean{BVOA}),
  $U=V^{-11/3}(\ltwentyeightsl)$ (\lean{UVOA}).\\
$L$, $c$ & the Lee--Yang VOA $\ltwentyeightLY$, $c=c_{2,5}=-22/5$ (Lean \lean{leeYangC}).\\
$b^a,c^a$; $\mathrm{gh}$ & the three ghost pairs (antighost $b$, weight $1$, ghost number $-1$; ghost $c$, weight $0$,
  ghost number $+1$); the ghost number.\\
$\ltwentyeightQ$, $Q$ & the BRST differential (Lean \lean{semiInfiniteQ}); in Section~\ref{l28:sec:hodge} the bold
  $\ltwentyeightQ^-=Q^-+S^-$ and $\ltwentyeightQ^+=Q^++S^+$ are Beem--Garner's pair of differentials and their
  homogeneous pieces.\\
$R$; $2R$ & the $R$-filtration degree; Lean indexes it by $2R\in\Z$ to allow half-integers.\\
$\ltwentyeightGU(X)$ & graded unitarity of $X$ for its stated form, conjugation and flag
  (Definition~\ref{l19:def:gu}).\\
$F_r$; $\ltwentyeightgr_r$ & the step $r$ of a filtration; the associated graded piece $F_r/F_{r-1}$.\\
$\ltwentyeightsl[t]$ & the positive current algebra spanned by $J^a_n$, $n\ge0$.\\
$\C[u_j,v_j]$; $X,Y,Z,J$ & the jet ring of two formal series $u(s)=\sum_ju_js^j$, $v(s)=\sum_jv_js^j$; $X=u^2$, $Y=uv$,
  $Z=v^2$, $J=u'v-uv'$ (as series in $s$).\\
$A_r$ & the span of products of $r$ coefficients of $X,Y,Z,J$ (Lean \lean{jetImage r}).\\
$(q)_n$, $(a;q)_\infty$ & $\prod_{i=1}^n(1-q^i)$; $\prod_{i\ge0}(1-aq^i)$. $\left[\begin{smallmatrix}n\\k\end{smallmatrix}\right]_q$
  the Gaussian binomial.\\
$p(N)$; $G_r(N)$; $\mathrm{rc}(N)$ & partitions of $N$; partitions of $N$ into $r$ parts $\ge2$ with gaps $\ge2$ (Lean
  \lean{gapTwoCoefficient N r}); $\sum_k(-1)^kp(N-k(5k+3)/2)$ (Lean \lean{rc N}).\\
\bottomrule
\end{longtable}
\end{center}

%======================================================================
\section{The mathematics in outline}\label{l28:sec:math}
%======================================================================

\subsection{Graded unitarity, recalled}

Let $V$ be a VOA with an invariant bilinear form $B$, an anti-linear conjugation $\rho$ commuting with all vertex
operators, and a \emph{graded flag} $F$: for each conformal weight $N$ an increasing exhaustive sequence of subspaces
$0=F_{-1}V_N\subseteq F_0V_N\subseteq F_1V_N\subseteq\cdots$, with $F_0V_0=\C\ltwentyeightvac$ and $F_0V_N=0$ for $N>0$.
Put $h(x,y)=B(\rho x,y)$, a sesquilinear form. The \emph{layer} $\mathcal L_p V_N$ is the $B$-orthogonal complement of
$F_{p-1}V_N$ inside $F_pV_N$. $V$ is \emph{graded unitary} (Definition~\ref{l19:def:gu}) when, at every weight,
(iv) $B$ is nondegenerate on every $F_pV_N$, and (vi) $(-1)^ph$ is positive definite on every layer
$\mathcal L_pV_N$. For the Virasoro VOA $L(c,0)$ the flag is pinned: $F_p$ is spanned by words
$L_{-n_1}\cdots L_{-n_q}\ltwentyeightvac$ with $q\le p$ (Definition~\ref{l19:def:lengthflag}).

\begin{example}[Lee--Yang at low weight]\label{l28:ex:lylow}
At $c=-22/5$ the weights $0,\dots,12$ of $L$ have dimensions $1,0,1,1,1,1,2,2,3,3,4,4,6$: these are the numbers of
partitions of $N$ into parts $\ge2$ with consecutive parts differing by at least $2$ (for $N=12$: $12$, $10+2$, $9+3$,
$8+4$, $7+5$, $6+4+2$). At weight $2$ the only state is $L_{-2}\ltwentyeightvac$, of length one, with
$h(L_{-2}\ltwentyeightvac,L_{-2}\ltwentyeightvac)=c/2=-11/5<0$: clause (vi) asks for $(-1)^1h>0$, which holds. At
weight $4$ the universal module has the two words $L_{-4}\ltwentyeightvac$, $L_{-2}^2\ltwentyeightvac$, but $L_4$ is
one-dimensional: the Lee--Yang singular vector $\bigl(L_{-2}^2-\tfrac35L_{-4}\bigr)\ltwentyeightvac$ vanishes in $L$
(Theorem~\ref{l19:thm:lylevel4}). So $F_1L_4=F_2L_4=L_4$, the layer at $p=2$ is zero, and the only condition is
$h(L_{-4}\ltwentyeightvac,L_{-4}\ltwentyeightvac)=\tfrac{4^3-4}{12}c=5c=-22<0$, again the right sign. At weight $6$,
$F_1L_6=\C L_{-6}\ltwentyeightvac$ and $F_2L_6=L_6$ is two-dimensional; the length-two layer is the first place where
the sign is a genuine constraint. Chapter~\ref{ch:L19} proves the length $\le2$ part of graded unitarity for every
boundary model; the full statement for Lee--Yang is what this chapter reduces to $\ltwentyeightGU(A)$.
\end{example}

\subsection{The coset}

For $k\notin\{-2,-3\}$ the diagonal currents $J^a\otimes1+1\otimes J^a$ generate an affine algebra of level $k+1$
inside $L_k(\ltwentyeightsl)\otimes L_1(\ltwentyeightsl)$, and the difference of Sugawara vectors
$\omega_{\mathrm{coset}}=\omega_k\otimes1+1\otimes\omega_1-\omega^{\mathrm{diag}}_{k+1}$ commutes with the diagonal
currents. Its central charge is
\begin{equation*}
c=\frac{3k}{k+2}+1-\frac{3(k+1)}{k+3}=1-\frac{6}{(k+2)(k+3)}.
\end{equation*}
At $k=-4/3$: $\tfrac{3k}{k+2}=-6$, the diagonal level is $-1/3$ with Sugawara charge $\tfrac{3(-1/3)}{5/3}=-\tfrac35$, and
$c=-6+1+\tfrac35=-\tfrac{22}5$. This is the Goddard--Kent--Olive (GKO) construction, continued to a negative level. The
\emph{coset VOA} is the commutant: all states killed by the non-negative diagonal current modes. Arakawa, Creutzig and
Linshaw proved that this commutant is the simple Lee--Yang VOA. The library does not use that theorem.

\subsection{Relative semi-infinite cohomology}

The commutant can be computed as a BRST cohomology. Take three pairs of fermionic ghosts $b^a,c^a$ in the adjoint of
$\ltwentyeightsl$; they carry an affine $\ltwentyeightsl$ action $J^a_{\mathrm{gh}}$ at level $2h^\vee=4$. Take a matter
VOA $M$ with an affine $\ltwentyeightsl$ action at level $k_M$, and the BRST current
\begin{equation*}
j_{\mathrm{BRST}}=\sum_a J^a_M\,c^a+\tfrac12\sum_a{:}c^aJ^a_{\mathrm{gh}}{:},\qquad
\ltwentyeightQ=\oint j_{\mathrm{BRST}}.
\end{equation*}
Then $\ltwentyeightQ^2=0$ if and only if $k_M+4=0$, i.e.\ $k_M=-2h^\vee$ (Theorem~\ref{l28:thm:qsq}). The \emph{relative}
complex is the subspace killed by the zero modes of the total currents $J^a_M+J^a_{\mathrm{gh}}$ and of the antighosts
$b^a$. For the matter $M=A\otimes B\otimes U$ with $A=L_{-4/3}$, $B=L_1$ and the auxiliary $U=V^{-11/3}$, the diagonal
level is $-\tfrac43+1-\tfrac{11}3=-4$. Physically the auxiliary factor ``absorbs'' the diagonal $\ltwentyeightsl$: the
relative cohomology of $(A\otimes B)\otimes V^{-4-(k+1)}$ at ghost number zero is, for good matter, the commutant of
the diagonal currents in $A\otimes B$. Here $-4-(-1/3)=-11/3$.

\subsection{Beem--Garner: cohomology inherits graded unitarity}

Beem and Garner (arXiv:2509.10364) prove that if the matter is graded unitary and the BRST action is \emph{good} (the
differential has no component of $R$-degree zero), then the relative cohomology is graded unitary for the induced
filtration. The mechanism is Hodge theory: there is a positive definite \emph{modified} Hermitian form on the complex,
built from the graded-unitarity forms of the factors with the sign $i^{2R}$; the differential splits as
$\ltwentyeightQ^-=Q^-+S^-$ with $Q^-$ raising $R$ by $\tfrac12$ and $S^-$ lowering it; a second differential
$\ltwentyeightQ^+$ is related to the adjoint of $\ltwentyeightQ^-$ by an $\mathrm{Sp}(2)$ (``Kähler'') identity; and
every cohomology class has a unique harmonic representative, which is \emph{homogeneous} in $R$. Homogeneity is what
makes the induced flag on cohomology split into layers, and positivity of the modified form on harmonic vectors of
$R$-degree $p$ is exactly $(-1)^ph>0$. Section~\ref{l28:sec:hodge} formalizes this at one conformal weight, and
Section~\ref{l28:sec:package} constructs all of it on the actual complex.

\subsection{The squeeze}

Graded unitarity of the cohomology does not yet say anything about Lee--Yang. The last step compares dimensions. The
harmonic vectors of $R$-degree $s$ inject into the cohomology $H(Q^-)$ of the \emph{leading} differential
(Theorem~\ref{l28:thm:harmbound}), and a Koszul argument identifies that leading cohomology with the
$\ltwentyeightsl[t]$-invariants in $\ltwentyeightgr_sA\otimes B$ (Section~\ref{l28:sec:koszul}). The associated graded
$\ltwentyeightgr A$ is modelled by a ring of jet polynomials (Section~\ref{l28:sec:jet}), and the invariants are
counted from above by $G_s(N)$. From below, the universal Virasoro algebra maps into the cohomology through the coset
stress tensor, and its image surjects onto $L$ filtration step by filtration step. The Lee--Yang character
(Section~\ref{l28:sec:character}) says that $\dim F_rL_N=\sum_{s\le r}G_s(N)$. The two bounds meet.

\begin{physics}
\textbf{Four-dimensional origin.} The Lee--Yang VOA is the chiral algebra, in the sense of Beem, Lemos, Liendo, Peelaers,
Rastelli and van Rees, of the simplest Argyres--Douglas theory $(A_1,A_2)$; $L_{-4/3}(\ltwentyeightsl)$ is the chiral algebra
of the $(A_1,A_3)=(A_1,D_3)$ Argyres--Douglas theory, whose $SU(2)$ flavour symmetry has two-dimensional level
$k_{2d}=-\tfrac43$. ABLR showed that unitarity of the four-dimensional theory forces graded unitarity of its chiral algebra,
with the filtration coming from the $SU(2)_R$ charge. So both $\ltwentyeightGU(A)$ and $\ltwentyeightGU(L)$ are predicted by
physics, and neither is proved. The BRST construction with ghosts at level $+4$ is the chiral-algebra image of
\emph{gauging} an $SU(2)$ flavour symmetry, and the condition $k_M=-4=-2h^\vee$ is the vanishing of the gauge coupling's beta
function; Beem--Garner is the chiral-algebra shadow of the fact that gauging preserves unitarity. The factors $B$ and $U$
in our complex are mathematical devices, not chiral algebras of known four-dimensional theories, and this physical
picture is a heuristic only.
\end{physics}

%======================================================================
\section{The two endpoint algebras and the hypothesis}\label{l28:sec:endpoints}
%======================================================================

The headline theorem is about concrete objects. This section defines them: the three affine factors, the corrected
filtration of $A$, the compact conjugation, and the Lee--Yang side.

\begin{definition}[The three affine factors]\label{l28:def:factors}
$A=L_{-4/3}(\ltwentyeightsl)$ is the simple quotient of $V^{-4/3}(\ltwentyeightsl)$ by its maximal graded ideal;
$B=L_1(\ltwentyeightsl)$ likewise at level $1$; $U=V^{-11/3}(\ltwentyeightsl)$ is the \emph{universal} affine VOA at level
$-11/3$. The level $-4/3$ is the first member $k_1$ of the boundary family $k_m=-2+\frac{2}{2m+1}$.
\end{definition}

\noindent File \leanpath{VOA/\allowbreak Examples/\allowbreak Sl2CosetConcreteFactors.lean}.
\begin{leancode}
abbrev ALevel : ℂ := AffineSl2BoundarySingularFamily.level 1

theorem aLevel_eq : ALevel = (-4 / 3 : ℂ) :=
  AffineSl2BoundarySingularFamily.level_one

theorem aNoncritical : ALevel + 2 ≠ 0 :=
  AffineSl2BoundarySingularFamily.level_noncritical 1

theorem bNoncritical : (1 : ℂ) + 2 ≠ 0 := by norm_num

theorem uNoncritical : (-11 / 3 : ℂ) + 2 ≠ 0 := by norm_num

abbrev AVOA := AffineSl2GeneralRigidity.simple ALevel aNoncritical

abbrev BVOA := AffineSl2GeneralRigidity.simple 1 bNoncritical

abbrev UVOA := AffineSl2GeneralRigidity.universal (-11 / 3) uNoncritical
\end{leancode}

\begin{dictionary}
\lean{AffineSl2BoundarySingularFamily.level 1} is $-2+2/(2\cdot1+1)=-4/3$, and \lean{aLevel_eq} records it.
\lean{AffineSl2GeneralRigidity.simple k hk} is the quotient of the universal affine VOA by its maximal graded ideal
\lean{QuotVOA.maxGradedIdeal}, i.e.\ $L_k(\ltwentyeightsl)$ of Definition~\ref{l09:def:Lk}; \lean{universal k hk} is
$V^k(\ltwentyeightsl)$ (Definition~\ref{l08:def:Vk}). The hypotheses \lean{hk : k + 2 ≠ 0} exclude the critical level,
where the Sugawara construction fails; the three proofs \lean{aNoncritical}, \lean{bNoncritical}, \lean{uNoncritical} supply
them. An \lean{abbrev} is a definition that Lean unfolds freely (Chapter~\ref{ch:01}, Section~\ref{ch01:sec:synonyms}).
Exactly the definition.

\textbf{Why $U$ is universal.} The auxiliary factor is never quotiented: its role is to supply the level $-11/3$ and a
graded-unitary structure, which is proved for the universal algebra (Theorem~\ref{l28:thm:aux}). Whether $V^{-11/3}$ is
simple plays no role.
\end{dictionary}

\begin{physics}
$A$ is the chiral algebra of the $(A_1,D_3)$ Argyres--Douglas theory, $B$ is the $SU(2)$ level-one WZW model (one compact
boson at the self-dual radius), and $U$ is an auxiliary current algebra whose level is fixed by the anomaly cancellation
$-\tfrac43+1+k_U=-4$; it has no four-dimensional interpretation that the library uses.
\end{physics}

\begin{status}
Definitions; \lean{aLevel_eq} and the three non-criticality facts are proved. Axioms standard (re-checked for the
objects built on them, Section~\ref{l28:sec:verification}).
\end{status}

\begin{definition}[The corrected filtration]\label{l28:def:corrected}
A \emph{corrected word} is a finite list of letters $(g,m)$, where $g$ is a current $e,f,h$ or the Sugawara vector $\omega$,
and $m\in\N$ is a mode depth. The letter $(g,m)$ acts by $g_{(-m-1)}$: for a current this is $J^g_{-m-1}$, for $\omega$ it is
$L_{-m-2}$. The \emph{corrected filtration} $F_pV^k$ is the span of the states $g^1_{(-m_1-1)}\cdots g^\ell_{(-m_\ell-1)}\ltwentyeightvac$
with $\ell\le p$ letters, intersected with each conformal weight space. On a quotient $V^k/I$ by a graded ideal the flag is the
image. The flag on $A$ is the image of the corrected flag of $V^{-4/3}$.
\end{definition}

\noindent File \leanpath{VOA/\allowbreak Examples/\allowbreak AffineSl2GoodFiltration.lean} (a file of Chapter~\ref{ch:L08}).
\begin{leancode}
inductive Generator
  | current (a : Fin 3)
  | sugawara
  deriving DecidableEq

def generatorState (k : ℂ) (hk : k + 2 ≠ 0) : Generator → WakModule k
  | .current a => currentState k a
  | .sugawara => (affineSl2VOA_closed k hk).conformalVector

def correctedWordState (k : ℂ) (hk : k + 2 ≠ 0) :
    List (Generator × ℕ) → WakModule k
  | [] => (affineSl2VOA_closed k hk).vacuum
  | gr :: rest =>
      (affineSl2VOA_closed k hk).smul (generatorState k hk gr.1)
        (-(gr.2 : ℤ) - 1) (correctedWordState k hk rest)

def correctedAffineRPiece (k : ℂ) (hk : k + 2 ≠ 0) (p : ℕ) :
    Submodule ℂ (WakModule k) :=
  Submodule.span ℂ (Set.range (fun w : {w : List (Generator × ℕ) // w.length ≤ p} =>
    correctedWordState k hk w.1))

noncomputable def correctedAffineRQuotientFlag (k : ℂ) (hk : k + 2 ≠ 0)
    (I : Submodule ℂ (WakModule k))
    (hI : (affineSl2VOA_closed k hk).IsIdeal I)
    (hgr : QuotVOA.IsGradedIdeal (affineSl2VOA_closed k hk) I) :
    GradedFlag (QuotVOA.quotientVOA (affineSl2VOA_closed k hk) I hI hgr) :=
  (correctedAffineRFlag k hk).quotient I hI hgr

\end{leancode}
\noindent File \leanpath{VOA/\allowbreak Examples/\allowbreak LeeYangTransferAssemblyAffine.lean}.
\begin{leancode}
def affineFlag : GradedFlag AVOA :=
  CorrectedR.correctedAffineRQuotientFlag ALevel aNoncritical
    (QuotVOA.maxGradedIdeal (universal ALevel aNoncritical))
    (QuotVOA.maxGradedIdeal_isIdeal _)
    (QuotVOA.maxGradedIdeal_isGradedIdeal _)
\end{leancode}

\begin{dictionary}
\textbf{Decoding.} \lean{WakModule k} is the carrier of $V^k(\ltwentyeightsl)$ and \lean{affineSl2VOA_closed k hk} its VOA
structure (Chapter~\ref{ch:L08}). \lean{(affineSl2VOA_closed k hk).smul u n v} is the $n$-th product $u_{(n)}v$, so
\lean{correctedWordState k hk ((g, m) :: rest)} is $g_{(-m-1)}$ applied to the word \lean{rest}. For a current $a$ (weight one)
$a_{(-m-1)}=J^a_{-m-1}$; for $\omega$ (weight two) $\omega_{(-m-1)}=L_{-m-2}$. \lean{currentState k a} is $J^a_{-1}\ltwentyeightvac$.
\lean{correctedAffineRPiece k hk p} is the span of all words with \lean{w.length ≤ p} letters, of any weight; the graded
flag \lean{correctedAffineRFlag} (not quoted) intersects it with each conformal weight space, and \lean{.quotient I hI hgr}
takes images in the quotient VOA by a graded ideal $I$. \lean{affineFlag} applies this to the maximal graded ideal, so it is a
flag on $A$.

\textbf{The cost of a letter is one, whatever its mode.} A derivative $J^a_{-2}\ltwentyeightvac=(\partial J^a)$ has degree one,
and so does $L_{-2}\ltwentyeightvac=\omega$, although in terms of currents $\omega$ is quadratic. This is the
``correction'' relative to the ordinary current-PBW filtration (where $\omega$ would have degree two). It is what ABLR's
$R$-filtration predicts: the stress tensor sits in the $SU(2)_R$ current multiplet, of $R$-charge one. The library proves that
this flag is a \emph{good} filtration (modes add degrees, \lean{correctedAffineRQuotientFlag_modeCompatible}); it does not
prove that it coincides with any intrinsically characterized filtration of $A$, and it does not need to.

\textbf{How it could fail to be the intended object.} The flag is an image under the quotient map; no basis of $A$ is assumed.
Had the flag been allowed to vary, the hypothesis ``$\ltwentyeightGU(A)$ for some flag'' would be nearly empty
(Warning~\ref{l19:warn:freeflag}); here it is pinned.
\end{dictionary}

\begin{physics}
$F_p$ counts the number of $SU(2)_R$-charged ``letters'' (flavour moment maps $J^a$ and the stress tensor $T$), each of
$R$-charge one, needed to build a Schur operator; derivatives cost nothing because they do not change $R$.
\end{physics}

\begin{status}
Definitions with all \lean{GradedFlag} fields proved (Chapter~\ref{ch:L08}, Section~\ref{l08:sec:infrastructure}); axioms
standard. Mode compatibility of the quotient flag is proved.
\end{status}

\begin{definition}[Compact conjugation]\label{l28:def:compact}
For real $k\ne-2$, the \emph{compact conjugation} of $V^k(\ltwentyeightsl)$ is the unique anti-linear map $\rho$ that fixes
$\ltwentyeightvac$, commutes with vertex operators in the sense $\rho(u_{(n)}v)=\rho(u)_{(n)}\rho(v)$, and acts on the currents by
$\rho(e)=-f$, $\rho(f)=-e$, $\rho(h)=-h$; that is, $J^a\mapsto-J^{a^*}$ with $a^*$ the ``adjoint colour''. It is an involution,
and it descends to every quotient by an ideal that it preserves; on $A$ we use the descended map.
\end{definition}

\noindent File \leanpath{VOA/\allowbreak Examples/\allowbreak AffineSl2BelowCriticalCompact.lean}.
\begin{leancode}
structure IsCompactRealStructure (R : RealStructure (universal k hk)) : Prop where
  current : ∀ a : Fin 3, R.rho (currentState k a) = -currentState k (adjointColour a)
  parity_ne_zero : ∀ n : ℤ, R.parity n ≠ 0
\end{leancode}
\noindent File \leanpath{VOA/\allowbreak Examples/\allowbreak AffineSl2BelowCriticalRealStructure.lean}.
\begin{leancode}
def compactRealStructure : RealStructure (universal k hk) where
\end{leancode}
\noindent File \leanpath{VOA/\allowbreak Examples/\allowbreak LeeYangTransferAssemblyAffine.lean}.
\begin{leancode}
def affineRealStructure : RealStructure AVOA where

theorem affineRealStructure_current (a : Fin 3) :
    affineRealStructure.rho (simpleCurrent ALevel aNoncritical a) =
      -simpleCurrent ALevel aNoncritical (adjointColour a) := by
\end{leancode}

\begin{dictionary}
\lean{adjointColour} swaps the colours $0\leftrightarrow1$ ($e\leftrightarrow f$) and fixes $2$ ($h$)
(\leanpath{AffineSl2VermaDual.lean}). \lean{IsCompactRealStructure R} says that the real structure $R$ (Definition~\ref{l19:def:real})
acts on the current states as $J^a\mapsto-J^{a^*}$ and that its parity scalars are nonzero. \lean{compactRealStructure} is
\emph{constructed} (the \lean{where} introduces its fields, whose proofs we omit) for real $k$ (\lean{hreal : starRingEnd ℂ k = k},
a section variable) and has parity $\equiv1$, i.e.\ $\rho^2=1$. \lean{affineRealStructure} is its descent to $A$: the file proves
that $\rho$ preserves the maximal ideal (the preimage of the maximal ideal is a proper ideal, so it lies in the maximal ideal) and
builds the quotient map; \lean{affineRealStructure_current} confirms that on $A$ it still sends the current $a$ to $-$(current $a^*$).
The parity of \lean{affineRealStructure} is \lean{fun _ => 1}.

\textbf{Why this conjugation.} $x\mapsto-x^{*}$, with $e^*=f$, $h^*=h$, is the Cartan involution whose fixed points form the compact
real form $\mathfrak{su}(2)$. With the invariance $a_n^\dagger=-a_{-n}$ of a weight-one primary under the VOA invariant form, it
gives the Hilbert-space adjoint $(J^a_n)^\dagger=J^{a^*}_{-n}$ of a unitary current algebra. Any other choice (for instance the split
real form) would make the hypothesis a different, and physically unmotivated, statement.
\end{dictionary}

\begin{physics}
$\rho$ is the CPT-like anti-unitary map of the $SU(2)$ flavour current algebra: it sends a flavour current to its Hermitian
conjugate. The compact real form is the one under which a unitary four-dimensional theory's flavour symmetry acts unitarily.
\end{physics}

\begin{status}
Constructed, all fields proved; axioms standard. Unconditional (given real $k$).
\end{status}

\begin{example}[Weight one of $A$]\label{l28:ex:weightone}
The weight-one space of $A$ is spanned by the three currents, all of corrected degree one, so $F_0A_1=0$, $F_1A_1=A_1$, and the
only layer is $p=1$. With the normalized invariant form, $B(J^a_{-1}\ltwentyeightvac,J^b_{-1}\ltwentyeightvac)=-k\,\kappa(a,b)$ (the sign
from $a_1^\dagger=-a_{-1}$ for weight-one primaries), so
\begin{equation*}
h(e,e)=B(\rho e,e)=B(-f,e)=k\,\kappa(f,e)=k=-\tfrac43,\qquad h(h,h)=B(-h,h)=2k=-\tfrac83 .
\end{equation*}
Clause (vi) at $p=1$ asks $-h>0$: it holds precisely because $k<0$.
At level $k>0$ the sign would be wrong, which is
why graded unitarity is a negative-level phenomenon and ordinary unitarity a positive-level one.
\end{example}

\begin{definition}[The Lee--Yang datum]\label{l28:def:lydatum}
On the Lee--Yang side the form is the canonical normalized form \lean{Bsimp} of $L(c,0)$, the flag is the length flag, and the real
structure is the partition-basis conjugation of Theorem~\ref{l19:thm:virreal}, at $c=c_{2,5}=1-12\lambda^2$ with the real background
charge $\lambda=\lambda_2$ of Theorem~\ref{l19:thm:boundaryreal}.
\end{definition}

\noindent File \leanpath{VOA/\allowbreak Examples/\allowbreak LeeYangTransferAssemblyAffine.lean}.
\begin{leancode}
abbrev affineForm : InvariantForm AVOA := simpleForm ALevel aNoncritical

def leeYangRealStructure : RealStructure (VirSimple.virasoroSimple LeeYang.leeYangC) :=
  virasoroSimpleRealStructure
    ((VirGradedUnitarity.boundaryBackgroundCharge_spec 2).trans
      VirGradedUnitarity.boundaryCharge_two_eq_leeYangC)
    (VirGradedUnitarity.boundaryBackgroundCharge_fixed 2)
\end{leancode}

\begin{dictionary}
\lean{simpleForm k hk} is the canonical invariant bilinear form of $L_k(\ltwentyeightsl)$ with $(\ltwentyeightvac|\ltwentyeightvac)=1$
(Chapter~\ref{ch:L09}); \lean{affineForm} is it at $k=-4/3$. \lean{leeYangRealStructure} is \lean{virasoroSimpleRealStructure}
(Theorem~\ref{l19:thm:virreal}) applied to the proof that $1-12\lambda_2^2$ equals the boundary charge of index two, which equals
\lean{leeYangC}$=-22/5$ (\lean{boundaryCharge_two_eq_leeYangC}), and to the reality of $\lambda_2$. So it fixes every Virasoro word
and conjugates coefficients. \lean{VirGradedUnitary c R} is \lean{GradedUnitary (Bsimp c) R (virGradedFlag c)}
(Definition~\ref{l19:def:virgu}): form and flag are pinned, only $R$ is a parameter, and here it is this specific $R$.
\end{dictionary}

\begin{physics}
The Lee--Yang stress tensor is a Hermitian field with real OPE coefficients, so conjugating coefficients in the basis of Virasoro words
is its anti-unitary conjugation. The statement $\ltwentyeightGU(L)$ with these data is ABLR's prediction for the $(A_1,A_2)$
Argyres--Douglas theory, whose Schur operators are the Virasoro descendants of the identity.
\end{physics}

\begin{status}
Definitions; axioms standard. These are exactly the data of Kaufmann's and ABLR's statement for $M(2,5)$.
\end{status}

%======================================================================
\section{The headline theorem}\label{l28:sec:headline}
%======================================================================

\begin{theorem}[Affine graded unitarity implies Lee--Yang graded unitarity]\label{l28:thm:headline}
Suppose that $A=L_{-4/3}(\ltwentyeightsl)$, with its normalized canonical invariant form, the descended compact conjugation
(Definition~\ref{l28:def:compact}) and the corrected flag (Definition~\ref{l28:def:corrected}), is graded unitary. Then the
Lee--Yang VOA $\ltwentyeightLY$, with its normalized canonical form, its partition-basis conjugation and its length flag
(Definition~\ref{l28:def:lydatum}), is graded unitary.
\end{theorem}

\noindent File \leanpath{VOA/\allowbreak Examples/\allowbreak AdamovicDescCyclic.lean}.
\begin{leancode}
theorem leeYang_gradedUnitary_of_affine
    (hA : GradedUnitary affineForm affineRealStructure affineFlag) :
    VirGradedUnitarity.VirGradedUnitary LeeYang.leeYangC leeYangRealStructure :=
  leeYang_gradedUnitary_of_affine_wordSymbols_jet hA fun r => Desc.wordSymbolDescends (r + 2)
\end{leancode}

\begin{dictionary}
\textbf{Every symbol.} \lean{GradedUnitary B R F} is Definition~\ref{l19:def:gu}: for every integer weight $n$, clause (iv)
(nondegeneracy of $B$ on every flag step) and clause (vi) (positivity of $(-1)^ph$ on every layer). \lean{affineForm},
\lean{affineRealStructure}, \lean{affineFlag} are the three data of Definitions~\ref{l28:def:lydatum},
\ref{l28:def:compact}, \ref{l28:def:corrected} on \lean{AVOA}$=A$. \lean{VirGradedUnitary LeeYang.leeYangC leeYangRealStructure}
unfolds to \lean{GradedUnitary (Bsimp (-22/5)) leeYangRealStructure (virGradedFlag (-22/5))} (Definition~\ref{l19:def:virgu}).
The namespace \lean{VirGradedUnitarity} is where \lean{VirGradedUnitary} lives; \lean{LeeYang.leeYangC} is the number $-22/5$.

\textbf{The hypotheses.} There is exactly one, \lean{hA}. It is the intended statement: graded unitarity of the \emph{simple}
quotient $A$ (not of the universal $V^{-4/3}$, which has a singular vector at weight $3$, Chapter~\ref{ch:L08}), for the
compact conjugation and the corrected (current-and-stress) flag. No other datum is free: the form, conjugation and flag are
the constructed objects of Section~\ref{l28:sec:endpoints}.

\textbf{Same, stronger or weaker.} The same as the paper statement of the implication, with one difference in the \emph{proof}: the paper proof cites an
external coset identification and an external invariant-symbol theorem; the Lean theorem uses neither as a hypothesis.

\textbf{The one-line proof term.} The proof is a composition: the rung \lean{leeYang_gradedUnitary_of_affine_wordSymbols_jet}
(Section~\ref{l28:sec:ladder}), which still takes the word-symbol descent \lean{desc} as an input, applied to the proof
\lean{Desc.wordSymbolDescends} of that input (Theorem~\ref{l28:thm:desc}).

\textbf{How it could have been vacuous, and why it is not.} (1) If \lean{hA} were false the theorem would be true and empty.
Nobody knows whether \lean{hA} is true; it is the four-dimensional prediction, and the project's exact-arithmetic Gram
computations (outside the library) find it satisfied through conformal weight $7$
(numerical evidence, not a proof). Example~\ref{l28:ex:weightone} checks weight one by hand. (2) The conclusion is not
trivially true: it fails if the form or real structure is changed (for example, the sign twist of
Warning~\ref{l28:warn:signtwist}), and only its length-$\le2$ part is proved independently (Chapter~\ref{ch:L19}). (3) The
flag of the hypothesis is pinned, so Warning~\ref{l19:warn:freeflag} does not apply.
\end{dictionary}

\paragraph{Proof on paper.} Fix a conformal weight $N$. The relative BRST complex $C_N$ of $A\otimes B\otimes U\otimes(\text{ghosts})$
is a finite-dimensional space (Section~\ref{l28:sec:brst}). The hypothesis, with the proved unitarity of $B$ and graded unitarity of
$U$, makes the Beem--Garner modified Hermitian form on $C_N$ positive definite, and the BRST differential is good for the
$R$-grading (Section~\ref{l28:sec:package}). The finite-dimensional Hodge theory of Section~\ref{l28:sec:hodge} then gives: every
class in $H^0_N$ has a unique harmonic representative, homogeneous in $R$; on the harmonic model, the induced flag is nondegenerate
for the invariant form and its layers have the sign $(-1)^p$; and $\dim F_rH^0_N\le\sum_{s\le r}\dim H^0(Q^-)_{s,N}$. The Koszul step
(Section~\ref{l28:sec:koszul}) identifies $H^0(Q^-)_{s,N}$ with the $\ltwentyeightsl[t]$-invariants of $\ltwentyeightgr_sA\otimes B$
in weight $N$, and the jet model (Section~\ref{l28:sec:jet}) bounds their dimension by $G_s(N)$. On the other side, the universal
Virasoro VOA maps into $H^0$ through the coset stress tensor (Section~\ref{l28:sec:coset}); its kernel is a proper ideal, hence lies
in the maximal ideal, so the map factors through a surjection from (a subspace of) the image onto $L$ which respects the length
filtration. Lee--Yang's PBW character (Section~\ref{l28:sec:character}) gives $\dim F_rL_N=\sum_{s\le r}G_s(N)$. The chain of
Remark~\ref{l28:rem:arc} closes, forcing the image of $F_rV(c,0)_N$ to be all of $F_rH^0_N$ and the map to $L$ to be injective
on it. The invariant form and the conjugation on harmonic representatives agree with \lean{Bsimp} and the Lee--Yang conjugation
(Theorem~\ref{l28:thm:repmetric}). Transporting the nondegenerate sign-alternating flag of $H^0_N$ to $L_N$ gives clauses (iv)
and (vi) at weight $N$ (Theorem~\ref{l28:thm:classfamily}). \qed

\begin{physics}
Read physically: if the $SU(2)$ flavour current algebra of the $(A_1,D_3)$ Argyres--Douglas theory obeys the positivity constraints
that four-dimensional unitarity predicts, then so does the Virasoro algebra of the $(A_1,A_2)$ theory. The bridge is a
chiral-algebra version of a four-dimensional operation (gauging an $SU(2)$ with a level-one current algebra and an auxiliary sector),
not a known four-dimensional duality between the two theories; the implication should be read as a statement about chiral algebras,
with the four-dimensional story as motivation.
\end{physics}

\begin{status}
Proved; \textbf{conditional on \lean{hA}}, which is open. \texttt{\#print axioms} \lean{VOA.AdamovicSymbol.leeYang_gradedUnitary_of_affine}:
\lean{[propext, Classical.choice, Quot.sound]} (re-checked), so it is \texttt{\#print axioms}-clean, conditional on \lean{hA}; the theorem
is also in the library's maintained axiom audit. A second, independent derivation of the same implication, through the
general $(2,2n+3)$ machinery, is \lean{VOA.TwoSeries.leeYang_gradedUnitary_of_affine_viaTwoSeries} (Chapter~\ref{ch:L29}); it has
the same single hypothesis and the standard axioms (re-checked).
\end{status}

\begin{warning}[title={Name clash: two different theorems called \texttt{leeYang\_gradedUnitary\_of\_affine}}]\label{l28:warn:clash}
The short name \lean{leeYang_gradedUnitary_of_affine} denotes two different theorems. The headline is
\lean{VOA.AdamovicSymbol.leeYang_gradedUnitary_of_affine}. The other, \lean{VOA.LeeYangTransferAssembly.leeYang_gradedUnitary_of_affine},
is the first conditional assembly of the argument, now superseded:
\par\noindent File \leanpath{VOA/\allowbreak Examples/\allowbreak LeeYangTransferAssembly.lean}.
\begin{leancode}
theorem leeYang_gradedUnitary_of_affine
    (hA : GradedUnitary affineForm affineRealStructure affineFlag)
    (inputs : OpenInputs) : VirGradedUnitary leeYangC leeYangRealStructure := by
\end{leancode}
\noindent with \leanpath{VOA/\allowbreak Examples/\allowbreak LeeYangTransferAssemblyInputs.lean}:
\begin{leancode}
structure OpenInputs where
  identification : CosetIdentification
  kahler : ConcreteKahlerGeometry
  weakGU : WeakGUTensorPositivity kahler
  gu : ConcreteGUData kahler
  metric : HarmonicMetricAgreement identification kahler gu
  filtered : CosetFiltrationCompatibility identification kahler
  symbols : LeadingInvariantSymbolDimension kahler
  character : LeeYangPBWCharacter
\end{leancode}
\noindent and \leanpath{VOA/\allowbreak Examples/\allowbreak LeeYangTransferAssemblySimple.lean}:
\begin{leancode}
structure CosetIdentification where
  equiv : CosetSpace ≃ₗ[ℂ] LeeYangSpace
  hom : IsVOAHom cosetVOA (virasoroSimple leeYangC) equiv.toLinearMap
  conformal : equiv cosetVOA.conformalVector =
    (virasoroSimple leeYangC).conformalVector
\end{leancode}
It takes, besides \lean{hA}, a structure \lean{inputs : OpenInputs} with eight fields, the first of which is a VOA isomorphism
between the coset VOA and Lee--Yang (\lean{CosetIdentification}); that identification is \emph{not} constructed anywhere in the
library. The superseded theorem is a correct conditional statement, but it must not be cited as the result. Always use full names.
(A documentation defect, severity S4 on the scale of Section~\ref{l28:sec:defects}.)
\end{warning}

\subsection{The reduction ladder, described collectively}\label{l28:sec:ladder}

The proof was built top-down: a first theorem with many named inputs, then a sequence of theorems each discharging some inputs
and sometimes introducing sharper ones. Every stage was kept as a theorem. The library therefore contains about fifty theorems
named \lean{leeYang_gradedUnitary_of_affine_*}, all with the conclusion \lean{VirGradedUnitary leeYangC leeYangRealStructure} and
the first hypothesis \lean{hA}, in the files of this chapter, plus two in the files of Chapter~\ref{ch:L29}
(\lean{..._viaTwoSeries}, \lean{..._viaInvariants}). They are not separate results: each is a step toward the headline. We list the
ones that mark the main stages, in the order in which the headline depends on them (each rung is proved from the next one down);
the reader needs this table only to decode the library, not to understand the mathematics.
\begin{center}\small
\begin{longtable}{p{0.40\textwidth}p{0.52\textwidth}}
\toprule
rung (namespace, suffix) & inputs besides \lean{hA}\\
\midrule\endhead
\lean{AdamovicSymbol}, (none) & none: the headline\\
\lean{AdamovicSymbol}, \lean{_wordSymbols_jet} & \lean{desc : ∀ r, WordSymbolDescends (r + 2)}\\
\lean{AdamovicSymbol}, \lean{_jetFermionic} & \lean{desc}, \lean{hq : JetFermionicLower}\\
\lean{AdamovicSymbol}, \lean{_wordSymbols_final} & \lean{desc}, \lean{chi : JetProductBound}\\
\lean{AdamovicSymbol}, \lean{_wordSymbols_gram}, \lean{_wordSymbols_rc}, \lean{_wordSymbols_two} & \lean{desc}, \lean{chi}, and a Lee--Yang character input (rank bound, Rocha-Caridi bound, or Rogers--Ramanujan bound)\\
\lean{KoszulResidual}, \lean{_leadingSymbol} & \lean{sym : ∀ r, JetLeadingSymbol (r + 1)}, \lean{chi}, \lean{lower}\\
\lean{KoszulResidual}, \lean{_strict} & \lean{strict : ∀ r, GrJetStrict (r + 1)}, \lean{lower}\\
\lean{KoszulResidual}, \lean{_symbolResidual}; \lean{KoszulConcrete}, \lean{_koszulResidual} & \lean{strict}, the symbol maps or the full Koszul residual, \lean{lower}\\
\lean{SymbolModel}, \lean{_jetStrict}, \lean{_grLeadKoszul}, \lean{_grSymbols}, \lean{_grStrict} & strictness or cyclic vectors, the leading Koszul data, \lean{lower}\\
\lean{LevelOneHeis}, \lean{_symbolModels}; \lean{CyclicEval}, \lean{_heisenberg*}, \lean{_symbolCyclic} & symbol models, level-one facts, \lean{lower}\\
\lean{ConcreteHodge}, \lean{_cyclicEvaluation}, \lean{_levelOne_rr_bound}, \lean{_levelOne_rr} & an invariant-symbol bound or equality, \lean{lower}\\
\lean{ConcreteHodge}, \lean{_levelOne} & \lean{symbols}, \lean{character}\\
\lean{ConcreteHodge}, \lean{_harmonic}, \lean{_reps}, \lean{_bgConj}, \lean{_ghostConj} & \lean{hB} (unitarity of $L_1$), \lean{symbols}, \lean{character}, and metric or ghost-conjugation identities\\
\lean{ConcreteHodge}, \lean{_bgRemainderI}, \lean{_bgPackageI}, \lean{_sp2*}, \lean{_concrete}, \ldots & further Beem--Garner package data\\
\lean{LeeYangShortcut}, \lean{_of_affine'}, \lean{_of_imageMetric} & a \lean{ShortcutInputs} structure (no coset identification)\\
\lean{LeeYangTransferAssembly}, (none) & \lean{OpenInputs}, including the coset identification (Warning~\ref{l28:warn:clash})\\
\bottomrule
\end{longtable}
\end{center}
Each input name in the table is a \lean{Prop} or a structure defined in the files of this chapter; the ones that the headline
ultimately needs are presented, and proved, in the sections that follow.

\begin{warning}[title={Probably vacuous rungs}]\label{l28:warn:vacuous}
Some intermediate rungs take inputs that are most likely unsatisfiable on the actual complex, so they are true but empty, and the
headline does not use them.
\begin{itemize}
\item \lean{VOA.ConcreteHodge.leeYang_gradedUnitary_of_affine_kahler} (\leanpath{BRSTGoodnessCreation.lean}) and
  \lean{..._ghost} (\leanpath{ConcreteGoodGenerators.lean}) take \lean{hK : InterpretedKahlerAnticomm (ghostRGrading w) ...}, a
  Kähler identity for a ghost form that is the same on all three colours. Invariance under $\ltwentyeightsl$ needs the colour
  $h$ to carry half the norm of $e$ and $f$ (because $\kappa(h,h)=2\kappa(e,f)$), so \lean{hK} is most likely false. This is a
  hand argument; it is not machine-checked.
\item \lean{..._bgRemainder} (\leanpath{BGPackageBFactor.lean}) and \lean{..._bgPackage} take a ghost-conjugation datum
  \lean{BGGhostRemainder} whose skewness constant $-(-1)^{\mathrm{gh}}$, together with a scalar $\rho^2$, admits solutions only at a
  value $\kappa=-i$ of a normalization constant that must be real; numerical evidence (project computations outside the library) only.
  The corrected datum (\lean{BGGhostRemainderI}, with $\rho^2=1$ and $B(Eu,v)=i(-1)^{\mathrm{gh}}B(u,Ev)$) is the one used, and it is
  proved (Theorem~\ref{l28:thm:ghostconj}).
\end{itemize}
As statements these are severity-S1 defects (wrong or vacuous; scale of Section~\ref{l28:sec:defects}); they are harmless for the headline, which does not use them.
\end{warning}

%======================================================================
\section{The relative semi-infinite complex}\label{l28:sec:brst}
%======================================================================

\begin{definition}[Affine $\ltwentyeightsl$ action on a super vertex algebra]\label{l28:def:action}
Let $V$ be a super vertex algebra. An \emph{affine $\ltwentyeightsl$ action of level $k$} on $V$ is a triple of even states $J^a$
of weight one ($a\in\{e,f,h\}$) whose modes satisfy
\begin{equation*}
[J^a_m,J^b_n]=\sum_cf_{ab}{}^c\,J^c_{m+n}+m\,k\,\kappa(a,b)\,\delta_{m+n,0}\qquad(m,n\in\Z).
\end{equation*}
A \emph{translated} action carries in addition a translation operator $T$ with $T\ltwentyeightvac=0$ and $[T,Y(u,z)]=\partial_zY(u,z)$.
\end{definition}

\noindent File \leanpath{VOA/\allowbreak Examples/\allowbreak Sl2SemiInfiniteAction.lean}.
\begin{leancode}
structure AffineSl2Action (V : SuperVertexAlgebra M) (k : ℂ) where
  current : Fin 3 → M
  current_weight : ∀ a, current a ∈ V.grading 1
  current_even : ∀ a, current a ∈ V.parity 0
  current_comm : ∀ (a b : Fin 3) (m n : ℤ),
    mode (V.Y (current a)) m * mode (V.Y (current b)) n
      - mode (V.Y (current b)) n * mode (V.Y (current a)) m
      = (∑ c, fabc a b c • mode (V.Y (current c)) (m + n))
        + if m + n = 0 then ((m : ℂ) * k * kappa a b) • LinearMap.id else 0
\end{leancode}
\noindent Files \leanpath{VOA/\allowbreak Examples/\allowbreak Sl2SemiInfiniteTranslation.lean}, \leanpath{Sl2SemiInfiniteDifferential.lean}.
\begin{leancode}
structure TranslationCovariant (V : SuperVertexAlgebra 𝒱) where
  T : End ℂ 𝒱
  vacuum_zero : T V.vacuum = 0
  covariant : ∀ u, Reconstruction.DCovariant T (V.Y u)

structure TranslatedAffineSl2Action (V : SuperVertexAlgebra M) (k : ℂ)
    extends AffineSl2Action V k where
  translation : TranslationCovariant V
\end{leancode}

\begin{dictionary}
\lean{mode (V.Y (current a)) m} is the mode $(J^a)_{(m)}$, which for a weight-one state is $J^a_m$. The commutator is stated for all
integer modes (not a truncated OPE); \lean{if m + n = 0 then ... else 0} is the Kronecker delta, and \lean{LinearMap.id} the identity
operator. \lean{V.parity 0} is the even part. \lean{Reconstruction.DCovariant T (V.Y u)} is $[T,Y(u,z)]=\partial_zY(u,z)$ in modes.
\lean{extends} adds the field \lean{translation} to \lean{AffineSl2Action}. Translation is required separately because a library
\lean{SuperVertexAlgebra} has no translation axiom (Chapter~\ref{ch:L17}, Theorem~\ref{l17:thm:nocov}). The universal affine VOA
inhabits the structure at every non-critical level (\lean{translatedAffineAction}), so it is not vacuous.
\end{dictionary}

\begin{physics}
A matter sector with an $SU(2)$ current algebra of level $k$: the data one needs in order to gauge the $SU(2)$.
\end{physics}

\begin{status}
Definitions. Instances for the universal affine VOA, for tensor products (\lean{tensorTranslatedAffineAction}) and for the simple
quotients are constructed and used below.
\end{status}

\begin{definition}[Ghosts, BRST current and BRST charge]\label{l28:def:charge}
The ghost system is the tensor product of three $bc$ systems (Definition~\ref{l17:def:bc}), one per colour $a$; $b^a$ has weight
$1$ and ghost number $-1$, $c^a$ weight $0$ and ghost number $+1$. The ghost currents are the adjoint bilinears $J^a_{\mathrm{gh}}=\sum_{b,c}f_{ab}{}^{c}\,{:}b^{c}c^{b}{:}$. For a matter action $J^a_M$ on $V$, the BRST current on $V\otimes(\text{ghosts})$ is
$j=\sum_aJ^a_M\otimes c^a+\ltwentyeightvac\otimes\tfrac12\sum_a{:}c^aJ^a_{\mathrm{gh}}{:}$ and the BRST charge is its residue
$\ltwentyeightQ=j_{(0)}$.
\end{definition}

\noindent Files \leanpath{VOA/\allowbreak Examples/\allowbreak Sl2SemiInfiniteGhosts.lean}, \leanpath{Sl2SemiInfiniteCurrents.lean},
\leanpath{Sl2SemiInfiniteCharge.lean}.
\begin{leancode}
abbrev GhostPairSpace := SympFerm.FockSF ⊗[ℂ] SympFerm.FockSF

abbrev Sl2GhostSpace := GhostPairSpace ⊗[ℂ] SympFerm.FockSF

noncomputable def bGhost (a : Fin 3) : Sl2GhostSpace :=
  ghostEmbed a (SympFerm.stWordSF [((0, 0) : ℕ × Fin 2)])

noncomputable def cGhost (a : Fin 3) : Sl2GhostSpace :=
  ghostEmbed a (SympFerm.stWordSF [((0, 1) : ℕ × Fin 2)])

noncomputable def ghostCurrent (a : Fin 3) : Sl2GhostSpace :=
  ∑ b : Fin 3, ∑ c : Fin 3, ghostCurrentCoefficient a b c • ghostBilinear c b

noncomputable def ghostBRSTState : Sl2GhostSpace :=
  (1 / 2 : ℂ) • ∑ a : Fin 3,
    mode (sl2GhostSVA.Y (cGhost a)) (-1) (ghostCurrent a)

noncomputable def brstState (A : AffineSl2Action V k) : M ⊗[ℂ] Sl2GhostSpace :=
  (∑ a : Fin 3, A.current a ⊗ₜ[ℂ] cGhost a) + V.vacuum ⊗ₜ[ℂ] ghostBRSTState

noncomputable def semiInfiniteQ (A : AffineSl2Action V k) : End ℂ (M ⊗[ℂ] Sl2GhostSpace) :=
  mode (TensorVA.tenY V sl2GhostSVA (brstState A)) 0
\end{leancode}

\begin{dictionary}
\lean{SympFerm.FockSF} is the Fock space of one fermionic pair, made into the $bc$ system \lean{Ghost.bcGhostSVA}
(Chapter~\ref{ch:L17}); the one-letter words \lean{stWordSF [((0, 0))]} and \lean{stWordSF [((0, 1))]} are the states of the fields
$b(z)$ (weight $1$) and $c(z)$ (weight $0$). \lean{ghostEmbed a x} puts $x$ in the pair of colour $a$ and vacua elsewhere. \lean{ghostCurrentCoefficient a b c} is the
structure constant \lean{fabc a b c}, and \lean{ghostBilinear c b} is ${:}b^cc^b{:}$, the $(-1)$-product of the two ghost states
(\leanpath{Sl2SemiInfiniteCurrents.lean}).
\lean{mode (Y x) (-1) y} is the normal-ordered product ${:}xy{:}$. \lean{⊗ₜ[ℂ]} is the elementary tensor. \lean{TensorVA.tenY} is the
tensor-product vertex operator with Koszul signs (Definition~\ref{l04:def:tenY}), and \lean{mode (...) 0} is the zeroth product,
i.e.\ the residue $\oint j$. No nilpotency is assumed in the definition: it is a theorem below.

\textbf{Conventions.} The $c^a$ are dual to the Chevalley basis indexing the matter currents, so no inverse Killing form appears
(the file's docstring). The ghost states of the three pairs anticommute by the Koszul signs of the tensor product.
\end{dictionary}

\begin{physics}
These are the Faddeev--Popov ghosts of an $SU(2)$ gauge field in the chiral algebra; $\ltwentyeightQ$ is the BRST charge of the
gauging, $\ltwentyeightQ=\oint\bigl(c^aJ^a_M+\tfrac12c^aJ^a_{\mathrm{gh}}\bigr)$, familiar from string theory, where the same
construction with the Virasoro algebra gives the critical dimension.
\end{physics}

\begin{status}
Definitions. Weight $1$, odd parity and ghost number $+1$ of \lean{brstState}, and the corresponding properties of
\lean{semiInfiniteQ} (weight $0$, odd, ghost number $+1$), are proved in the same file.
\end{status}

\begin{theorem}[Ghost level and nilpotency]\label{l28:thm:qsq}
(1) The ghost currents form an affine $\ltwentyeightsl$ action of level $+4=2h^\vee$ on the ghost system.
(2) For a translated affine action of level $k$ on a matter super vertex algebra $V$ with nonzero vacuum (and the parity and grading
conditions needed to form the tensor product), $\ltwentyeightQ^2=0$ if and only if $k=-4$.
\end{theorem}

\noindent Files \leanpath{VOA/\allowbreak Examples/\allowbreak Sl2SemiInfiniteGhostAffine.lean}, \leanpath{Sl2SemiInfiniteAnomalyCriterion.lean}.
\begin{leancode}
noncomputable def ghostAffineAction : AffineSl2Action sl2GhostSVA 4 :=

theorem semiInfiniteQ_sq_zero_iff (A : AffineSl2Action V k)
    (H : TranslationCovariant V) (hV : TensorVA.ParityCompatible V)
    (hg : TensorVA.TensorGradingHyp V sl2GhostSVA) (hvac : V.vacuum ≠ 0) :
    semiInfiniteQ A * semiInfiniteQ A = 0 ↔ k = -4 := by
\end{leancode}

\begin{dictionary}
\lean{AffineSl2Action sl2GhostSVA 4}: the ghost currents satisfy Definition~\ref{l28:def:action} at level $4$; the construction
\lean{affineActionOfProducts} derives all mode commutators from the explicit singular products. In the second statement, \lean{A} is a
bare affine action (the translation \lean{H} is supplied separately), \lean{hV} and \lean{hg} are the parity compatibility and the
finite-dimensionality of tensor weight spaces needed by \lean{TensorVA.tensorSVA} (Chapter~\ref{ch:L04}), and \lean{hvac} excludes the
zero algebra. \lean{semiInfiniteQ A * semiInfiniteQ A} is the composite $\ltwentyeightQ\circ\ltwentyeightQ$. The \lean{↔} is a genuine
equivalence: at any other level $\ltwentyeightQ^2\neq0$.

\textbf{Why the hypotheses are harmless.} All of them hold for the concrete matter below (\lean{matterParityCompatible},
\lean{matterGhostTensorGrading}, \lean{matterVacuum_ne_zero}).
\end{dictionary}

\paragraph{Proof on paper.} $\ltwentyeightQ^2=\tfrac12\{\ltwentyeightQ,\ltwentyeightQ\}$ is the residue of the self-OPE of $j$. The
$c^ac^b$-terms combine the matter commutator with the ghost cubic term; their simple-pole parts cancel by the Jacobi identity of
$\ltwentyeightsl$ (this is the role of the factor $\tfrac12$), and what remains is the central term. The total current
$J^a_M+J^a_{\mathrm{gh}}$ has level $k+4$, and the anomaly is proportional to it: the library proves the exact formula
$\ltwentyeightQ^2=\tfrac12(k+4)\,(\ltwentyeightvac\otimes\alpha)_{(0)}$ for the ghost state $\alpha=\sum_{a,b}\kappa(a,b)\,{:}\partial c^a\,c^b{:}$
(\lean{ghostAnomalyState}, \leanpath{Sl2SemiInfiniteChargeProducts.lean}), in \lean{semiInfiniteQ_square_anomaly} (\leanpath{Sl2SemiInfiniteNilpotency.lean}). The ghost level $4$ comes from the double
contraction $\sum_{c,d}f_{ac}{}^df_{bd}{}^c=2h^\vee\kappa(a,b)$ with $h^\vee=2$. The operator $(\ltwentyeightvac\otimes\alpha)_{(0)}$ is nonzero when the matter vacuum is nonzero
(\leanpath{Sl2SemiInfiniteAnomalyCriterion.lean}), which gives the converse direction. \qed

\begin{physics}
Gauge anomaly cancellation: the BRST charge of a gauged current algebra is nilpotent exactly when the total matter level is
$-2h^\vee$. In four dimensions this is the condition that the gauge coupling be exactly marginal (vanishing one-loop beta function),
translated by the chiral-algebra map $k_{2d}=-k_{4d}/2$.
\end{physics}

\begin{status}
Both proved, axioms standard (re-checked). Unconditional given the stated structural hypotheses.
\end{status}

\begin{definition}[Relative states]\label{l28:def:relative}
The \emph{relative} subspace of $V\otimes(\text{ghosts})$ is the simultaneous kernel of the six operators $(J^a_M+J^a_{\mathrm{gh}})_0$
and $b^a_0$ ($a=e,f,h$).
\end{definition}

\noindent File \leanpath{VOA/\allowbreak Examples/\allowbreak Sl2SemiInfiniteRelative.lean}.
\begin{leancode}
noncomputable def totalCurrentZero (A : AffineSl2Action V k) (a : Fin 3) :
    End ℂ (M ⊗[ℂ] Sl2GhostSpace) :=
  mode (TensorVA.tenY V sl2GhostSVA (totalCurrentState A a)) 0

noncomputable def antighostZero (V : SuperVertexAlgebra M) (a : Fin 3) :
    End ℂ (M ⊗[ℂ] Sl2GhostSpace) :=
  mode (TensorVA.tenY V sl2GhostSVA (tensorAntighost V a)) 0

noncomputable def relativeCarrier (A : AffineSl2Action V k) :
    Submodule ℂ (M ⊗[ℂ] Sl2GhostSpace) :=
  (⨅ a : Fin 3, LinearMap.ker (totalCurrentZero A a)) ⊓
    ⨅ a : Fin 3, LinearMap.ker (antighostZero V a)
\end{leancode}

\begin{dictionary}
\lean{totalCurrentState A a} is $J^a_M\otimes\ltwentyeightvac+\ltwentyeightvac\otimes J^a_{\mathrm{gh}}$ and \lean{tensorAntighost V a} is
$\ltwentyeightvac\otimes b^a$; \lean{mode (...) 0} is the zero mode of a weight-one field. \lean{⨅} is the intersection of a family of
subspaces and \lean{⊓} the intersection of two. Exactly the definition. The file proves that $\ltwentyeightQ$ preserves the relative
subspace (\lean{semiInfiniteQ_mem_relativeCarrier}) and that the vacuum is relative.
\end{dictionary}

\begin{physics}
Relative cohomology imposes the Gauss law ($\ltwentyeightsl$-invariance under the total zero modes) and removes the ghost zero modes
$c^a_0$, which would otherwise double the cohomology; it is the chiral analogue of computing gauge-invariant states.
\end{physics}

\begin{status}
Definition; stability under $\ltwentyeightQ$ proved (axioms standard).
\end{status}

\begin{theorem}[The concrete Lee--Yang BRST complex]\label{l28:thm:complex}
The matter $(A\otimes B)\otimes U$ carries the diagonal affine action of level $-\tfrac43+1-\tfrac{11}3=-4$; tensoring with the ghosts
gives a BRST complex in the sense of Definition~\ref{l04:def:brst} whose differential is $\ltwentyeightQ$. Its cohomology vertex
superalgebra has nonzero vacuum class.
\end{theorem}

\noindent File \leanpath{VOA/\allowbreak Examples/\allowbreak Sl2CosetConcreteFactors.lean}.
\begin{leancode}
def leeYangMatterAction : TranslatedAffineSl2Action matterSVA (-4) :=
  Sl2SemiInfinite.leeYangMatterAction aAction bAction uAction
    aParityCompatible bParityCompatible uParityCompatible abTensorGrading abuTensorGrading

def leeYangBRSTComplex : BRST.BRSTComplex brstSVA :=
  semiInfiniteBRSTComplex leeYangMatterAction matterParityCompatible matterGhostTensorGrading

theorem leeYangBRSTComplex_Q :
    leeYangBRSTComplex.Q = semiInfiniteQ leeYangMatterAction.toAffineSl2Action := rfl

theorem leeYangBRSTComplex_cohomology_vacuum_ne_zero :
    leeYangBRSTComplex.cohomology.vacuum ≠ 0 :=
\end{leancode}

\begin{dictionary}
\lean{TranslatedAffineSl2Action matterSVA (-4)}: an action of level exactly $-4$ on \lean{matterSVA}, the three-fold tensor product
$(A\otimes B)\otimes U$ (\lean{tensorThreeSVA}). Its currents are $J^a\otimes1\otimes1+1\otimes J^a\otimes1+1\otimes1\otimes J^a$
(\lean{Sl2SemiInfinite.leeYangMatterAction}, \leanpath{Sl2CosetTensor.lean}); the level adds. \lean{brstSVA} is
$(\text{matter})\otimes(\text{ghosts})$. \lean{BRST.BRSTComplex brstSVA} bundles a ghost-number grading, the differential, and the
proofs that it has weight $0$, is odd, raises ghost number by one, squares to zero and is a derivation of all products
(Definition~\ref{l04:def:brst}); \lean{semiInfiniteBRSTComplex} supplies them, using Theorem~\ref{l28:thm:qsq} at $k=-4$.
\lean{leeYangBRSTComplex_Q} confirms that the differential is literally \lean{semiInfiniteQ}. \lean{cohomology.vacuum ≠ 0} says the
vacuum is not exact, so the cohomology is not the zero algebra.

\textbf{No hypotheses.} This is a definition with every field proved for the concrete factors: nothing about the complex is assumed.
\end{dictionary}

\paragraph{Proof on paper.} The three actions are proved for the simple quotients and the universal algebra
(\lean{aAction}, \lean{bAction}, \lean{uAction}); levels add under tensor products; Theorem~\ref{l28:thm:qsq} gives $\ltwentyeightQ^2=0$.
The vacuum is not exact because $\ltwentyeightQ$ preserves weight and ghost number and the weight-zero, ghost-number-$(-1)$ space is
zero (no state of weight $0$ contains a $b$). \qed

\begin{physics}
The chiral algebra of ``$SU(2)$ gauge theory'' coupled to the three matter sectors; its ghost-number-zero cohomology is the
algebra of gauge-invariant chiral operators, and the vacuum survives gauging.
\end{physics}

\begin{status}
Constructed and proved; axioms standard (re-checked: \lean{leeYangBRSTComplex}, \lean{leeYangBRSTComplex_cohomology_vacuum_ne_zero},
\lean{Sl2CosetConcrete.relativeVacuum_survives}). Unconditional.
\end{status}

%======================================================================
\section{The coset VOA and the Virasoro image}\label{l28:sec:coset}
%======================================================================

\begin{definition}[Commutant and coset stress tensor]\label{l28:def:coset}
For an affine action $J^a$ on $V$, the \emph{commutant} is $\{x\in V: J^a_nx=0\ \text{for all }a\text{ and all }n\ge0\}$. On
$A\otimes B$ take the diagonal action $J^a\otimes1+1\otimes J^a$ (level $-1/3$), and the coset stress tensor
$\omega_{\mathrm{coset}}=\omega_A\otimes\ltwentyeightvac+\ltwentyeightvac\otimes\omega_B-\omega^{\mathrm{diag}}$, where $\omega^{\mathrm{diag}}$ is
the Sugawara vector of the diagonal currents.
\end{definition}

\noindent Files \leanpath{VOA/\allowbreak Examples/\allowbreak Sl2CosetSubspace.lean}, \leanpath{Sl2CosetConcreteOmega.lean},
\leanpath{Sl2CosetConcreteCohomology.lean}.
\begin{leancode}
def cosetSubspace (A : AffineSl2Action V k) : Submodule ℂ M where
  carrier := {x | ∀ (a : Fin 3) (n : ℤ), 0 ≤ n → mode (V.Y (A.current a)) n x = 0}

def ambientConformal : ASpace ⊗[ℂ] BSpace :=
  AVOA.conformalVector ⊗ₜ[ℂ] BSVA.vacuum + ASVA.vacuum ⊗ₜ[ℂ] BVOA.conformalVector

def diagonalSugawara : ASpace ⊗[ℂ] BSpace := sugawaraState diagonalAction.toAffineSl2Action

def cosetOmega : ASpace ⊗[ℂ] BSpace := ambientConformal - diagonalSugawara

abbrev CosetSpace := cosetSubspace diagonalAction.toAffineSl2Action
\end{leancode}

\begin{dictionary}
\lean{cosetSubspace A} is a \lean{Submodule} given by its carrier (the remaining fields, closure under $+$ and scalars, are proved);
the carrier is literally the commutant. \lean{diagonalAction} is the tensor-product action on $A\otimes B$ at level $-4/3+1$
(\leanpath{Sl2CosetConcreteFactors.lean}); \lean{sugawaraState} is $\frac{1}{2(k+2)}\sum_{a,b}\kappa^{ab}{:}J^aJ^b{:}$ for that action.
\lean{ASVA.vacuum} and \lean{BSVA.vacuum} are the vacua. Exactly the definition.
\end{dictionary}

\begin{physics}
The commutant is the GKO coset: the operators of the $A\otimes B$ theory that are neutral under the diagonal $SU(2)$ current algebra. The
difference of stress tensors is the stress tensor of that neutral sector.
\end{physics}

\begin{status}
Definitions; that \lean{cosetOmega} lies in the commutant and is a conformal vector there is proved (\lean{cosetConformalData},
\lean{cosetOmega_virasoro}).
\end{status}

\begin{theorem}[The coset VOA]\label{l28:thm:cosetvoa}
$c_A+c_B-c_{\mathrm{diag}}=-6+1-(-\tfrac35)=-\tfrac{22}{5}$, and the commutant, with $\omega_{\mathrm{coset}}$, is a VOA of central charge
$-22/5$.
\end{theorem}

\noindent File \leanpath{VOA/\allowbreak Examples/\allowbreak Sl2CosetConcreteVOA.lean}.
\begin{leancode}
theorem cosetCentralCharge :
    AVOA.centralCharge + BVOA.centralCharge -
      3 * ((-4 / 3 : ℂ) + 1) / (((-4 / 3 : ℂ) + 1) + 2) = -22 / 5 := by

def cosetVOA : VertexOperatorAlgebra CosetSpace where
  toPreVertexAlgebra := cosetVertexAlgebra
  conformalVector := cosetOmegaState
  conformal_mem := cosetOmegaState_weight
  centralCharge := -22 / 5
\end{leancode}

\begin{dictionary}
\lean{AVOA.centralCharge} is the Sugawara value $3k/(k+2)=-6$ at $k=-4/3$ (\lean{aCentralCharge}); \lean{BVOA.centralCharge}$=1$.
The third term is the Sugawara central charge $3k'/(k'+2)$ at $k'=-1/3$. The \lean{VertexOperatorAlgebra} structure
(Chapter~\ref{ch:L04}) is built by its fields; we show the first five: the underlying vertex algebra on \lean{CosetSpace}, the conformal
vector, its weight, the central charge $-22/5$, and the Virasoro relations (whose proof is the commutation argument below). The remaining
fields ($L_0$-grading, translation, the vacuum grade, positivity of the grading) are proved in the file.

\textbf{What is \emph{not} claimed.} Nothing here identifies \lean{cosetVOA} with $L$, and simplicity of \lean{cosetVOA} is not proved.
The library reduces injectivity of the coset-to-cohomology map to simplicity
(\lean{leeYangCosetToH0_injective_of_cosetSimple}) but never uses it.
\end{dictionary}

\paragraph{Proof on paper.} The diagonal Sugawara modes $L^{\mathrm{diag}}_n$ commute with every operator that commutes with the diagonal
currents; in particular with $L_n-L^{\mathrm{diag}}_n$, where $L_n$ are the modes of $\omega_A\otimes\ltwentyeightvac+\ltwentyeightvac\otimes\omega_B$.
If two Virasoro families $L_n$ (central charge $c_1$) and $D_n$ (central charge $c_2$) satisfy $[D_m,L_n-D_n]=0$, then $L_n-D_n$ is a Virasoro
family of central charge $c_1-c_2$ (\lean{virasoro_difference_of_commute}, a direct computation). The arithmetic is
$\tfrac{3(-4/3)}{2/3}=-6$, $\tfrac{3(-1/3)}{5/3}=-\tfrac35$. \qed

\begin{physics}
GKO: $\mathrm{SU}(2)_{k}\times\mathrm{SU}(2)_1/\mathrm{SU}(2)_{k+1}$ has $c=1-6/((k+2)(k+3))$; for $k=m-2$ a positive integer this is the
unitary series, and the analytic continuation to $k=-4/3$ lands on Lee--Yang.
\end{physics}

\begin{status}
Proved; axioms standard (re-checked: \lean{Sl2CosetConcrete.cosetVOA}, \lean{Sl2CosetConcrete.cosetCentralCharge}). Unconditional.
\end{status}

\begin{theorem}[Coset states are relative cocycles; the map to $H^0$]\label{l28:thm:cosetmap}
For a commutant state $x$, the state $x\otimes\ltwentyeightvac_U\otimes\ltwentyeightvac_{\mathrm{gh}}$ is relative and $\ltwentyeightQ$-closed, and
$x\mapsto[x\otimes\ltwentyeightvac\otimes\ltwentyeightvac]$ is a linear map into the ghost-number-zero cohomology which intertwines the vertex
products.
\end{theorem}

\noindent File \leanpath{VOA/\allowbreak Examples/\allowbreak Sl2CosetConcreteCohomology.lean}.
\begin{leancode}
theorem cosetToComplex_apply (x : CosetSpace) :
    cosetToComplex x = (x.val ⊗ₜ[ℂ] USVA.vacuum) ⊗ₜ[ℂ] sl2GhostSVA.vacuum := rfl

theorem cosetToComplex_closed (x : CosetSpace) : leeYangBRSTComplex.Q (cosetToComplex x) = 0 :=

theorem cosetToComplex_relative (x : CosetSpace) :
    cosetToComplex x ∈ relativeCarrier leeYangMatterAction.toAffineSl2Action :=

theorem cosetToComplex_mode (x y : CosetSpace) (n : ℤ) :
    cosetToComplex (cosetMode diagonalAction.toAffineSl2Action diagonalAction.translation x n y) =
      mode (brstSVA.Y (cosetToComplex x)) n (cosetToComplex y) :=

def leeYangCosetToH0 : CosetSpace →ₗ[ℂ] leeYangBRSTComplex.H0 :=
\end{leancode}

\begin{dictionary}
\lean{x.val} is the underlying vector of $x\in$ \lean{CosetSpace}. \lean{leeYangBRSTComplex.H0} is the ghost-number-zero part of the
cohomology \lean{Coh} $=\ker\ltwentyeightQ/\operatorname{im}\ltwentyeightQ$ (Chapter~\ref{ch:L04}). \lean{cosetMode} is the vertex product of the
coset VOA. The last line only names the map (its body, a specialization of the generic \lean{Sl2Coset.leeYangCosetToH0}, is omitted).
\end{dictionary}

\paragraph{Proof on paper.} Write $j=\sum_aJ^a_{\mathrm{diag}+U}c^a+\cdots$. On $x\otimes\ltwentyeightvac\otimes\ltwentyeightvac$ every term of
$\ltwentyeightQ$ contains either a non-negative mode of a diagonal current acting on $x$ (zero, since $x$ is in the commutant), a non-negative
mode of a $U$-current or of a ghost acting on a vacuum (zero), or a positive ghost mode. So $\ltwentyeightQ$ kills it; the zero modes of the
total current and of $b^a$ kill it for the same reason. Products are preserved because the auxiliary and ghost vacua are units. \qed

\begin{physics}
Gauge-neutral operators of the matter, dressed with nothing, are BRST-closed: the coset sits inside the gauged chiral algebra.
\end{physics}

\begin{status}
Proved; axioms standard. Unconditional.
\end{status}

\begin{theorem}[The Virasoro image surjects onto Lee--Yang]\label{l28:thm:virimage}
Let $\phi:V(-\tfrac{22}{5},0)\to H^0$ be the composite of the VOA map $V(c,0)\to(\text{coset})$ determined by $\omega\mapsto\omega_{\mathrm{coset}}$
with the map of Theorem~\ref{l28:thm:cosetmap}. Then $\ker\phi$ is contained in the maximal ideal of $V(c,0)$.
\end{theorem}

\noindent Files \leanpath{VOA/\allowbreak Examples/\allowbreak LeeYangShortcutVirasoro.lean},
\leanpath{LeeYangShortcutImage.lean}.
\begin{leancode}
theorem conformalEval_hom (W : VertexOperatorAlgebra E) {c : ℂ}
    (hc : W.centralCharge = c) :
    IsVOAHom (virasoroVOA c) W (conformalEval W hc) := by

def virasoroToCoset : VirUniv LeeYang.leeYangC →ₗ[ℂ] Sl2CosetConcrete.CosetSpace :=
  conformalEval Sl2CosetConcrete.cosetVOA rfl

def virasoroToH0 : VirUniv LeeYang.leeYangC →ₗ[ℂ] leeYangBRSTComplex.H0 :=
  leeYangCosetToH0.comp virasoroToCoset

theorem virasoroToH0_ker_le_maxIdeal :
    LinearMap.ker virasoroToH0 ≤ maxIdealVir LeeYang.leeYangC :=
\end{leancode}

\begin{dictionary}
\lean{conformalEval W hc} is the linear map $V(c,0)\to W$ sending $L_{-n_1}\cdots L_{-n_k}\ltwentyeightvac$ to the corresponding product of modes of
the conformal vector of $W$ applied to the vacuum; \lean{conformalEval_hom} says it is a VOA homomorphism (\lean{IsVOAHom}: vacuum and
all products preserved) whenever $W$ has central charge $c$. \lean{VirUniv c} is the carrier of $V(c,0)$ (Chapter~\ref{ch:L06}),
\lean{maxIdealVir c} its maximal proper ideal, so $V(c,0)/$\lean{maxIdealVir c}$=L(c,0)$. \lean{LinearMap.ker} is the kernel.

\textbf{Consequence used later.} Since $\ker\phi\subseteq\mathcal M_c$, the quotient map $V(c,0)\to L$ factors through the image of $\phi$:
there is a surjection $\operatorname{im}\phi\to L$, and it maps the image of the length-$\le r$ words onto $F_rL$. This is all that the
squeeze (Theorem~\ref{l28:thm:classfamily}) needs from the coset.
\end{dictionary}

\paragraph{Proof on paper.} $\ker\phi$ is an ideal of $V(c,0)$ because $\phi$ intertwines all products. It does not contain the vacuum, since the
vacuum class of $H^0$ is nonzero (Theorem~\ref{l28:thm:complex}). Every proper ideal of $V(c,0)$ is contained in the maximal one
(Chapter~\ref{ch:L06}). \qed

\begin{physics}
The stress tensor of the gauged theory generates a Virasoro algebra of central charge $-22/5$; whatever it generates is at least as big as
Lee--Yang, because the only way to be smaller would be to kill the identity operator.
\end{physics}

\begin{status}
Proved; axioms standard (re-checked: \lean{LeeYangShortcut.virasoroToH0_ker_le_maxIdeal}). Unconditional.
\end{status}

\begin{remark}[Why the coset identification is avoided]\label{l28:rem:noacl}
The paper proof identifies the coset VOA with $L$ (Arakawa--Creutzig--Linshaw) and then the coset with $H^0$. Both identifications are hard
to formalize. The ``shortcut'' route (\leanpath{LeeYangShortcut*}) needs only the two one-sided facts of Theorems~\ref{l28:thm:cosetmap} and
\ref{l28:thm:virimage}: $L$ is a quotient of a subspace of $H^0$. Under the hypothesis $\ltwentyeightGU(A)$ the dimension
squeeze then shows, as a \emph{consequence}, that the image of $V(c,0)$ is all of $H^0$ and that $\ker\phi=\mathcal M_c$ exactly; in particular
$H^0\cong L$. This is a consequence on paper of the proof (inside it, \lean{maps_eq_of_finrank_squeeze} gives the equality of each filtration step);
no Lean declaration states ``$\operatorname{im}\phi=H^0$'' or ``$\ker\phi=\mathcal M_c$'', and without $\ltwentyeightGU(A)$ nothing beyond
Theorems~\ref{l28:thm:cosetmap} and \ref{l28:thm:virimage} is claimed. The superseded assembly of
Warning~\ref{l28:warn:clash} assumed \lean{CosetIdentification} instead.
\end{remark}

%======================================================================
\section{Hodge theory of a weakly graded-unitary complex}\label{l28:sec:hodge}
%======================================================================

This section is pure finite-dimensional linear algebra: Beem--Garner's positivity argument, at one conformal weight, for an abstract
complex. Section~\ref{l28:sec:package} instantiates it.

\begin{definition}[Kähler complex]\label{l28:def:kahler}
A \emph{Kähler complex} is a finite-dimensional complex vector space $K$ with:
a positive definite sesquilinear form $P$ (conjugate-linear in the first slot); a decomposition $K=\bigoplus_{(j,g)\in\Z^2}K_{(j,g)}$
into $P$-orthogonal pieces, $j=2R$ and $g$ the ghost number; four linear maps $Q^-,S^-,Q^+,S^+$ shifting the bidegree by $(1,1)$,
$(-1,1)$, $(1,-1)$, $(-1,-1)$ respectively; a sign $\varepsilon\in\{\pm1\}$ with
\begin{equation*}
P(Q^-x,y)=\varepsilon P(x,S^+y),\qquad P(S^-x,y)=-\varepsilon P(x,Q^+y);
\end{equation*}
and the relations $(\ltwentyeightQ^-)^2=0$ and $\ltwentyeightQ^-\ltwentyeightQ^++\ltwentyeightQ^+\ltwentyeightQ^-=0$ for
$\ltwentyeightQ^\pm=Q^\pm+S^\pm$.
\end{definition}

\noindent File \leanpath{VOA/\allowbreak Examples/\allowbreak GradedUnitaryCohomologyKahler.lean}.
\begin{leancode}
structure KahlerComplex (K : Type u) [AddCommGroup K] [Module ℂ K] where
  P : SesqForm K
  pos : IsPosOn P ⊤
  deg : ℤ × ℤ → Submodule ℂ K
  span : iSup deg = ⊤
  orth : ∀ a b, a ≠ b → ∀ x ∈ deg a, ∀ y ∈ deg b, P x y = 0
  Qm : K →ₗ[ℂ] K
  Sm : K →ₗ[ℂ] K
  Qp : K →ₗ[ℂ] K
  Sp : K →ₗ[ℂ] K
  Qm_deg : ∀ a, ∀ x ∈ deg a, Qm x ∈ deg (a + shQm)
  Sm_deg : ∀ a, ∀ x ∈ deg a, Sm x ∈ deg (a + shSm)
  Qp_deg : ∀ a, ∀ x ∈ deg a, Qp x ∈ deg (a + shQp)
  Sp_deg : ∀ a, ∀ x ∈ deg a, Sp x ∈ deg (a + shSp)
  c : ℝ
  c_sq : c * c = 1
  adj_Qm : ∀ x y, P (Qm x) y = (c : ℂ) * P x (Sp y)
  adj_Sm : ∀ x y, P (Sm x) y = -(c : ℂ) * P x (Qp y)
  sq_zero : (Qm + Sm).comp (Qm + Sm) = 0
  anticomm : (Qm + Sm).comp (Qp + Sp) + (Qp + Sp).comp (Qm + Sm) = 0
\end{leancode}

\begin{dictionary}
\lean{SesqForm K} is a sesquilinear form and \lean{IsPosOn P ⊤} means $P(x,x)$ is a positive real for every nonzero $x$
(Definition~\ref{l19:def:isposon}). \lean{iSup deg = ⊤}: the pieces span; with \lean{orth} and positivity they are independent, so the sum is
direct. \lean{shQm}$=(1,1)$, \lean{shSm}$=(-1,1)$, \lean{shQp}$=(1,-1)$, \lean{shSp}$=(-1,-1)$. \lean{c : ℝ} with \lean{c * c = 1} is
$\varepsilon$. \lean{(Qm + Sm).comp (Qm + Sm) = 0} is $(\ltwentyeightQ^-)^2=0$.

\textbf{Against Beem--Garner.} \lean{P} is their modified form $\langle x|y\rangle=((\sigma\circ\rho)x,y)$, and \lean{pos} is their ``weak
graded unitarity'' of the complex (their Theorem~3.20); the four pieces and the two adjoint relations are their Lemma~3.16, in which
\emph{goodness} of the action is exactly the absence of a component of $\ltwentyeightQ^-$ preserving $R$. The other two adjoint relations
$(Q^+)^\dagger=-\varepsilon S^-$ and $(S^+)^\dagger=\varepsilon Q^-$ are derived (\lean{adj_Qp}, \lean{adj_Sp}); Hermitian symmetry of $P$
is derived from positivity (\lean{herm}); $(\ltwentyeightQ^+)^2=0$ is not needed.

\textbf{Non-vacuity.} The structure is inhabited by a five-dimensional example with nonzero differentials and a nonzero harmonic vector
(Example~\ref{l28:ex:quartet}), and by the actual complex (Section~\ref{l28:sec:package}).
\end{dictionary}

\begin{physics}
$P$ is the positive inner product that the four-dimensional theory induces on chiral-algebra states after twisting by $i^{2R}$; the two
differentials $\ltwentyeightQ^\pm$ are the BRST charge and its $SU(2)_R$ partner, and the identities are the chiral shadow of an $\mathrm{Sp}(2)$
(Kähler) structure, as in Hodge theory on a compact Kähler manifold, where $\partial$, $\bar\partial$ and their adjoints satisfy analogous
relations.
\end{physics}

\begin{status}
Definition (a \lean{structure}). Instances: Example~\ref{l28:ex:quartet} and the concrete complex.
\end{status}

\begin{theorem}[Harmonic vectors; Beem--Garner Remark 4.4 in strong form]\label{l28:thm:harm}
Put $D=\ltwentyeightQ^-$ and $D^\dagger=\varepsilon(S^+-Q^+)$, its $P$-adjoint. A vector is harmonic ($Dx=0=D^\dagger x$) if and only if
$Q^-x=S^-x=Q^+x=S^+x=0$. In particular the harmonic space is bigraded, and every class of $\ker D/\operatorname{im}D$ has a unique harmonic
representative.
\end{theorem}

\noindent Files \leanpath{VOA/\allowbreak Examples/\allowbreak GradedUnitaryCohomologyKahler.lean},
\leanpath{LeeYangCosetTransferHodge.lean}.
\begin{leancode}
def D : K →ₗ[ℂ] K := C.Qm + C.Sm

def Dadj : K →ₗ[ℂ] K := (C.c : ℂ) • (C.Sp - C.Qp)

def harm : Submodule ℂ K := harmonicSpace C.D C.Dadj

theorem mem_harm_iff (x : K) :
    x ∈ C.harm ↔ C.Qm x = 0 ∧ C.Sm x = 0 ∧ C.Qp x = 0 ∧ C.Sp x = 0 := by

def harmonicSpace (Q Qadj : K →ₗ[ℂ] K) : Submodule ℂ K :=
  LinearMap.ker Q ⊓ LinearMap.ker Qadj
\end{leancode}

\begin{dictionary}
\lean{D} is $\ltwentyeightQ^-$, \lean{Dadj} is $\varepsilon(S^+-Q^+)$, and \lean{adj_D} proves $P(Dx,y)=P(x,D^\dagger y)$. \lean{harmonicSpace Q Qadj}
is $\ker Q\cap\ker Q^\dagger$. The statement is exactly the first sentence. The unique-representative statement is
\lean{harmonicEquivCohomology} (a linear equivalence between the harmonic space and $\ker D/\operatorname{im}D$, same file, with
\lean{harmonicToCohomology_injective} and \lean{_surjective}).
\end{dictionary}

\paragraph{Proof on paper.} If $Dx=0$ and $D^\dagger x=0$, then $Q^-x=-S^-x$ and $S^+x=Q^+x$. Consider $E=Q^+x+S^+x=\ltwentyeightQ^+x$ and
$F=Q^-x-S^-x$. Using the adjoint relations and $\ltwentyeightQ^-\ltwentyeightQ^++\ltwentyeightQ^+\ltwentyeightQ^-=0$ split into bidegrees (the
``fine'' relations \lean{fine_QmQp}, \lean{fine_SmSp}, \lean{Qm_sq}, obtained by projecting onto bidegrees), one finds
$P(E,E)+P(F,F)=0$. By positivity $E=F=0$, and combined with $Q^-x+S^-x=0$, $S^+x-Q^+x=0$ all four pieces vanish. Conversely if all four
vanish, so do $D$ and $D^\dagger$. Homogeneous components of a harmonic vector are then harmonic, because each piece is homogeneous. The
unique-representative statement is the usual finite-dimensional Hodge decomposition $K=\operatorname{im}D\oplus\mathcal H\oplus\operatorname{im}D^\dagger$
for a positive form. \qed

\begin{physics}
Harmonic representatives are the ``on-shell'' gauge-invariant states with definite $R$-charge; the strong form says that they are
simultaneously annihilated by both supercharges and their adjoints.
\end{physics}

\begin{status}
Proved; axioms standard (re-checked: \lean{mem_harm_iff}). Unconditional (for any \lean{KahlerComplex}).
\end{status}

\begin{theorem}[The harmonic bound]\label{l28:thm:harmbound}
Let $H(Q^-)_{(j,g)}$ be the cohomology of the leading piece $Q^-$ alone in bidegree $(j,g)$. Harmonic vectors of bidegree $(j,g)$ inject into
$H(Q^-)_{(j,g)}$. Consequently, for a finite set $S$ of values of $2R$, the classes of $\ker D/\operatorname{im}D$ with a representative of
ghost number $g$ and $2R\in S$ span a space of dimension at most $\sum_{j\in S}\dim H(Q^-)_{(j,g)}$.
\end{theorem}

\noindent Files \leanpath{VOA/\allowbreak Examples/\allowbreak GradedUnitaryCohomologyKahler.lean},
\leanpath{GradedUnitaryCohomologyFiltered.lean}.
\begin{leancode}
def leadCycles (a : ℤ × ℤ) : Submodule ℂ K := LinearMap.ker C.Qm ⊓ C.deg a

def leadBoundaries (a : ℤ × ℤ) : Submodule ℂ (C.leadCycles a) :=
  (Submodule.map C.Qm (C.deg (a - shQm))).comap (C.leadCycles a).subtype

abbrev leadCohom (a : ℤ × ℤ) := C.leadCycles a ⧸ C.leadBoundaries a

def harmDeg (a : ℤ × ℤ) : Submodule ℂ K := C.harm ⊓ C.deg a

theorem harmToLead_injective (a : ℤ × ℤ) : Function.Injective (C.harmToLead a) := by

def cohomFilt (A : Set (ℤ × ℤ)) : Submodule ℂ (cycleQuotient C.D) :=
  Submodule.map (boundariesInCycles C.D).mkQ
    (Submodule.comap (LinearMap.ker C.D).subtype (C.degOn A))

def rowSet (S : Finset ℤ) (g : ℤ) : Set (ℤ × ℤ) := {a | a.1 ∈ S ∧ a.2 = g}

theorem finrank_cohomFilt_row_le_lead (S : Finset ℤ) (g : ℤ) :
    finrank ℂ (C.cohomFilt (rowSet S g)) ≤ ∑ j ∈ S, finrank ℂ (C.leadCohom (j, g)) := by
\end{leancode}

\begin{dictionary}
\lean{leadCycles a} $=\ker Q^-\cap K_a$; \lean{leadBoundaries a} $=Q^-(K_{a-(1,1)})$ viewed inside the cycles; \lean{leadCohom a} is the
quotient, $H(Q^-)_a$. \lean{harmToLead a} sends a harmonic vector of bidegree $a$ to its $Q^-$-class. \lean{cohomFilt A} is the image in
$\ker D/\operatorname{im}D$ (\lean{cycleQuotient C.D}) of the closed vectors supported on the bidegrees in $A$; \lean{rowSet S g} is
$\{(j,g):j\in S\}$. Exactly the theorem. At ghost number $g=0$ and $S=\{0,1,\dots,2r\}$ it reads
$\dim F_rH^0\le\sum_{j=0}^{2r}\dim H(Q^-)_{(j,0)}$, inequality (c) of Remark~\ref{l28:rem:arc}.
\end{dictionary}

\paragraph{Proof on paper.} If a harmonic $x$ of bidegree $a$ were $Q^-y$, then
$P(x,x)=P(Q^-y,x)=\varepsilon P(y,S^+x)=0$, since $S^+x=0$ by Theorem~\ref{l28:thm:harm}; so $x=0$. For the second statement: a class with a
representative supported on $A$ has its harmonic representative supported on $A$ (the harmonic projection respects the bigrading,
\lean{exists_harm_rep_mem_degOn}), so the space of such classes is isomorphic to the harmonic vectors supported on $A$, whose dimension is
$\sum_{a\in A}\dim\mathcal H_a\le\sum_{a\in A}\dim H(Q^-)_a$. \qed

\begin{physics}
Counting: gauge-invariant states of $R$-charge $\le r$ are at most as many as the states of the ``free'' (leading-order) BRST problem. The
leading problem is computable, which is why this bound is useful.
\end{physics}

\begin{status}
Proved; axioms standard (re-checked: \lean{harmToLead_injective}, \lean{finrank_cohomFilt_row_le_lead}). Unconditional.
\end{status}

\begin{definition}[Form and conjugation data on a Kähler complex]\label{l28:def:gudata}
A \emph{GU datum} on a Kähler complex $C$ is a bilinear form $B$, an anti-linear map $\rho$ and the sesquilinear form $h(x,y)=B(\rho x,y)$, such
that $P(x,y)=i^{j}h(x,y)$ for $x$ of bidegree $(j,g)$, and $\rho$ maps the harmonic vectors of bidegree $(j,0)$ onto themselves. On the harmonic
model of $H^0$ define the flag $G_j=\bigoplus_{L\le i\le j}\mathcal H_{(i,0)}$ and its layers $G_j\cap B\text{-}\mathrm{orth}(G_{j-1})$.
\end{definition}

\noindent File \leanpath{VOA/\allowbreak Examples/\allowbreak GradedUnitaryCohomologyTransfer.lean}.
\begin{leancode}
structure GUData (C : KahlerComplex K) where
  B : LinearMap.BilinForm ℂ K
  ρ : K →ₗ⋆[ℂ] K
  h : SesqForm K
  h_eq : ∀ x y, h x y = B (ρ x) y
  P_eq : ∀ j g : ℤ, ∀ x ∈ C.deg (j, g), ∀ y, C.P x y = Complex.I ^ j * h x y
  ρ_maps : ∀ j : ℤ, ∀ x ∈ C.harmDeg (j, 0), ρ x ∈ C.harmDeg (j, 0)
  ρ_onto : ∀ j : ℤ, ∀ y ∈ C.harmDeg (j, 0), ∃ x ∈ C.harmDeg (j, 0), ρ x = y

def flag (L j : ℤ) : Submodule ℂ K :=
  C.harm ⊓ C.degOn (KahlerComplex.rowSet (Finset.Icc L j) 0)
\end{leancode}

\begin{dictionary}
\lean{K →ₗ⋆[ℂ] K} is an anti-linear map. \lean{Complex.I ^ j} is $i^j$, $j=2R$; for integral $R=p$ it is $(-1)^p$. \lean{flag L j} is the span
of harmonic vectors of ghost number $0$ with $2R\in[L,j]$ (\lean{flag_eq_iSup}). The \lean{layer} (\lean{GUData.layer}, same file) is
\lean{C.flag L j ⊓ G.B.orthogonal (C.flag L (j - 1))}, matching \lean{GradedFlag.layer} of Definition~\ref{l19:def:gu}.
\end{dictionary}

\begin{physics}
$i^{2R}$ is the phase relating the four-dimensional inner product to the chiral-algebra (BPZ) pairing; for integer $R$ it is ABLR's sign
$(-1)^R$.
\end{physics}

\begin{status}
Definition. Instantiated on the actual complex by the Beem--Garner package (Section~\ref{l28:sec:package}).
\end{status}

\begin{theorem}[Cohomology inherits graded unitarity; Beem--Garner Theorem 4.1 at one weight]\label{l28:thm:bg}
For a GU datum on a Kähler complex: (1) $B$ is nondegenerate on every step $G_j$ of the harmonic flag; (2) the layer at $j$ is exactly
$\mathcal H_{(j,0)}$; (3) on the layer at $j=2p$, the form $(-1)^ph$ is positive definite.
\end{theorem}

\noindent File \leanpath{VOA/\allowbreak Examples/\allowbreak GradedUnitaryCohomologyTransfer.lean}.
\begin{leancode}
theorem flag_nondegenerate [FiniteDimensional ℂ K] (L j : ℤ) :
    (G.B.restrict (C.flag L j)).Nondegenerate := by

theorem layer_eq {L j : ℤ} (hLj : L ≤ j) : G.layer L j = C.harmDeg (j, 0) := by

theorem layer_signed_pos {L : ℤ} (p : ℕ) (hL : L ≤ 2 * (p : ℤ)) :
    IsPosOn (signedForm G.h p) (G.layer L (2 * (p : ℤ))) := by
\end{leancode}

\begin{dictionary}
\lean{(G.B.restrict S).Nondegenerate} is clause (iv) on the step; \lean{IsPosOn (signedForm G.h p) S} is clause (vi), with
\lean{signedForm h p}$=(-1)^ph$ (Chapter~\ref{ch:L19}). The section variables are a \lean{KahlerComplex C} and \lean{G : GUData C}; the hypothesis
\lean{hL : L ≤ 2 * p} says the layer index is in the range of the flag. Exactly clauses (iv) and (vi) for the harmonic model.
\end{dictionary}

\paragraph{Proof on paper.} The pieces $\mathcal H_{(m,0)}$ are $B$-orthogonal for different $m$: for $x\in\mathcal H_{(m,0)}$ and $y$ of a different
bidegree, $B(\rho x,y)=h(x,y)=i^{-m}P(x,y)=0$, and $\rho$ maps each piece onto itself. So $G_j=\mathcal H_{(j,0)}\oplus G_{j-1}$ is a $B$-orthogonal sum, which gives (2). On
$\mathcal H_{(j,0)}$, $i^jh=P$ is positive definite, so $h$ is nondegenerate there, hence so is $B$ (as $\rho$ is bijective); an orthogonal sum
of nondegenerate pieces is nondegenerate, which gives (1). For $j=2p$, $i^{2p}=(-1)^p$ and (3) is positivity of $P$. \qed

\begin{physics}
This is the statement that BRST reduction (gauging) preserves the ABLR positivity constraints, because they come from a positive
four-dimensional inner product which survives the restriction to gauge-invariant states.
\end{physics}

\begin{status}
Proved; axioms standard (re-checked: \lean{flag_nondegenerate}, \lean{layer_signed_pos}). Unconditional (for any GU datum).
\end{status}

\begin{example}[The quartet, a non-vacuity witness]\label{l28:ex:quartet}
In \leanpath{GradedUnitaryCohomologyExample.lean} the library builds a Kähler complex on $\C^5$ with basis $e_0=\tilde x$ in bidegree $(0,0)$,
$e_1=Q^-\tilde x$ in $(1,1)$, $e_2=Q^+\tilde x$ in $(1,-1)$, $e_3=Q^-Q^+\tilde x$ in $(2,0)$, and a singlet $e_4$ in $(0,0)$, with the standard
Hermitian form and $\varepsilon=1$. The four vectors $e_0,\dots,e_3$ form Beem--Garner's ``quartet'' (their Figure~3.1), which is
$\ltwentyeightQ^-$-acyclic; the singlet $e_4$ is harmonic, $e_0$ is not, $D\neq0$, and a GU datum with $B$ weighted by $i^{-j}$ makes the
singlet's layer positive (\lean{singlet_signed_pos}). So the hypotheses of Theorem~\ref{l28:thm:bg} are consistent and not degenerate.
\end{example}

\begin{warning}[title={Dead ends in the abstract Hodge theory}]\label{l28:warn:signtwist}
\begin{itemize}
\item \textbf{One differential is not enough.} A single differential $D$ with an $R$-filtration does not make harmonic representatives
  homogeneous: on $K=\langle a,b,c\rangle$ with $R(a)=0$, $R(b)=\tfrac12$, $R(c)=-\tfrac12$ and $Da=b+c$, the harmonic vector $b-c$ is not
  homogeneous. The second differential $\ltwentyeightQ^+$ and the Kähler identity are essential (project notes; hand counterexample).
\item \textbf{A GU datum cannot force the vacuum normalization.} $(B,\rho)\mapsto(-B,-\rho)$ preserves every axiom of \lean{GUData} (since
  $h$ is unchanged), so the normalization $B(\text{vacuum},\text{vacuum})=1$ is extra information
  (\lean{LeeYangShortcut.signTwistGUData_not_both_normalized}, \leanpath{LeeYangShortcutMetric.lean}, proved). It is supplied on the actual
  complex by Theorem~\ref{l28:thm:repmetric}.
\end{itemize}
\end{warning}

%======================================================================
\section{Transfer by a dimension squeeze}\label{l28:sec:squeeze}
%======================================================================

\begin{theorem}[The single-piece squeeze]\label{l28:thm:squeeze}
Let $i:C\to H$ be injective and linear, $P_r\le C$, $F_r\le H$ finite-dimensional, with $i(P_r)\subseteq F_r$ and $\dim F_r\le d_r=\dim P_r$ for
every $r$. Then $i(P_r)=F_r$ for every $r$.
\end{theorem}

\noindent File \leanpath{VOA/\allowbreak Examples/\allowbreak LeeYangCosetTransferDimension.lean}.
\begin{leancode}
theorem maps_eq_of_finrank_squeeze
    (i : C →ₗ[𝕜] H) (hi : Function.Injective i)
    (P : ι → Submodule 𝕜 C) (F : ι → Submodule 𝕜 H) (d : ι → ℕ)
    (hfiltered : ∀ r, (P r).map i ≤ F r)
    (hupper : ∀ r, finrank 𝕜 (F r) ≤ d r)
    (hexact : ∀ r, d r = finrank 𝕜 (P r)) : ∀ r, (P r).map i = F r := by
\end{leancode}

\begin{dictionary}
\lean{(P r).map i} is the image $i(P_r)$. \lean{d : ι → ℕ} is the numerical bound; in the application it is $\sum_{s\le r}G_s(N)$. The section
variable \lean{[FiniteDimensional 𝕜 H]} makes all pieces finite-dimensional. Exactly the theorem.
\end{dictionary}

\paragraph{Proof on paper.} $\dim i(P_r)=\dim P_r=d_r\ge\dim F_r$ and $i(P_r)\subseteq F_r$. \qed

\begin{physics}
A purely numerical step: if the states one can construct are at least as many as an upper bound on the states that can exist, one has all
of them.
\end{physics}

\begin{status}
Proved; axioms standard (re-checked). Unconditional.
\end{status}

\begin{theorem}[Transfer of graded unitarity through a filtered isometry]\label{l28:thm:filterediso}
If a weight-preserving isometric embedding $E$ of $A$ into $T$ that intertwines the real structures maps every flag step of $A$ \emph{onto} the
corresponding step of $T$, and $T$ is graded unitary, then $A$ is graded unitary.
\end{theorem}

\noindent File \leanpath{VOA/\allowbreak Examples/\allowbreak LeeYangCosetTransferFiltered.lean}.
\begin{leancode}
theorem gradedUnitary_of_filteredIso
    (E : GradedIsometricEmbedding A T BA BT RA RT)
    (FA : GradedFlag A) (FT : GradedFlag T)
    (hpiece : ∀ n p, (FA.pieceIn p n).map (E.weightMap n) = FT.pieceIn p n)
    (hambient : GradedUnitary BT RT FT) : GradedUnitary BA RA FA := by
\end{leancode}

\begin{dictionary}
\lean{GradedIsometricEmbedding A T BA BT RA RT} is Definition~\ref{l19:def:gie}: weightwise injective maps preserving the forms and commuting
with the conjugations. \lean{E.weightMap n} is its weight-$n$ part and \lean{pieceIn p n} the step $F_pV_n$. \lean{hpiece} is ``onto every
step''. Exactly the theorem. It strengthens Theorem~\ref{l19:thm:transfer} (which assumes orthogonal reduction): ontoness gives orthogonal
reduction for free (\lean{orthogonallyReducingAt_of_surjective}).
\end{dictionary}

\paragraph{Proof on paper.} If $E$ maps $F_pA_n$ onto $F_pT_n$ for every $p$ and the flags exhaust, $E$ is onto at weight $n$, hence an
isometric isomorphism of filtered spaces; layers go to layers and both clauses transfer. \qed

\begin{physics}
Two chiral algebras with the same states, the same inner products and the same $R$-filtration have the same positivity properties.
\end{physics}

\begin{status}
Proved; axioms standard (re-checked). Unconditional.
\end{status}

\begin{theorem}[Graded unitarity at one weight from a family of cohomology classes]\label{l28:thm:classfamily}
Let $C$ be a Kähler complex with a GU datum, $W$ a VOA with a flag $F$ at weight $n$ and forms $\beta$, $h$. Let $f:U\to\ker D/\operatorname{im}D$
and $g:U\to W_n$ be linear, $g$ surjective, $\ker f\subseteq\ker g$, and $P_p\le U$ an exhaustive filtration with $g(P_p)=F_pW_n$ and
$f(P_p)\subseteq F_{p}(\ker D/\operatorname{im}D)$ (classes with representatives of $R\le p$ at ghost number zero). Suppose
$\sum_{j=0}^{2p}\dim H(Q^-)_{(j,0)}=\dim F_pW_n$ for every $p$, that $\beta$ and $h$ pulled back along $g$ agree with $B$ and $h$ on the harmonic
lifts of $f$, and that harmonic vectors of odd $2R$ at ghost number zero vanish. Then $W$ is graded unitary at weight $n$.
\end{theorem}

\noindent File \leanpath{VOA/\allowbreak Examples/\allowbreak LeeYangShortcutHodge.lean}.
\begin{leancode}
theorem gradedUnitaryAt_of_class_family {U : Type*} [AddCommGroup U] [Module ℂ U]
    (C : KahlerComplex K) (G : GUData C) (F : GradedFlag W) (n : ℤ)
    (β : LinearMap.BilinForm ℂ (W.grading n)) (h : SesqForm (W.grading n))
    (f : U →ₗ[ℂ] cycleQuotient C.D) (g : U →ₗ[ℂ] W.grading n)
    (hgs : Function.Surjective g) (hker : LinearMap.ker f ≤ LinearMap.ker g)
    (P : ℕ → Submodule ℂ U)
    (hquotient : ∀ p, (P p).map g = F.pieceIn p n)
    (hexhaustive : ∀ x : U, ∃ p, x ∈ P p)
    (hfiltered : ∀ p, (P p).map f ≤
      C.cohomFilt (KahlerComplex.rowSet (Finset.Icc 0 (2 * (p : ℤ))) 0))
    (hcharacter : ∀ p : ℕ,
      (∑ j ∈ Finset.Icc 0 (2 * (p : ℤ)), Module.finrank ℂ (C.leadCohom (j, 0))) =
        Module.finrank ℂ (F.pieceIn p n))
    (hB : ∀ x y, β (g x) (g y) = G.B (harmonicClassLift C f x) (harmonicClassLift C f y))
    (hh : ∀ x y, h (g x) (g y) = G.h (harmonicClassLift C f x) (harmonicClassLift C f y))
    (hintegral : IntegralHarmonicSupport C) : GradedUnitaryAt F n β h := by
\end{leancode}
\noindent File \leanpath{VOA/\allowbreak Examples/\allowbreak LeeYangTransferAssemblyHodge.lean}.
\begin{leancode}
def IntegralHarmonicSupport (C : KahlerComplex K) : Prop :=
  ∀ p : ℕ, C.harmDeg (2 * (p : ℤ) + 1, 0) = ⊥
\end{leancode}

\begin{dictionary}
\textbf{Hypotheses one by one.} \lean{f} (in the application: universal Virasoro words $\mapsto$ their cohomology classes, $U=V(c,0)_n$) and
\lean{g} (the quotient map onto $L_n$). \lean{hgs}: $g$ onto. \lean{hker}: $\ker f\subseteq\ker g$, which is Theorem~\ref{l28:thm:virimage}.
\lean{P}: the length filtration of $V(c,0)_n$; \lean{hquotient}: its image is the length flag of $L_n$; \lean{hexhaustive}: it exhausts.
\lean{hfiltered}: words of length $\le p$ give classes with representatives of $2R\in[0,2p]$ at ghost number zero; in the application
this is \lean{endpointKahler_filtered hA} (Theorem~\ref{l28:thm:filtered}). \lean{hcharacter}: the dimension equality, which in the application is (d)--(f) of
Remark~\ref{l28:rem:arc}. \lean{hB}, \lean{hh}: the forms agree on harmonic lifts (\lean{harmonicClassLift C f x} is the harmonic representative
of $f(x)$; Theorem~\ref{l28:thm:repmetric}). \lean{hintegral}: \lean{IntegralHarmonicSupport C}, no harmonic vectors at odd $2R$ and ghost number
$0$, which follows from the invariant-symbol bound (\lean{integralHarmonicSupport_of_symbols}) because the leading cohomology vanishes at odd $j$.

\textbf{Injectivity of $f$ is not assumed}, and neither is injectivity of $g$ on the image of $f$: the theorem proves it.
\end{dictionary}

\paragraph{Proof on paper.} Factor $g$ through the range of $f$ (possible since $\ker f\subseteq\ker g$). The range of $f$, filtered by the images of
$P_p$, maps into the harmonic flag (via harmonic representatives) and onto $F_pW_n$. Then
$\dim F_pW_n\le\dim f(P_p)\le\dim F_p(\ker D/\operatorname{im}D)\le\sum_{j\le2p}\dim H(Q^-)_{(j,0)}=\dim F_pW_n$
(the third inequality is Theorem~\ref{l28:thm:harmbound}), so all are equal; by Theorem~\ref{l28:thm:squeeze} the harmonic lifts of $f(P_p)$ are the
whole harmonic flag step, and $g$ is injective on the range of $f$. The resulting map from the harmonic flag of $H^0$ to $F_\bullet W_n$ is a
filtered isometric isomorphism (by \lean{hB}, \lean{hh}); odd layers vanish by \lean{hintegral}, so the harmonic layer at $j=2p$ corresponds to
$\mathcal L_pW_n$. Theorem~\ref{l28:thm:bg} gives clauses (iv) and (vi) on the harmonic side, and they transfer. \qed

\begin{physics}
The Lee--Yang states are realized, one for one and with the same norms and $R$-charges, as gauge-invariant states of the BRST problem; their
positivity is then inherited from the latter.
\end{physics}

\begin{status}
Proved; axioms standard (re-checked: \lean{LeeYangShortcut.gradedUnitaryAt_of_class_family}). Unconditional (all hypotheses explicit).
\end{status}

%======================================================================
\section{The Beem--Garner package on the actual complex}\label{l28:sec:package}
%======================================================================

To apply Section~\ref{l28:sec:hodge} at weight $N$ one needs, on the actual relative complex: a positive definite modified form, an
$R$-bigrading for which $\ltwentyeightQ$ is good, a second differential with the Kähler identities, and a GU datum whose invariant form and
conjugation agree with those of Lee--Yang on the coset representatives. Positivity needs each matter factor to be graded unitary (or
unitary) and the ghost form to be positive. This section presents the inputs and the theorems that discharge them.

\begin{definition}[Unitary VOA]\label{l28:def:unitary}
A VOA with invariant form $B$ and real structure $\rho$ is \emph{unitary} if $h(x,x)=B(\rho x,x)$ is a positive real for every nonzero
homogeneous $x$.
\end{definition}

\noindent File \leanpath{VOA/\allowbreak Examples/\allowbreak ConcreteHodgeAssembly.lean}.
\begin{leancode}
def Unitary (B : InvariantForm W) (R : RealStructure W) : Prop :=
  ∀ n : ℤ, IsPosOn (gradedSesq B R n) ⊤
\end{leancode}

\begin{dictionary}
\lean{gradedSesq B R n} is $h$ restricted to $V_n$ (Definition~\ref{l19:def:gradedsesq}) and \lean{⊤} the whole weight space. This is
Dong--Lin unitarity (Definition~\ref{l19:def:donglin}) up to the requirement that $\rho$ be an involution, which is part of the
\lean{RealStructure} via its parity. It is a separate predicate, not \lean{GradedUnitary B R F} for any graded flag $F$: a
\lean{GradedFlag} must have $F_0V_N=0$ for $N>0$ and $F_0V_0=\C\ltwentyeightvac$, so ``everything in step $0$'' is not a graded flag. What
corresponds to it is the Beem--Garner $R$-grading that puts all of $B$ in $R$-degree $0$ (\lean{trivialDeg}, Definition~\ref{l28:def:good}).
\end{dictionary}

\begin{physics}
The ordinary Hilbert-space positivity of a unitary CFT.
\end{physics}

\begin{status}
Definition.
\end{status}

\begin{theorem}[$L_1(\ltwentyeightsl)$ is unitary]\label{l28:thm:bunitary}
With its normalized canonical form and the descended compact conjugation, $B=L_1(\ltwentyeightsl)$ is unitary. Moreover its conformal vector
is $\omega=\tfrac14 h_{-1}h_{-1}\ltwentyeightvac$.
\end{theorem}

\noindent Files \leanpath{VOA/\allowbreak Examples/\allowbreak BGPackageBFactor.lean}, \leanpath{LevelOneUnitaryPositivity.lean},
\leanpath{LevelOneHeisSingular.lean}.
\begin{leancode}
def bRealStructure : RealStructure BVOA where

abbrev bForm : InvariantForm BVOA := simpleForm 1 bNoncritical

theorem bUnitary : ConcreteHodge.Unitary bForm bRealStructure := by

abbrev hB : BSpace := bAction.current 2

theorem bConformal_eq : BVOA.conformalVector = (1 / 4 : ℂ) • BVOA.smul hB (-1) hB := by
\end{leancode}

\begin{dictionary}
\lean{bRealStructure} is the compact conjugation descended to $L_1$ (built exactly as \lean{affineRealStructure}, with parity $1$);
\lean{bForm} is \lean{simpleForm 1}; \lean{hB} is the $h$-current state $h_{-1}\ltwentyeightvac$ of $L_1$; \lean{BVOA.smul hB (-1) hB} is
$h_{-1}h_{-1}\ltwentyeightvac$. No hypotheses.

\textbf{Why the second statement is interesting.} At level one the Sugawara vector $\tfrac1{6}(e_{-1}f_{-1}+f_{-1}e_{-1}+\tfrac12h_{-1}h_{-1})\ltwentyeightvac$
equals the free-boson stress tensor $\tfrac14h_{-1}^2\ltwentyeightvac$ in the simple quotient: the coset $L_1/\text{Heisenberg}$ has zero stress
tensor. This is true only at $k=1$, and it is what makes the level-one charge sectors ``Heisenberg-separated''
(\lean{levelOneHeisenbergSeparation}), an input of the invariant-symbol bound.
\end{dictionary}

\paragraph{Proof on paper.} Work in $V^1(\ltwentyeightsl)$ with $h(x,y)=B(\rho x,y)$ and the maximal ideal $I$. For a vector $y$ of energy $E$ and
$h_0$-charge $\mu\ge0$ the Sugawara relation $L_0=E$ and $[e_0,f_0]=h_0$ give
$(6E-\mu-\tfrac{\mu^2}{2})\,h(y,y)=2h(e_0y,e_0y)+\sum_c\sum_{1\le m\le E}2w_c\,h(J^c_my,J^c_my)$ with positive weights $w_c$. Every vector on the
right has lower energy, or the same energy and larger charge. Integrability of $L_1$ gives the charge bound $\mu^2\le4E$ off $I$, and with
$\mu\le2E$ the coefficient $6E-\mu-\mu^2/2$ is positive. Induction (on energy, and downward on charge) shows that $h(y,y)\ge0$, with equality
only on $I$. Negative charges follow by the conjugation, and distinct charges are orthogonal. Pass to the quotient. The identity
$\omega=\tfrac14h_{-1}^2\ltwentyeightvac$ is checked by showing that the difference is a singular vector of weight $2$, hence zero in the simple
quotient. \qed

\begin{physics}
$SU(2)$ level one is the free boson at the self-dual radius; its unitarity is the positivity of a free-field Fock space, and
$\omega=\tfrac14h_{-1}^2\ltwentyeightvac$ is the Frenkel--Kac statement that the WZW stress tensor is that of one boson.
\end{physics}

\begin{status}
Both proved outright; axioms standard (re-checked: \lean{LevelOneUnitary.bUnitary}, \lean{LevelOneHeis.bConformal_eq}).
\end{status}

\begin{theorem}[The auxiliary factor is graded unitary]\label{l28:thm:aux}
For every real $k<-2$, $V^k(\ltwentyeightsl)$ with its normalized form, the compact conjugation and the ordinary current-PBW flag (words of at
most $p$ currents) is graded unitary. In particular $U=V^{-11/3}(\ltwentyeightsl)$ is.
\end{theorem}

\noindent Files \leanpath{VOA/\allowbreak Examples/\allowbreak AffineSl2BelowCriticalAnchor.lean} (a file of Chapter~\ref{ch:L19}),
\leanpath{LeeYangTransferAssemblyAffine.lean}.
\begin{leancode}
theorem universal_gradedUnitary_belowCritical_unconditional (k : ℝ) (hk : k < -2) :
    GradedUnitary (universalForm (k : ℂ) (noncritical_of_lt hk)) (compactRealFamily k hk)
      (ordinaryFlag (k : ℂ) (noncritical_of_lt hk)) :=

abbrev auxiliaryForm : InvariantForm UVOA := universalForm (-11 / 3) uNoncritical

abbrev auxiliaryRealStructure : RealStructure UVOA :=
  compactRealStructure auxiliaryLevel_real uNoncritical

abbrev auxiliaryFlag : GradedFlag UVOA := ordinaryFlag (-11 / 3) uNoncritical

theorem auxiliaryGradedUnitary :
    GradedUnitary auxiliaryForm auxiliaryRealStructure auxiliaryFlag := by
\end{leancode}

\begin{dictionary}
\lean{noncritical_of_lt hk} proves $k+2\ne0$ from $k<-2$; \lean{compactRealFamily k hk} is the compact real structure at that real level;
\lean{ordinaryFlag} is the flag by number of current letters (Sugawara counted as two). \lean{auxiliaryGradedUnitary} specializes to $k=-11/3$
(the proof, omitted, is a cast from $\mathbb R$ to $\C$). No hypotheses. Note the auxiliary flag is the \emph{ordinary} flag, unlike the corrected flag of
$A$; the BRST $R$-grading uses each factor's own flag.
\end{dictionary}

\paragraph{Proof on paper.} (Developed in the files of Chapter~\ref{ch:L19}.) For $k$ very negative the Gram matrices, rescaled, tend to those of
a free (central) limit in which positivity on the layers is manifest. The ``no-wall'' argument shows that along $k\in(-\infty,-2)$ no prefix Gram
determinant vanishes (a prefix-radical and Sugawara induction), so signs cannot change; continuity then fixes the signs at every $k<-2$.
\qed

\begin{physics}
Heuristically, below the critical level the affine current algebra behaves like a free (Gaussian) sector in the large-$|k|$ limit, and every
layer has the sign predicted by ABLR; this is an analogy, not a statement about a specific four-dimensional theory. That this holds for the universal algebra, with no quotient, is special to $k<-2$.
\end{physics}

\begin{status}
Proved; axioms standard (re-checked: both). Unconditional.
\end{status}

\begin{definition}[Goodness of the actual differential]\label{l28:def:good}
Give each factor its $R$-grading: on $A$ the orthogonal layers of the corrected flag ($R=p$ on layer $p$), on $B$ everything in $R=0$, on $U$ the
layers of the ordinary flag, and on the ghosts a positive grading by weights. The differential is \emph{good} on relative states if it maps
bidegree $(2R,\mathrm{gh})$ into $(2R\pm1,\mathrm{gh}+1)$, with no $R$-preserving component.
\end{definition}

\noindent File \leanpath{VOA/\allowbreak Examples/\allowbreak ConcreteGoodInterpreted.lean}.
\begin{leancode}
def aDeg : (ℤ × ℤ) × ℂ → Submodule ℂ ASpace := layerDeg affineForm affineFlag

def bDeg : (ℤ × ℤ) × ℂ → Submodule ℂ BSpace := trivialDeg BVOA

def uDeg : (ℤ × ℤ) × ℂ → Submodule ℂ USpace := layerDeg auxiliaryForm auxiliaryFlag

def matterDeg : (ℤ × ℤ) × ℂ → Submodule ℂ MatterSpace := tensorDeg (tensorDeg aDeg bDeg) uDeg

def brstDeg {PG : SesqForm Sl2GhostSpace} (Gh : GhostRGrading PG) :
    (ℤ × ℤ) × ℂ → Submodule ℂ BRSTSpace :=
  tensorDeg matterDeg Gh.deg

def RelativeDifferentialGood {PG : SesqForm Sl2GhostSpace} (Gh : GhostRGrading PG) : Prop :=
  ∀ (n j g : ℤ) (x : BRSTSpace), x ∈ relativeWeight n → x ∈ brstDeg Gh ((j, g), (n : ℂ)) →
    leeYangBRSTComplex.Q x ∈
      brstDeg Gh ((j + 1, g + 1), (n : ℂ)) ⊔ brstDeg Gh ((j - 1, g + 1), (n : ℂ))

\end{leancode}

\begin{dictionary}
\lean{layerDeg B F} places the layer $\mathcal L_p$ of the flag $F$ (orthogonal complements for $B$) in bidegree $(2p,0)$ and weight $n$;
\lean{trivialDeg} puts everything in $R=0$; \lean{tensorDeg} is the tensor-product grading (degrees add). \lean{GhostRGrading PG} is a
grading of the ghost space by $((2R,\mathrm{gh}),\text{weight})$ with pieces orthogonal for the ghost form \lean{PG}. The index type
\lean{(ℤ × ℤ) × ℂ} is $((2R,\mathrm{gh}),\text{weight})$; \lean{relativeWeight n} is the relative complex at weight $n$
(\leanpath{LeeYangTransferAssemblyComplex.lean}), and \lean{⊔} is the sum of subspaces.

\textbf{A subtlety.} The layers of the corrected flag on $A$ form a decomposition of $A_n$ only if $B$ is nondegenerate on every step, which is
clause (iv) of \lean{hA}; so the $R$-grading of the complex itself depends on \lean{hA} (\lean{matterDeg_span hA}). This is legitimate: the grading
is needed only under that hypothesis.
\end{dictionary}

\begin{physics}
Goodness is Beem--Garner's Definition~3.12: the gauge current has $R$-charge one, so each term of $c^aJ^a$ changes the $R$-charge by exactly
$\pm\tfrac12$ after the ghost's own charge is accounted for.
\end{physics}

\begin{status}
Definition (docstring still says ``OPEN''; it is proved, next theorem).
\end{status}

\begin{theorem}[The actual differential is good]\label{l28:thm:good}
Assume $\ltwentyeightGU(A)$ and that $B$ is unitary for some form and real structure. Then for every positive weight function $w$ on ghost modes,
$\ltwentyeightQ$ is good on relative states for the ghost grading built from $w$.
\end{theorem}

\noindent File \leanpath{VOA/\allowbreak Examples/\allowbreak BRSTGoodnessDifferential.lean}.
\begin{leancode}
theorem relativeDifferentialGood_of_affine
    (hA : GradedUnitary affineForm affineRealStructure affineFlag)
    {BB : InvariantForm BVOA} {RB : RealStructure BVOA} (hB : Unitary BB RB)
    (w : Finset ModeIdx → ℝ) (hw : ∀ s, 0 < w s) :
    RelativeDifferentialGood (ghostRGrading w) := by
\end{leancode}

\begin{dictionary}
\lean{ModeIdx} indexes ghost modes; \lean{ghostRGrading w} is the ghost $R$-grading with weights \lean{w}. The hypotheses \lean{hA} and \lean{hB}
are used only to have the matter grading at all (previous dictionary) and adjointness of current modes with respect to the layers.
\end{dictionary}

\paragraph{Proof on paper.} Refine the grading by the number of $c_0$'s and set $R'=R+\tfrac12n_{c_0}$. Every term of
$\ltwentyeightQ=\sum_a\sum_nJ^a_n\otimes c^a_{-n}+\ltwentyeightQ_{\mathrm{gh}}$ shifts $2R'$ by exactly $\pm1$: the matter modes because the flags
are good and the layers are adjoint-compatible, the $c$-modes by the sign of their index ($c_0$ included), and the cubic term by a direct
computation. Relative states contain no $c_0$ (since $\sum_ac_0^ab_0^a$ acts as the $c_0$-count and kills states annihilated by the $b^a_0$), so
$R'=R$ there, and $\ltwentyeightQ$ preserves relative states. \qed

\begin{physics}
The gauged chiral algebra inherits a well-defined $R$-filtration, and the BRST charge has definite $R$-charge $\pm\tfrac12$.
\end{physics}

\begin{status}
Proved; axioms standard (re-checked). Conditional on \lean{hA} and \lean{hB} as displayed; in the headline \lean{hB} is discharged by
Theorem~\ref{l28:thm:bunitary}.
\end{status}

\begin{definition}[The two ghost-conjugation identities]\label{l28:def:ghostconj}
Let $\rho=\rho_M\otimes\rho_G$ and $B=B_M\otimes B_G$ be the product conjugation and form on the BRST space, $\Lambda$ Beem--Garner's operator
(a weighted sum of ghost bilinears) and $E=[\Lambda,\ltwentyeightQ]$. The identities are: $\rho\,\ltwentyeightQ=(-1)^{\mathrm{gh}}E\,\rho$ and
$B(Eu,v)=i(-1)^{\mathrm{gh}(u)}B(u,Ev)$ on relative states.
\end{definition}

\noindent File \leanpath{VOA/\allowbreak Examples/\allowbreak GhostConjEndpoint.lean}.
\begin{leancode}
def ghConjBRST : BRSTSpace →ₗ⋆[ℂ] BRSTSpace :=
  TensorProduct.map (matterRho bRealStructure) ghostConj

def ghBBRST : LinearMap.BilinForm ℂ BRSTSpace := (matterB bForm).tmul ghostB

def GhostConjRhoD : Prop :=
  ∀ (n j g : ℤ), ∀ x ∈ relativeWeight n, x ∈ brstDeg bgGh ((j, g), (n : ℂ)) →
    ghConjBRST (brstQ x) = ((-1 : ℂ) ^ g) • brstE n (ghConjBRST x)

def GhostConjBE : Prop :=
  ∀ (n j g : ℤ), ∀ u ∈ relativeWeight n, u ∈ brstDeg bgGh ((j, g), (n : ℂ)) →
    ∀ v ∈ relativeWeight n,
      ghBBRST (brstE n u) v = Complex.I * (-1 : ℂ) ^ g * ghBBRST u (brstE n v)
\end{leancode}

\begin{dictionary}
\lean{matterRho bRealStructure} and \lean{matterB bForm} are the product conjugation and form of the three matter factors; \lean{ghostConj},
\lean{ghostB} those of the ghosts. \lean{brstQ} is $\ltwentyeightQ$ and \lean{brstE n}$=\Lambda\ltwentyeightQ-\ltwentyeightQ\Lambda$
(\leanpath{BGPackageReduced.lean}); \lean{bgGh} is the standard ghost $R$-grading. Both are statements about the BRST operator only: \lean{hA}
does not occur.
\end{dictionary}

\begin{physics}
$E$ is (up to normalization) the second supercharge $\ltwentyeightQ^+$, the $SU(2)_R$ partner of the BRST charge. The identities say that the
anti-unitary conjugation exchanges the two supercharges, and that $E$ is skew for the bilinear pairing with the phase $i$; together they make
$\ltwentyeightQ^+$ the adjoint partner required by Definition~\ref{l28:def:kahler}.
\end{physics}

\begin{status}
Definitions (docstrings: ``OPEN (numerically verified)''; both are proved, next theorem).
\end{status}

\begin{theorem}[The ghost-conjugation identities hold]\label{l28:thm:ghostconj}
Both identities of Definition~\ref{l28:def:ghostconj} hold.
\end{theorem}

\noindent Files \leanpath{VOA/\allowbreak Examples/\allowbreak BRSTConjIdentityProved.lean}, \leanpath{BRSTConjIdentityBE.lean}.
\begin{leancode}
theorem ghostConjRhoD_proved : GhostConjRhoD := ghostConjRhoD_of_cubic ghostCubicConj_proved

theorem ghostConjBE_proved : GhostConjBE := by
\end{leancode}

\begin{dictionary}
Exactly the two identities, with no hypotheses. The first is the composition of a cubic-term identity (\lean{ghostCubicConj_proved}) with the
reduction \lean{ghostConjRhoD_of_cubic}.
\end{dictionary}

\paragraph{Proof on paper.} Both sides are computed on the explicit ghost Fock basis tensored with matter: the matter part uses
$\rho_M J^a_n\rho_M^{-1}=-J^{a^*}_n$ and the adjointness of current modes, the ghost part the explicit action of $\rho_G$ and $\Lambda$ on
$b$- and $c$-modes. The cubic ghost term needs a separate computation (\leanpath{BRSTConjIdentityCubic*.lean}). The identities were first
checked numerically to $10^{-12}$ at relative weights $2,3,4$ (project computations outside the library); the Lean proof is independent of that. \qed

\begin{physics}
The gauged theory is $SU(2)_R$-covariant: CPT conjugation maps the BRST charge to its $R$-partner.
\end{physics}

\begin{status}
Proved outright; axioms standard (re-checked: \lean{ConcreteHodge.ghostConjRhoD_proved}, \lean{ConcreteHodge.ghostConjBE_proved}).
\end{status}

\begin{theorem}[The metric on the coset representatives]\label{l28:thm:repmetric}
Let $\mathrm{rep}(x)=\phi_{\mathrm{coset}}(x)\otimes\ltwentyeightvac_U\otimes\ltwentyeightvac_{\mathrm{gh}}$ be the explicit representative of a universal
Virasoro state $x\in V(-\tfrac{22}5,0)$. With $B=B_M\otimes B_G$ and $\rho=\rho_M\otimes\rho_G$: the vacuum representative has norm one;
$L_{-i-2}$ and $L_{i+2}$ are adjoint on representatives; and $\rho$ fixes the representative of every Virasoro word.
\end{theorem}

\noindent Files \leanpath{VOA/\allowbreak Examples/\allowbreak HarmonicMetricBridge.lean}, \leanpath{HarmonicMetricVacuum.lean},
\leanpath{HarmonicMetricAdjoint.lean}, \leanpath{HarmonicMetricReal.lean}.
\begin{leancode}
def repVec (x : VirUniv leeYangC) : BRSTSpace := cosetToComplex (virasoroToCoset x)

def RepVacuumNormalized : Prop := ghBBRST (repVec (vacVec leeYangC)) (repVec (vacVec leeYangC)) = 1

def RepCreationAdjoint : Prop :=
  ∀ (i : ℕ) (u v : VirUniv leeYangC),
    ghBBRST (repVec (Lop leeYangC (-(i : ℤ) - 2) u)) (repVec v) =
      ghBBRST (repVec u) (repVec (Lop leeYangC ((i : ℤ) + 2) v))

def RepWordsReal : Prop :=
  ∀ w : List ℕ, ghConjBRST (repVec (negWordVir leeYangC w)) = repVec (negWordVir leeYangC w)

theorem repVacuumNormalized_proved : RepVacuumNormalized := by

theorem repCreationAdjoint_proved : RepCreationAdjoint := by

theorem repWordsReal_proved : RepWordsReal := by
\end{leancode}

\begin{dictionary}
\lean{repVec x} is \lean{cosetToComplex (virasoroToCoset x)}, i.e.\ $\mathrm{rep}(x)$. \lean{vacVec leeYangC} is the vacuum of $V(c,0)$,
\lean{Lop c m} the Virasoro mode $L_m$ on $V(c,0)$, \lean{negWordVir c w} the creation word of the list $w$. Exactly the three statements.

\textbf{Why these three suffice.} The harmonic representative of a class differs from $\mathrm{rep}(x)$ by a boundary, and the form pairs harmonic
vectors trivially with boundaries (\leanpath{HarmonicMetricGeneric.lean}), so the three statements hold for harmonic representatives too
(\lean{harmonicVacuumNormalized_of_rep} and companions). A bilinear form on the weight spaces of $V(c,0)$ that is normalized on the vacuum and for which
$L_{-m}$ is adjoint to $L_m$ is the canonical form (\lean{LeeYangShortcut.weightForm_eq_Bsimp}, re-checked); and a conjugation fixing all words is
the partition-basis conjugation. These are hypotheses \lean{hB}, \lean{hh} of Theorem~\ref{l28:thm:classfamily}.
\end{dictionary}

\paragraph{Proof on paper.} The representative is a matter state tensored with vacua, so all three are statements about the form and conjugation of
$A\otimes B\otimes U$ on states of the form $x\otimes\ltwentyeightvac$. The vacuum norm is the product of the normalized vacuum norms. The coset
Virasoro modes are built from the modes of $\omega_A,\omega_B$ and the diagonal currents, whose adjoints are again of the same type
(invariance of each factor form), and the coset projection commutes with them. The conjugation fixes $\omega_{\mathrm{coset}}$ (it fixes each
Sugawara vector) and commutes with all modes, so it fixes every word. \qed

\begin{physics}
The inner product of Lee--Yang states computed in the gauged theory equals the intrinsic Lee--Yang (BPZ) inner product: the stress tensor's
Hermiticity survives gauging.
\end{physics}

\begin{status}
All proved outright; axioms standard (re-checked: the three \lean{rep*_proved}). The rung \lean{ConcreteHodge.leeYang_gradedUnitary_of_affine_levelOne}
assembles everything of this section; its remaining inputs are the invariant-symbol and character statements, treated next (re-checked:
standard axioms).
\end{status}

\begin{theorem}[The Virasoro classes are filtered]\label{l28:thm:filtered}
Assume $\ltwentyeightGU(A)$. On the actual complex, at every conformal weight $n$ and for every $p$, the cohomology class of a universal
Virasoro word of length $\le p$ (through the coset stress tensor, Theorem~\ref{l28:thm:cosetmap}) has a representative of $2R\in[0,2p]$ at
ghost number zero. This is step (b) of Remark~\ref{l28:rem:arc}.
\end{theorem}

\noindent Files \leanpath{VOA/\allowbreak Examples/\allowbreak LeeYangShortcutAssembly.lean}, \leanpath{InvariantSymbolBridge.lean}.
\begin{leancode}
def UniversalClassFiltered (S : ConcreteKahlerGeometry) : Prop :=
  ∀ n (hp : IsPosOn (S.geometry n).P ⊤) p,
    (universalPrefix p n).map (modelUniversalClass S n hp) ≤
      (S.model n hp).cohomFilt (KahlerComplex.rowSet (Finset.Icc 0 (2 * (p : ℤ))) 0)

abbrev endpointKahler (hA : GradedUnitary affineForm affineRealStructure affineFlag) :
    ConcreteKahlerGeometry :=
  sp2Kahler hA LevelOneUnitary.bUnitary bgPG_pos bgGh bgGh_isStd
    ((bgGhostRemainderI hA ghostConjRhoD_proved ghostConjBE_proved).pkg
      LevelOneUnitary.bUnitary).sp2

theorem endpointKahler_filtered (hA : GradedUnitary affineForm affineRealStructure affineFlag) :
    UniversalClassFiltered (endpointKahler hA) :=
\end{leancode}

\begin{dictionary}
\lean{ConcreteKahlerGeometry} is a weightwise Kähler geometry on the relative complex (Definition~\ref{l28:def:gaptwo}); \lean{S.model n hp} is
the Kähler complex at weight $n$ once the positivity \lean{hp} of the modified form is supplied. \lean{universalPrefix p n} is the span of the
universal Virasoro words of length $\le p$ in $V(c,0)_n$ (\leanpath{LeeYangShortcutRelative.lean}), \lean{modelUniversalClass S n hp} sends a
word to its cohomology class, and \lean{cohomFilt (rowSet (Icc 0 (2p)) 0)} is the span of the classes with a representative in bidegrees
$(j,0)$, $0\le j\le 2p$. \lean{endpointKahler hA} is the Kähler geometry actually used by the headline: the Beem--Garner package of this
section with the standard ghost grading \lean{bgGh}, the proved unitarity of $B$ and the proved ghost-conjugation identities. So
\lean{endpointKahler_filtered hA} is exactly the hypothesis \lean{hfiltered} of Theorem~\ref{l28:thm:classfamily} in the application. It is
conditional on \lean{hA} only because the complex's $R$-grading is built from \lean{hA} (Definition~\ref{l28:def:good}).
\end{dictionary}

\paragraph{Proof on paper.} The proof (\lean{universalClassFiltered_of_generators}, with the two generator tests of
\lean{interpretedGenerators_std}) is an induction on word length. The vacuum class is represented by the vacuum of the BRST space, which
sits in bidegree $(0,0)$ for the standard ghost grading. If the class of a word $x$ of length $\le p$ has a representative in $2R\le 2p$, then
the class of $L_{-i-2}x$ has one in $2R\le 2p+2$: the coset Virasoro mode is built from the Sugawara vectors of $A$ and $B$ and the
diagonal currents, each of corrected degree at most one (Definition~\ref{l28:def:corrected}), so it raises $R$ by at most one. This is where
the convention that $\omega$ costs one letter, not two, is used. Positivity of the modified form at every weight, needed to speak of the
Kähler complex at all, is \lean{concreteKahler_weakGU} applied to \lean{hA}, the unitarity of $B$ and \lean{auxiliaryGradedUnitary}. \qed

\begin{physics}
The coset stress tensor has $R$-charge one, like the four-dimensional stress tensor it models, so a product of $p$ of its modes has
$R$-charge at most $p$.
\end{physics}

\begin{status}
Proved; axioms standard (re-checked: \lean{ConcreteHodge.endpointKahler_filtered}). Conditional on \lean{hA}, as displayed.
\end{status}

\begin{warning}[title={Dead ends in the concrete package}]\label{l28:warn:package}
\begin{itemize}
\item \textbf{$c_0$ breaks goodness off relative states.} With $c_0$ placed at $R=0$, the cubic terms $c_0c_{-j}b_j$, $c_{-i}c_ib_0$ preserve $R$; the
  fix is the shifted $R'=R+\tfrac12n_{c_0}$ and the restriction to relative states (Theorem~\ref{l28:thm:good}).
\item \textbf{No global transposes of $b_0$ (conditional on \lean{hA}).} \lean{antighostZero_transpose_inconsistent}
  (\leanpath{BGPackageReduced.lean}) takes \lean{hA}, a unitarity datum \lean{hB : Unitary BB RB} for $B$, a positive ghost form with a standard
  ghost grading, and then a bilinear form on the BRST space for which the antighost zero modes have a global transpose of the expected sign and
  which is compatible with the positive form at weight zero; from all of these it concludes \lean{False} (it forces the vacuum to have norm
  zero). So the theorem proves ``$\ltwentyeightGU(A)$ implies that no such form exists''; if $\ltwentyeightGU(A)$ failed it would say nothing. Under
  the standing hypothesis of the headline, the package must therefore be built on $c_0$-free (relative) states.
\item \textbf{The full ghost current is not skew.} Once $c_0$ is present, the full ghost current is not skew for the ghost form (a pairing sum equals
  $4$ instead of $0$); the skewness lemma is proved on the projection to relative states instead.
\item \textbf{Colour-uniform ghost forms and the old remainder.} See Warning~\ref{l28:warn:vacuous}.
\end{itemize}
\end{warning}

%======================================================================
\section{The Koszul step: leading cohomology is a space of invariants}\label{l28:sec:koszul}
%======================================================================

At ghost number zero the leading differential $Q^-$ has two parts: a Koszul differential $d_K$ that contracts the auxiliary creation modes of
$U$ against the antighosts $b^a_{-n}$, and a Chevalley--Eilenberg part $d_{CE}$ that inserts a ghost $c^a_n$ ($n\ge0$) together with the
positive mode $J^a_n$ of $\ltwentyeightgr A\otimes B$. Since the auxiliary $U$ is freely generated (PBW), $d_K$ is exact in positive antighost
number, and the spectral sequence of the $c$-count collapses onto one row: $H^0(Q^-)$ is the kernel of the positive-mode action on
$\ltwentyeightgr A\otimes B$.

\begin{theorem}[Leading cohomology equals invariants]\label{l28:thm:koszul}
Let $C$ be a Kähler complex whose pieces of bidegrees $a-(1,1)$, $a$, $a+(1,1)$ are modelled by the diagonals of a Koszul--CE bicomplex, with
Koszul exactness in positive antighost degree, and suppose the zeroth Koszul homologies of the two lowest rows are identified with $N$ and
$N^\kappa$ so that the induced CE map is $v\mapsto(\rho_kv)_k$. Then $H(Q^-)_a\cong\bigcap_k\ker\rho_k$, and harmonic vectors of bidegree $a$ have
dimension at most $\dim\bigcap_k\ker\rho_k$.
\end{theorem}

\noindent File \leanpath{VOA/\allowbreak Examples/\allowbreak LeadingKoszulKahler.lean}.
\begin{leancode}
def invariants (ρ : κ → N →ₗ[R] N) : Submodule R N := ⨅ k, LinearMap.ker (ρ k)

structure LeadModel where
  eX : C.deg a ≃ₗ[ℂ] (∀ p, X p)
  eY : C.deg (a - shQm) ≃ₗ[ℂ] (∀ p, Y p)
  eZ : C.deg (a + shQm) →ₗ[ℂ] (∀ p, Z p)
  eZ_inj : Function.Injective eZ
  hQX : ∀ (x : C.deg a) (h : C.Qm x ∈ C.deg (a + shQm)), eZ ⟨C.Qm x, h⟩ = B.D0 (eX x)
  hQY : ∀ (y : C.deg (a - shQm)) (h : C.Qm y ∈ C.deg a), eX ⟨C.Qm y, h⟩ = B.Dm (eY y)

noncomputable def leadCohomEquivInvariants {N : Type*} [AddCommGroup N] [Module ℂ N]
    {κ : Type*}
    (e0 : B.KH0 ≃ₗ[ℂ] N) (e1 : B.KH1 ≃ₗ[ℂ] (κ → N)) (ρ : κ → N →ₗ[ℂ] N)
    (hρ : ∀ v k, e1 (B.dbar v) k = ρ k (e0 v)) : C.leadCohom a ≃ₗ[ℂ] invariants ρ :=

theorem finrank_harmDeg_le_invariants [FiniteDimensional ℂ K] (M : LeadModel C B a)
    {N : Type*} [AddCommGroup N] [Module ℂ N] {κ : Type*}
    (e0 : B.KH0 ≃ₗ[ℂ] N) (e1 : B.KH1 ≃ₗ[ℂ] (κ → N)) (ρ : κ → N →ₗ[ℂ] N)
    (hρ : ∀ v k, e1 (B.dbar v) k = ρ k (e0 v)) :
    finrank ℂ (C.harmDeg a) ≤ finrank ℂ (invariants ρ) :=
\end{leancode}

\begin{dictionary}
\lean{B : LeadingBicomplex ℂ X Y Z G} is the bicomplex: \lean{X p}$=T^{p,p}$ (ghost number $0$), \lean{Y p}$=T^{p+1,p}$ ($-1$), \lean{Z p}$=T^{p,p+1}$
($+1$), \lean{G p}$=T^{p,p+2}$, where $T^{p,q}$ has $p$ antighosts and $q$ ghosts; its fields are the six components of $d_K$, $d_{CE}$, the used
components of $(d_K+d_{CE})^2=0$, and Koszul exactness at two positions (\leanpath{LeadingKoszulBicomplex.lean}). \lean{B.D0}, \lean{B.Dm} are the total
differentials out of and into ghost number $0$ (on the \emph{product} totalization \lean{∀ p, X p}). \lean{LeadModel} identifies the three
bidegrees of the Kähler complex with these totalizations, intertwining $Q^-$. \lean{B.KH0}, \lean{B.KH1} are the zeroth Koszul homologies of rows
$0$ and $1$; \lean{B.dbar} the map induced by $d_{CE}$. \lean{invariants ρ} is $\bigcap_k\ker\rho_k$. The \lean{finrank} statement additionally takes
\lean{M : LeadModel C B a}.
\end{dictionary}

\paragraph{Proof on paper.} (\lean{LeadingBicomplex.equivKer}.) A ghost-number-zero cocycle $(x_p)_p$ with $x_p\in T^{p,p}$ has leading component
$x_0$ whose class in $H^K_0(\text{row }0)$ is killed by $\bar d_{CE}$; conversely, given such a class, Koszul exactness lets one solve for
$x_1,x_2,\dots$ in turn (a zigzag), and two solutions differ by a coboundary, again by exactness. So $H^0\cong\ker\bar d_{CE}$. With the
identifications, $\ker\bar d_{CE}=\{v\in N:\rho_kv=0\ \forall k\}$. The dimension bound combines this with Theorem~\ref{l28:thm:harmbound}. \qed

\begin{physics}
The ``Gauss law'' at leading order: after the auxiliary sector and the antighosts cancel in pairs (Koszul), the surviving gauge-invariant states
are those annihilated by the positive gauge-current modes.
\end{physics}

\begin{status}
Proved; axioms standard (re-checked: \lean{LeadingKoszul.LeadingBicomplex.equivKer}). Unconditional for the abstract datum.
\end{status}

\begin{theorem}[Koszul exactness and the symbol maps on the actual complex]\label{l28:thm:kexact}
On the actual relative complex at every weight $N$ at which the modified Hermitian form is positive definite (hypothesis \lean{hp}; under
$\ltwentyeightGU(A)$ this holds at every weight), the blocks of $Q^-$ by antighost number satisfy Koszul exactness
in positive antighost degree, and the two symbol maps from the antighost-free blocks onto $\ltwentyeightgr_rA\otimes B$ exist, with kernels the
$d_K$-boundaries and intertwining $d_{CE}$ with the positive-mode action.
\end{theorem}

\noindent Files \leanpath{VOA/\allowbreak Examples/\allowbreak KoszulResidualExact.lean}, \leanpath{KoszulResidualSymbolMaps.lean}.
\begin{leancode}
theorem kExact_endpoint {n : ℤ} (hp : IsPosOn ((endpointKahler hA).geometry n).P ⊤)
    (a : ℤ × ℤ) : (endpointBlocks hA n hp).KExact a :=

def koszulResidual_endpoint : EndpointKoszulResidual hA :=
  koszulResidual_of_symbols hA (symbolResidual_endpoint hA)
\end{leancode}

\begin{dictionary}
\lean{endpointKahler hA} is the actual Kähler geometry (assembled from the package of Section~\ref{l28:sec:package}), \lean{endpointBlocks hA n hp}
its block decomposition by antighost number, and \lean{KExact a} the two exactness statements of the bicomplex at bidegree $a$
(\leanpath{KoszulConcreteBlocks.lean}). \lean{EndpointKoszulResidual hA} bundles, for every weight and degree, Koszul exactness and the two symbol
maps $\pi_0,\pi_1$ (\leanpath{KoszulConcreteEndpoint.lean}); \lean{hA} is a section variable. \textbf{Hypotheses.} \lean{kExact_endpoint}
takes \lean{hA} and, explicitly, \lean{hp : IsPosOn ((endpointKahler hA).geometry n).P ⊤}, positivity of the modified form at weight $n$ (the
Kähler complex at weight $n$ is only defined once \lean{hp} is given). \lean{hp} is not an extra assumption in the headline: it is discharged at
every weight from \lean{hA}, the unitarity of $B$ and \lean{auxiliaryGradedUnitary} by \lean{concreteKahler_weakGU} (as in the proof of
Theorem~\ref{l28:thm:filtered}). \lean{koszulResidual_endpoint} is a \lean{def} whose type is a structure of proofs (\lean{EndpointKoszulResidual hA}),
not a \lean{theorem}; constructing it is the proof. \lean{hA} itself is needed only because the complex's grading is built from it.
\end{dictionary}

\paragraph{Proof on paper.} The auxiliary factor $U$ is the universal affine VOA, whose PBW basis (Chapter~\ref{ch:L02}, now a theorem) makes it a free
polynomial algebra in the creation modes; $d_K$ is the Koszul differential of the regular sequence of these modes against the antighosts, hence exact
in positive degree, uniformly at each weight (\leanpath{KoszulResidualPBW.lean}, \leanpath{KoszulResidualHomotopy.lean}). The symbol maps send a
$b$-free state to its $U$-vacuum component's class in $\ltwentyeightgr A$. \qed

\begin{physics}
The auxiliary current algebra is ``pure gauge'': its oscillators cancel against the antighosts, leaving the gauge-invariant content of $A\otimes B$.
\end{physics}

\begin{status}
Proved, conditional on \lean{hA} and (for \lean{kExact_endpoint}) on the weightwise positivity \lean{hp}, which itself follows from \lean{hA};
axioms standard (re-checked: \lean{KoszulResidual.kExact_endpoint}, \lean{KoszulResidual.koszulResidual_endpoint}).
\end{status}

\begin{definition}[The gap-two coefficient and the invariant-symbol statements]\label{l28:def:gaptwo}
$G_r(n)$ is the coefficient of $z^n$ in $z^{r(r+1)}/\prod_{j=1}^r(1-z^j)$, i.e.\ the number of partitions of $n$ into exactly $r$ parts, all $\ge2$,
with consecutive parts differing by at least $2$. In the Kähler coordinate $j=2R$ put $\mathrm{lc}(n,j)=G_{j/2}(n)$ for even $j\ge0$ and $0$
otherwise. The \emph{invariant-symbol bound} for a concrete Kähler geometry says $\dim H(Q^-)_{(j,0)}\le\mathrm{lc}(n,j)$ at every weight $n$; the
\emph{dimension} statement says equality. The \emph{Lee--Yang PBW character} says $\dim F_pL_n=\sum_{j=0}^{2p}\mathrm{lc}(n,j)$.
\end{definition}

\noindent File \leanpath{VOA/\allowbreak Examples/\allowbreak LeeYangTransferAssemblyCharacter.lean}.
\begin{leancode}
def gapTwoCoefficient (n : ℤ) (r : ℕ) : ℕ := by
  classical
  exact Fintype.card {a : Fin r → Fin (n.toNat + 1) //
    ((r * (r + 1) + ∑ j : Fin r, (j.val + 1) * (a j).val : ℕ) : ℤ) = n}

def leadingCoefficient (n j : ℤ) : ℕ :=
  if 0 ≤ j ∧ j % 2 = 0 then gapTwoCoefficient n (j.toNat / 2) else 0

def LeeYangPBWCharacter : Prop :=
  ∀ (n : ℤ) (p : ℕ),
    Module.finrank ℂ ((virGradedFlag leeYangC).pieceIn p n) =
      ∑ j ∈ Finset.Icc 0 (2 * (p : ℤ)), leadingCoefficient n j
\end{leancode}
\noindent Files \leanpath{VOA/\allowbreak Examples/\allowbreak LeeYangTransferAssemblyInputs.lean}, \leanpath{InvariantSymbolBridge.lean}.
\begin{leancode}
def LeadingInvariantSymbolDimension (S : ConcreteKahlerGeometry) : Prop :=
  ∀ n (hp : IsPosOn (S.geometry n).P ⊤) j,
    Module.finrank ℂ ((S.model n hp).leadCohom (j, 0)) = leadingCoefficient n j

def LeadingInvariantSymbolBound (S : ConcreteKahlerGeometry) : Prop :=
  ∀ n (hp : IsPosOn (S.geometry n).P ⊤) j,
    Module.finrank ℂ ((S.model n hp).leadCohom (j, 0)) ≤ leadingCoefficient n j
\end{leancode}

\begin{dictionary}
\lean{gapTwoCoefficient n r} counts tuples $(a_1,\dots,a_r)$ with $r(r+1)+\sum_j j\,a_j=n$: $a_j$ is the multiplicity of the part $j$ in a partition
into parts $\le r$, and adding $r(r+1)=2+4+\dots+2r$ turns it into a gap-two partition with $r$ parts (the standard bijection). The bound
\lean{Fin (n.toNat + 1)} on multiplicities is harmless because $j\,a_j\le n$. \lean{ConcreteKahlerGeometry} is a weightwise Kähler geometry on
\lean{relativeWeight n} whose differential is the actual BRST differential and whose ghost coordinate is the actual ghost number
(\leanpath{LeeYangTransferAssemblyInputs.lean}); \lean{S.model n hp} is the Kähler complex once positivity \lean{hp} is supplied. In the headline the
geometry is the concrete \lean{endpointKahler hA}; nothing is free. The docstrings still say ``OPEN''; see Section~\ref{l28:sec:defects}.
\end{dictionary}

\begin{physics}
$\sum_rG_r(n)z^n$ is the generating function $\sum_rz^{r(r+1)}/(z)_r$ of the second Rogers--Ramanujan identity, $=\prod_{n\ge0}(1-z^{5n+2})^{-1}(1-z^{5n+3})^{-1}$,
the Lee--Yang vacuum character; $G_r$ refines it by the number of stress-tensor modes, i.e.\ by $R$-charge.
\end{physics}

\begin{status}
Definitions. \lean{gapTwoCoefficient_vacuum}: $G_0(0)=1$. \lean{LeeYangPBWCharacter} is proved (Theorem~\ref{l28:thm:pbwchar}); the
symbol bound is proved for the actual geometry inside the headline's proof (Remark~\ref{l28:rem:assemble}).
\end{status}

\begin{theorem}[The cyclic-evaluation bound]\label{l28:thm:cyclic}
Let $V$ be a vector space with an injective linear map $\mathrm{ev}:V\to\C[a_1,a_2,\dots]$ (oscillators, $a_m$ of weight $m$) landing in weighted
degree $n-r(r+1)$, whose image is annihilated, under the pairing $\langle F,P\rangle=F(2\partial_{a_1},4\partial_{a_2},\dots)P$, by every element of
matching degree of the Cartan ideal $I_r=\ker\bigl(\C[H_1,H_2,\dots]\to\C[s_1,\dots,s_r]^{S_r},\ H_m\mapsto2\sum_is_i^m\bigr)$. Then
$\dim V\le G_r(n)$.
\end{theorem}

\noindent File \leanpath{VOA/\allowbreak Examples/\allowbreak InvariantSymbolCyclic.lean}.
\begin{leancode}
theorem finrank_le_gapTwoCoefficient_of_cyclicEvaluation
    {V : Type*} [AddCommGroup V] [Module ℂ V] (N : ℤ) (r : ℕ)
    (ev : V →ₗ[ℂ] MvPolynomial ℕ ℂ) (hinj : Function.Injective ev)
    (hdeg : ∀ x, IsWeightedHomogeneous modeWeightZ (ev x) (N - r * (r + 1)))
    (hann : ∀ x (F : MvPolynomial ℕ ℂ), F ∈ cartanIdeal r →
      IsWeightedHomogeneous modeWeightZ F (N - r * (r + 1)) → modePairing F (ev x) = 0) :
    finrank ℂ V ≤ LeeYangTransferAssembly.gapTwoCoefficient N r := by
\end{leancode}

\begin{dictionary}
\lean{MvPolynomial ℕ ℂ} with variable $m$ of weight $m+1$ (\lean{modeWeightZ}) models both $\C[H_1,H_2,\dots]$ and $\C[a_1,a_2,\dots]$.
\lean{IsWeightedHomogeneous modeWeightZ P d} says $P$ has weighted degree $d$. \lean{cartanIdeal r} is $I_r$ and \lean{modePairing} the pairing.
Exactly the theorem.
\end{dictionary}

\paragraph{Proof on paper.} The pairing is perfect in each weighted degree, so the annihilator of $(I_r)_d$ has dimension
$\dim(\C[H]/I_r)_d$. The quotient $\C[H]/I_r$ is the image of $H_m\mapsto2\sum_is_i^m$, which is the whole ring of symmetric polynomials in $r$
variables (power sums generate it over $\C$); its degree-$d$ part has dimension the number of partitions of $d$ into at most $r$ parts, equivalently
into parts $\le r$. With $d=n-r(r+1)$ this is $G_r(n)$. An injective map into that
annihilator bounds $\dim V$. \qed

\begin{physics}
In the level-one lattice picture, a charge-$r$ invariant symbol is detected by its overlap with the $r$-fermion (charge $r$) Fock sector; the
Cartan relations express that $r$ fermions have only symmetric functions of their $r$ positions as degrees of freedom.
\end{physics}

\begin{status}
Proved outright; axioms standard (re-checked). Unconditional.
\end{status}

\begin{remark}[How inequality (e) is assembled]\label{l28:rem:assemble}
To apply Theorem~\ref{l28:thm:cyclic} to $V=H(Q^-)_{(2r,0)}\cong(\ltwentyeightgr_rA\otimes B)^{\ltwentyeightsl[t]}_n$ (Theorem~\ref{l28:thm:koszul}) the library
builds $\mathrm{ev}$ as follows. (i) A \emph{cyclic dual vector} on $\ltwentyeightgr_rA$: a functional $\beta$ such that $x\ne0$ implies
$\beta(J^{a_1}_{m_1}\cdots J^{a_\ell}_{m_\ell}x)\ne0$ for some word of non-negative modes. It is obtained from strictness of the Adamović symbol map
(Section~\ref{l28:sec:jet}), \lean{GrJetStrict}, and the proved cyclicity of the jet ring (Theorem~\ref{l28:thm:jetcyclic}). (ii) Heisenberg separation of
the charge sectors of $B=L_1$ (\lean{levelOneHeisenbergSeparation}), proved from two facts: in charge $2r$ the bottom energy $r^2$ is one-dimensional
(\lean{levelOneBottomDetected}), and Heisenberg-singular vectors of charge $2r$ sit at energy $r^2$ (\lean{levelOneHeisenbergSingular}), the latter using
$\omega=\tfrac14h_{-1}^2\ltwentyeightvac$. (iii) Combining, an invariant $\sum_ix_i\otimes y_i$ is evaluated by pairing with the cyclic vector on the $A$ side
and the Cartan monomials on the $B$ side, and invariance under $\ltwentyeightsl[t]$ forces the Cartan relations (\leanpath{CyclicEval*.lean}). All
three are proved; the assembled result is that \lean{LeadingInvariantSymbolBound (endpointKahler hA)} holds, inside the proof of the rung
\lean{SymbolModel.leeYang_gradedUnitary_of_affine_jetStrict}.
\end{remark}

%======================================================================
\section{The jet ring and the Adamovi\'c symbol}\label{l28:sec:jet}
%======================================================================

The invariant-symbol bound needs a concrete model of $\ltwentyeightgr_rA$, the associated graded of the corrected filtration, together with its
$\ltwentyeightsl[t]$-action. Adamović's free-field realization of $L_{-4/3}(\ltwentyeightsl)$ suggests the model: $e\mapsto u^2$, $f\mapsto-v^2$,
$h\mapsto-2uv$, and the Sugawara field $\mapsto-(u'v-uv')$, for two commuting formal series $u(s),v(s)$. The library does not use a free-field
realization; it proves directly that this assignment defines an isomorphism on associated graded pieces.

\begin{definition}[The jet ring and its image $A_r$]\label{l28:def:jet}
Let $\C[u_j,v_j]_{j\ge0}$ be the polynomial ring in two sequences of variables and $u(s)=\sum_ju_js^j$, $v(s)=\sum_jv_js^j$. Put
$X=u^2$, $Y=uv$, $Z=v^2$, $J=u'v-uv'$ (power series in $s$ with polynomial coefficients), and let $A_r$ be the span of all products of $r$ coefficients
of $X,Y,Z,J$. The conformal grading of $A_r$ puts $u_j,v_j$ in degree $j$ and adds $r$. The loop algebra $\ltwentyeightsl[s]$ acts by derivations:
$e_m\colon v\mapsto s^mu$, $f_m\colon u\mapsto s^mv$, $h_m\colon u\mapsto s^mu,\ v\mapsto-s^mv$ (and $e_mu=f_mv=0$).
\end{definition}

\noindent Files \leanpath{VOA/\allowbreak Examples/\allowbreak SymbolModelJet.lean}, \leanpath{SymbolModelJetImage.lean},
\leanpath{SymbolModelJetCyclic.lean}.
\begin{leancode}
abbrev JetPoly : Type := MvPolynomial (Fin 2 × ℕ) ℂ

def jetDer (k : Fin 3 × ℕ) : Derivation ℂ JetPoly JetPoly := mkDerivation ℂ (jetDerVal k)

def colSer (b : Fin 2) : JetSer := PowerSeries.mk fun j => X (b, j)

def xSer : JetSer := colSer 0 * colSer 0
def ySer : JetSer := colSer 0 * colSer 1
def zSer : JetSer := colSer 1 * colSer 1
def jSer : JetSer := dS (colSer 0) * colSer 1 - colSer 0 * dS (colSer 1)

def genSpan : Submodule ℂ JetPoly :=
  Submodule.span ℂ (Set.range fun p : Fin 4 × ℕ => genCoeff p.1 p.2)

def jetImage (r : ℕ) : Submodule ℂ JetPoly := genSpan ^ r

def jetDeg (r : ℕ) (d : ℤ) : Submodule ℂ (jetImage r) :=
  (weightedHomogeneousSubmodule ℂ jetWt (d - r)).comap (jetImage r).subtype

\end{leancode}

\begin{dictionary}
\lean{MvPolynomial (Fin 2 × ℕ) ℂ}: variables $u_j=$\lean{X (0, j)}, $v_j=$\lean{X (1, j)}. \lean{JetSer} is power series in $s$ with coefficients in
the jet ring, \lean{dS} is $d/ds$. \lean{genCoeff g a} is the $s^a$-coefficient of $X,Y,Z,J$ for $g=0,1,2,3$. \lean{genSpan ^ r} is the $r$-th power
of a subspace in the algebra (the span of $r$-fold products). \lean{jetDer k} is the derivation of $\C[u_j,v_j]$ determined by its values on
variables (\lean{jetDerVal}, which encodes the formulas above with \lean{k = (a, m)}). \lean{jetDeg r d} is the subspace of $A_r$ of jet weight
$d-r$, i.e.\ conformal weight $d$.
\end{dictionary}

\begin{physics}
$u,v$ are the two components of an $\ltwentyeightsl$ doublet of commuting fields (a ``$\beta\gamma$-type'' pair seen classically); $X,Y,Z$ are the
classical (arc-space) images of the currents and $J$ is the image of the stress tensor. $A_r$ is the classical phase space of $r$ letters.
\end{physics}

\begin{status}
Definitions.
\end{status}

\begin{definition}[Word symbols; descent; strictness]\label{l28:def:symbols}
The \emph{symbol} of a corrected letter $(g,m)$ (Definition~\ref{l28:def:corrected}) is the $s^m$-coefficient of $X$, $-Z$, $-2Y$ for $g=e,f,h$ and of
$-J$ for $g=\omega$; the symbol of a word is the product. The \emph{descent} statement in degree $r$ asks for a linear map $\sigma$ on $F_rA$ vanishing
on $F_{r-1}A$ and sending the class of each length-$r$ word to its symbol. \emph{Strictness} in degree $r$ asks for an injective, degree-preserving map
$\ltwentyeightgr_rA\to A_r$ intertwining the non-negative modes $J^a_m$ with the derivations of Definition~\ref{l28:def:jet}; a \emph{leading-symbol
map} is the same with ``surjective'' in place of ``injective''.
\end{definition}

\noindent Files \leanpath{VOA/\allowbreak Examples/\allowbreak AdamovicSymbolJet.lean}, \leanpath{AdamovicSymbolDescent.lean},
\leanpath{JetStrictSqueeze.lean}, \leanpath{SymbolModelJetCyclic.lean}, \leanpath{SymbolModelGr.lean}.
\begin{leancode}
def curSer : Fin 3 → JetSer
  | 0 => xSer
  | 1 => -zSer
  | 2 => -(ySer + ySer)

def letterSym : Generator × ℕ → JetPoly
  | (.current a, n) => PowerSeries.coeff n (curSer a)
  | (.sugawara, n) => PowerSeries.coeff n (-jSer)

def wordSym (w : List (Generator × ℕ)) : JetPoly := (w.map letterSym).prod

structure WordSymbolDescends (r : ℕ) where
  σ : aFil r →ₗ[ℂ] JetPoly
  lower : ∀ x : aFil r, (x : ASpace) ∈ aFilLT r → σ x = 0
  word : ∀ (w : List (Generator × ℕ)) (hw : w.length = r), σ (wordClass w hw.le) = wordSym w

structure JetLeadingSymbol (r : ℕ) where
  g : GrA r →ₗ[ℂ] jetImage r
  surjective : Function.Surjective g
  intertwine : ∀ k, g ∘ₗ grRho r k = jetImgRho r k ∘ₗ g
  degree : ∀ d, ∀ m ∈ grDeg r d, g m ∈ jetDeg r d

structure GrJetStrict (r : ℕ) where
  f : GrA r →ₗ[ℂ] jetImage r
  injective : Function.Injective f
  intertwine : ∀ k, f ∘ₗ grRho r k = jetImgRho r k ∘ₗ f
  degree : ∀ d, ∀ m ∈ grDeg r d, f m ∈ jetDeg r d

def aFil (p : ℕ) : Submodule ℂ ASpace := ⨆ m : ℤ, affineFlag.piece p m

def grKer (r : ℕ) : Submodule ℂ (aFil r) := (aFilLT r).comap (aFil r).subtype

abbrev GrA (r : ℕ) : Type := aFil r ⧸ grKer r
\end{leancode}

\begin{dictionary}
\lean{aFil p} is $F_pA=\sum_n F_pA_n$ (all weights at once), \lean{aFilLT r} is $F_{r-1}A$, \lean{GrA r} is $F_rA/F_{r-1}A=\ltwentyeightgr_rA$.
\lean{grRho r k} is the induced action of $J^a_m$, $m\ge0$, on $\ltwentyeightgr_rA$ (well defined because the flag is good), \lean{jetImgRho r k}
the derivation restricted to $A_r$, and \lean{grDeg r d} the weight-$d$ part of $\ltwentyeightgr_rA$. \lean{wordClass w hw} is the class in $F_rA$ of the
corrected word $w$; \lean{∘ₗ} is composition of linear maps. A \lean{structure} with data fields (\lean{σ}, \lean{f}, \lean{g}) is a \lean{Type}, not a
\lean{Prop}: having an element means having the map with its properties.

\textbf{Why ``descent'' is the hard part.} The words of length $r$ span $F_rA$ modulo $F_{r-1}A$ but satisfy relations in $A$ (the maximal ideal is
nonzero: there is a singular vector at weight $3$). Descent says that every relation among words, modulo shorter words, is a relation among their
symbols. This is the content of a free-field realization, and it is where the paper proof cites Adamović.
\end{dictionary}

\begin{physics}
Descent is the classical limit of the free-field realization: it says that the Schur operators of $(A_1,D_3)$, at leading order in the $R$-filtration,
are polynomials in two classical fields $u,v$, and that every null relation of the chiral algebra holds classically.
\end{physics}

\begin{status}
Definitions (the docstrings say ``OPEN''; all three are realized below for every $r$).
\end{status}

\begin{theorem}[Cyclicity of the jet ring]\label{l28:thm:jetcyclic}
Let $\mathrm{ev}$ be evaluation at the base point $(u,v)=(0,1)$, i.e.\ $u_j\mapsto0$, $v_0\mapsto1$, $v_j\mapsto0$ ($j\ge1$). If a jet polynomial $Q$ satisfies
$\mathrm{ev}(D_{k_1}\cdots D_{k_\ell}Q)=0$ for every word $D_{k_1}\cdots D_{k_\ell}$ in the non-negative-mode derivations, then $Q=0$.
\end{theorem}

\noindent File \leanpath{VOA/\allowbreak Examples/\allowbreak SymbolModelJet.lean}.
\begin{leancode}
theorem jet_cyclic (Q : JetPoly)
    (hQ : ∀ l : List (Fin 3 × ℕ), evalBase (word jetRho l Q) = 0) : Q = 0 := by
\end{leancode}

\begin{dictionary}
\lean{word jetRho l} is the composite of the operators \lean{jetRho k}, $k\in l$ (\leanpath{CyclicEvalCore.lean}); \lean{evalBase} is evaluation at
\lean{baseVal}, which is $1$ at $v_0$ and $0$ elsewhere. Exactly the theorem.
\end{dictionary}

\paragraph{Proof on paper.} Shift $v_0\mapsto v_0+1$, so that the base point becomes the origin. The derivations become derivations whose action on the
variables has a ``Taylor'' form: from $f_0u=v=1+\cdots$, words of derivations can produce any monomial's dual, so a polynomial with all iterated
derivatives vanishing at the origin has all coefficients zero (\lean{exists_taylor}). The library's proof is an algebraic Taylor argument replacing
the orbit-and-density argument of the source note. \qed

\begin{physics}
The orbit of the point $(u,v)=(0,1)$ under the loop group $SL_2[[s]]$ is dense in the space of jets: a classical observable that vanishes along the
whole orbit vanishes identically.
\end{physics}

\begin{status}
Proved outright; axioms standard (re-checked). The base point is $(0,1)$, not the $(1,0)$ of the source note; this matches the charge convention of the
dual vectors used later (file docstring).
\end{status}

\begin{theorem}[The Adamović symbols descend, in every degree]\label{l28:thm:desc}
For every noncritical $k\notin\N$ and every VOA ideal $I$ of $V^k(\ltwentyeightsl)$ that avoids conformal weight $0$: if $y\in I$ and $(y,Q)$ is in the
span of the pairs $(\text{word},\text{symbol})$ of length-$r$ corrected words, modulo $F_{r-1}V^k\times0$, then $Q=0$. Consequently, at $k=-4/3$ the
word symbols descend to $\ltwentyeightgr_rA$ for every $r$.
\end{theorem}

\noindent Files \leanpath{VOA/\allowbreak Examples/\allowbreak AdamovicSymbolWords.lean}, \leanpath{AdamovicDescCyclic.lean}.
\begin{leancode}
def symMod (q : ℕ) : Submodule ℂ (WakModule k × JetPoly) :=
  symPairs k hk q ⊔ (lowerP k hk q).prod ⊥

theorem symbol_eq_zero_of_mem (hkN : ∀ i : ℕ, k ≠ i)
    (hI : (affineSl2VOA_closed k hk).IsIdeal I) (hIN : I ≤ offZero k hk) (r : ℕ)
    {y : WakModule k} {Q : JetPoly} (hyQ : (y, Q) ∈ symMod k hk r) (hy : y ∈ I) : Q = 0 :=

def wordSymbolDescends (r : ℕ) : WordSymbolDescends r where
\end{leancode}

\begin{dictionary}
\lean{symPairs k hk q} is the span of the pairs $(\text{state of }w,\ \text{symbol of }w)$ over words of length $q$, and \lean{lowerP k hk q} is
$F_{q-1}V^k$; \lean{symMod} is their sum (\lean{⊔}) in $V^k\times\C[u_j,v_j]$, with \lean{.prod ⊥} meaning ``paired with $0$''. In
\lean{symbol_eq_zero_of_mem} the section variables are \lean{k}, \lean{hk : k + 2 ≠ 0} and an ideal \lean{I}; \lean{hkN : ∀ i : ℕ, k ≠ i} says
$k\notin\N$; \lean{hI} says $I$ is a VOA ideal; \lean{hIN : I ≤ offZero k hk} says $I$ is contained in $\bigoplus_{n\ne0}V^k_n$. The conclusion
\lean{Q = 0} is the single-valuedness of the symbol modulo $I$. \lean{wordSymbolDescends r} \emph{constructs} the descent datum of
Definition~\ref{l28:def:symbols} in every degree $r$, at $k=-4/3$ and $I$ the maximal ideal: the map $\sigma$ is the graph of the relation (fields
proved; we show only the header).
\end{dictionary}

\paragraph{Proof on paper.} (1) \emph{A detecting functional.} Put $E=e_1$. For $(y,Q)$ in the span, $E^ry\equiv(-1)^rc_r\,\mathrm{ev}(Q)\ltwentyeightvac$ modulo
$\bigoplus_{n\ne0}V_n$, with $c_r=\prod_{i<r}(i+1)(k-i)$. On a word, $\mathrm{ev}(\text{symbol})$ is $(-1)^r$ for $f_{-1}^r\ltwentyeightvac$ and $0$ otherwise, and
$E^rf_{-1}^r\ltwentyeightvac=c_r\ltwentyeightvac$ by the $\ltwentyeightsl$-triple $(e_1,h_0+k,f_{-1})$; every other word is sent off weight $0$ or into weight $0$ with
nonzero $h_0$-charge, where it vanishes since $V_0=\C\ltwentyeightvac$. (2) \emph{Detection on the ideal.} If $y\in I$ then $E^ry\in I\subseteq\bigoplus_{n\ne0}V_n$,
so $\mathrm{ev}(Q)=0$; $c_r\ne0$ because $k\notin\N$. (3) \emph{Cyclicity.} The span is stable under the pairs of operators $(J^a_\mu,D_{(a,\mu)})$ (a mode on the
state, the derivation on the symbol), and $I$ is stable under $J^a_\mu$; so $\mathrm{ev}(D_{k_1}\cdots D_{k_\ell}Q)=0$ for every word, and
Theorem~\ref{l28:thm:jetcyclic} gives $Q=0$. (4) \emph{Descent.} The image of the span in $(V/I)\times\C[u_j,v_j]$ is therefore the graph of a linear map on
$F_rA$, which is the descent datum. \qed

\begin{physics}
The vacuum expectation value of $e_1^r$ against a state detects its ``all-$f$'' component; invariance under the current algebra then propagates this single
measurement to all of the state's classical shadow. This replaces the free-field realization that the physics literature (Adamović) uses.
\end{physics}

\begin{status}
Proved outright; axioms standard (re-checked: \lean{AdamovicSymbol.Desc.symbol_eq_zero_of_mem}, \lean{AdamovicSymbol.Desc.wordSymbolDescends}). This is the
last input discharged before the headline.
\end{status}

\begin{example}[The singular vector has vanishing degree-two symbol]\label{l28:ex:singular}
At $k=-4/3$, $V^k$ has the weight-3 singular vector
$20e_{-3}-12e_{-2}h_{-1}-6h_{-2}e_{-1}+36e_{-1}e_{-1}f_{-1}+9e_{-1}h_{-1}h_{-1}$ (on $\ltwentyeightvac$; Chapter~\ref{ch:L08}, Theorem~\ref{l08:thm:sing43}).
Rewriting $2e_{-1}f_{-1}-h_{-2}+\tfrac12h_{-1}^2=\tfrac43L_{-2}$ (the Sugawara relation at $k=-4/3$; a hand computation, not formalized) gives the corrected
form $20e_{-3}-12e_{-2}h_{-1}-6h_{-2}e_{-1}+24e_{-1}L_{-2}+18e_{-1}h_{-2}$, whose length-two part must have zero symbol if descent holds. It does:
\par\noindent File \leanpath{VOA/\allowbreak Examples/\allowbreak AdamovicSymbolDescent.lean}.
\begin{leancode}
theorem singular_symbol_vanishes :
    -12 • wordSym [(.current 0, 1), (.current 2, 0)]
      - 6 • wordSym [(.current 2, 1), (.current 0, 0)]
      + 24 • wordSym [(.current 0, 0), (.sugawara, 0)]
      + 18 • wordSym [(.current 0, 0), (.current 2, 1)] = (0 : JetPoly) := by
\end{leancode}
In jet terms this is the identity $XJ=X'Y-XY'$ at $s=0$. A sanity check, proved by computation.
\end{example}

\begin{definition}[The jet character bound and the fermionic series]\label{l28:def:jetbound}
The \emph{jet character bound} asks, for every $N$, that $[q^N]\bigl((q^3;q^3)_\infty/(q;q)_\infty\bigr)^3\le\sum_{r\le N}\dim A_{r,N}$. The \emph{fermionic jet
series} is $\sum_{r\ge0}q^r\bigl(\sum_{k=0}^{2r}\left[\begin{smallmatrix}2r\\k\end{smallmatrix}\right]_q\bigr)/(q)_r$, and the \emph{fermionic identity} says it equals
$\bigl((q^3;q^3)_\infty/(q;q)_\infty\bigr)^3$.
\end{definition}

\noindent Files \leanpath{VOA/\allowbreak Examples/\allowbreak JetStrictEndpoint.lean}, \leanpath{JetProductCount.lean}, \leanpath{JetProductEndpoint.lean},
\leanpath{AffineSl2BoundaryVacuumCharacter.lean}.
\begin{leancode}
def JetProductBound : Prop :=
  ∀ N : ℕ, PowerSeries.coeff N (AffineSl2BoundaryCharacter.product 1) ≤
    ∑ r ∈ Finset.range (N + 1), (Module.finrank ℂ (jetDeg r (N : ℤ)) : ℚ)

def jetTerm (r : ℕ) : PowerSeries ℚ :=
  PowerSeries.X ^ r * (∑ k ∈ Finset.range (2 * r + 1), gaussian PowerSeries.X (2 * r) k) *
    reciprocalQFactorial r

def jetSeries : PowerSeries ℚ := VOA.LocallyFinitePowerSeriesSum.sum jetTerm

def JetFermionicIdentity : Prop := jetSeries = AffineSl2BoundaryCharacter.product 1

def product (m : ℕ) : PowerSeries ℚ :=
  (PowerSeries.expand (denominator m) (by dsimp [denominator]; omega) euler) ^ 3 *
    reciprocalEuler ^ 3
\end{leancode}

\begin{dictionary}
\lean{PowerSeries.coeff N} is $[q^N]$. \lean{product m} is $(\mathrm{expand}_{2m+1}\,\phi)^3\cdot\phi^{-3}$ with $\phi=(q;q)_\infty$ (\lean{euler}) and
\lean{reciprocalEuler}$=\phi^{-1}$; for $m=1$ (\lean{denominator 1}$=3$) it is $((q^3;q^3)_\infty/(q;q)_\infty)^3$. \lean{gaussian PowerSeries.X n k} is the
Gaussian binomial, \lean{reciprocalQFactorial r} is $1/(q)_r$, and \lean{LocallyFinitePowerSeriesSum.sum} is the sum of a family of power series of which
only finitely many contribute to each coefficient. The coefficients are in $\Q$ (exact rationals).
\end{dictionary}

\begin{physics}
$((q^3;q^3)_\infty/(q;q)_\infty)^3$ is the vacuum character of $L_{-4/3}(\ltwentyeightsl)$ (Chapter~\ref{ch:L20}), i.e.\ the Schur index of the $(A_1,D_3)$
theory. The fermionic series counts the same states by number of letters: a fermionic (quasi-particle) formula.
\end{physics}

\begin{status}
Definitions (docstring of \lean{JetProductBound}: ``OPEN''; it is proved, next theorem).
\end{status}

\begin{theorem}[The fermionic jet identity, and strictness]\label{l28:thm:jetfermionic}
(1) $\sum_{r\ge0}q^r\sum_{k=0}^{2r}\left[\begin{smallmatrix}2r\\k\end{smallmatrix}\right]_q/(q)_r=\bigl((q^3;q^3)_\infty/(q;q)_\infty\bigr)^3$.
(2) $[q^N]$ of the fermionic series is at most $\sum_{r\le N}\dim A_{r,N}$; hence the jet character bound holds. (3) Given leading-symbol maps in every degree
$r\ge1$ and the jet character bound, strictness holds in every degree $r\ge1$; with Theorem~\ref{l28:thm:desc} this gives strictness outright.
\end{theorem}

\noindent Files \leanpath{VOA/\allowbreak Examples/\allowbreak JetFermionicIdentity.lean}, \leanpath{JetProductCount.lean},
\leanpath{JetStrictEndpoint.lean}.
\begin{leancode}
theorem jetFermionicIdentity : JetFermionicIdentity := by

theorem coeff_jetSeries_le (N : ℕ) :
    PowerSeries.coeff N jetSeries ≤
      ∑ r ∈ Finset.range (N + 1), (Module.finrank ℂ (jetDeg r (N : ℤ)) : ℚ) := by

def grJetStrict_of_leadingSymbol_product (sym : ∀ r, JetLeadingSymbol (r + 1))
    (chi : JetProductBound) (r : ℕ) : GrJetStrict (r + 1) :=
  grJetStrict_succ_of_leadingSymbol sym (jetCharacterBound_of_product chi) r
\end{leancode}

\begin{dictionary}
Exactly (1), (2) and (3). \lean{grJetStrict_of_leadingSymbol_product} produces \lean{GrJetStrict (r + 1)} for every $r$ from \lean{sym} and \lean{chi};
it goes through \lean{jetCharacterBound_of_product}, which uses the Arakawa--Dai character theorem \lean{AffineSl2BoundaryCharacter.character_eq_product 1}
(Chapter~\ref{ch:L20}, Theorem~\ref{l20:thm:character}): $\dim A_N=[q^N]((q^3;q^3)_\infty/(q;q)_\infty)^3$. The leading-symbol maps come from the descent data
by \lean{WordSymbolDescends.toJetLeadingSymbol} (surjective because word symbols span $A_r$), with degrees $0$ and $1$ done by hand
(\lean{jetLeadingSymbol_one}).
\end{dictionary}

\paragraph{Proof on paper.} (1) Multiply both sides by $(q;q)_\infty^3$. The left side becomes a double sum which a Wilf--Zeilberger certificate
(\lean{wzSum_eq}) and a moment identity (\lean{euler_cube_mul_pairSum}) reduce to Jacobi's identity for $(q;q)_\infty^3$ (\lean{jacobi_cube}); the right side is
$(q^3;q^3)_\infty^3$. It is the case $k=1$, $t=q$, $z=1$ of a theorem of Chern, Tran and Wakhare combined with the Jacobi cube; the Lean proof uses no Bailey
pairs. (2) $A_r$ is a free module over a ring of symmetric functions with fibre of dimension $4^r$ (\lean{block_indep}, \lean{local_indep}), which gives a lower
bound for $\dim A_{r,N}$ by the $r$-th term of the series. (3) The squeeze: a leading-symbol map $g_r:\ltwentyeightgr_rA_N\twoheadrightarrow A_{r,N}$ gives
$\sum_r\dim A_{r,N}\le\sum_r\dim\ltwentyeightgr_rA_N=\dim A_N=[q^N]((q^3;q^3)/(q;q))^3\le\sum_r\dim A_{r,N}$, so each $g_r$ is bijective. \qed

\begin{physics}
The number of Schur operators of $(A_1,D_3)$ built from $r$ letters equals the number of classical jet polynomials of degree $r$: the $R$-filtered
Schur index is computed by the classical theory, with no quantum corrections to the counting.
\end{physics}

\begin{status}
(1) and (2) proved outright; (3) proved from its hypotheses, and those are discharged in the headline. Axioms standard (re-checked:
\lean{JetFermionic.jetFermionicIdentity}).
\end{status}

\begin{warning}[title={Dead ends for strictness}]
Recorded in the project notes, with counterexamples: (i) the subalgebra generated by $X,Y,Z$ alone is too small (dimension $8$ against $9$ at weight $2$), so
the arc-space or Li-filtration route cannot give strictness and the generator $J$ (the stress tensor) is essential; (ii) a leading-monomial basis of $A_r$
overcounts ($78$ against $81$ at weight $5$); (iii) $A_r$ is not free over $\C[X_a,Z_a]$; (iv) there is no naive multiplicity-$\le2$ monomial basis, because
$X_0Z_0=Y_0^2$. These are hand or numerical findings, not machine-checked statements.
\end{warning}

%======================================================================
\section{The Lee--Yang vacuum character}\label{l28:sec:character}
%======================================================================

\begin{definition}[The two sides of the character identity]\label{l28:def:rr}
The \emph{Rogers--Ramanujan lower bound} for Lee--Yang says $\sum_rG_r(N)\le\dim L_N$ for every $N$. The \emph{Rocha-Caridi} number is
$\mathrm{rc}(N)=\sum_{k\in\Z}(-1)^kp\bigl(N-\tfrac{k(5k+3)}2\bigr)$; the \emph{Rocha-Caridi lower bound} says $\mathrm{rc}(N)\le\dim L_N$, and the
\emph{Rocha-Caridi character} says equality. The \emph{Gram-rank bound} says $\mathrm{rc}(N)\le\operatorname{rank}G_N$, where $G_N$ is the $p(N)\times p(N)$
Shapovalov matrix of the Verma module $M(-\tfrac{22}5,0)$ at level $N$.
\end{definition}

\noindent Files \leanpath{VOA/\allowbreak Examples/\allowbreak LeeYangCharacterReduction.lean}, \leanpath{RRLowerLimit.lean},
\leanpath{RRLowerAndrews.lean}, \leanpath{RRLowerReduction.lean}, \leanpath{RRLowerGram.lean}.
\begin{leancode}
def LeeYangRRLowerBound : Prop :=
  ∀ N : ℕ, (∑ r ∈ Finset.range (N + 1), gapTwoCoefficient (N : ℤ) r) ≤
    Module.finrank ℂ ((virasoroSimple leeYangC).grading (N : ℤ))

def rc (N : ℕ) : ℤ := ∑ k ∈ Ico (-(N : ℤ) - 1) ((N : ℤ) + 1), sgn k * pp ((N : ℤ) - ee k)

def ee (k : ℤ) : ℤ := k * (5 * k + 3) / 2

def sgn (k : ℤ) : ℤ := if k % 2 = 0 then 1 else -1

def LeeYangRochaCaridiLowerBound : Prop :=
  ∀ N : ℕ, rc N ≤ (Module.finrank ℂ ((virasoroSimple leeYangC).grading (N : ℤ)) : ℤ)

def LeeYangRochaCaridiCharacter : Prop :=
  ∀ N : ℕ, (Module.finrank ℂ ((virasoroSimple leeYangC).grading (N : ℤ)) : ℤ) = rc N

def LeeYangGramRankBound : Prop :=
  ∀ N : ℕ, rc N ≤ ((vermaGramMatrix leeYangC 0 N).rank : ℤ)
\end{leancode}

\begin{dictionary}
\lean{Ico (-N - 1) (N + 1)} is the integer interval $[-N-1,N+1)$; terms outside it vanish. \lean{sgn k}$=(-1)^k$, \lean{ee k}$=k(5k+3)/2$ (an integer),
\lean{pp M}$=p(M)$, zero for $M<0$ (\leanpath{RRLowerQBinomial.lean}). \lean{vermaGramMatrix c h N} (namespace \lean{InverseDS.VirasoroVerma}) is the explicit Shapovalov matrix of $M(c,h)$ in the PBW
basis at level $N$; Gram matrices of Verma modules are introduced in Chapter~\ref{ch:L06}, Definition~\ref{l06:def:gramverma}. \lean{virasoroSimple leeYangC} is $L$. Exactly the
definitions; the casts \lean{(... : ℤ)} are only there to compare naturals with integers.

\textbf{Sanity check.} $\mathrm{rc}(N)=p(N)-p(N-1)-p(N-4)+p(N-7)+p(N-13)-\cdots$; at $N=6$: $11-7-2=2$, at $N=8$: $22-15-5+1=3$, matching Example~\ref{l28:ex:lylow}.
\end{dictionary}

\begin{physics}
$\mathrm{rc}(N)$ is the Rocha-Caridi character of the $(2,5)$ vacuum module: the Verma character minus the two embedded singular submodules at levels $1$ and $4$,
plus the next ones, and so on (the Feigin--Fuchs resolution). $\sum_rG_r(N)$ is the fermionic (Rogers--Ramanujan) form of the same character.
\end{physics}

\begin{status}
Definitions (all docstrings: ``OPEN''; all are proved, Theorem~\ref{l28:thm:lychar}).
\end{status}

\begin{theorem}[Gap-two spanning and the upper bound]\label{l28:thm:gapspan}
In $L$, the length-$\le p$ step of the PBW flag at weight $N$ is spanned by the gap-two words of length $\le p$ (words $L_{-n_1}\cdots L_{-n_q}\ltwentyeightvac$ with
$n_1\ge n_2+2\ge\dots$ and $n_q\ge2$). Hence $\dim F_pL_N\le\sum_{j=0}^{2p}\mathrm{lc}(N,j)$.
\end{theorem}

\noindent Files \leanpath{VOA/\allowbreak Examples/\allowbreak LeeYangCharacterSpanning.lean}, \leanpath{LeeYangCharacterReduction.lean}.
\begin{leancode}
def GapTwo (w : List ℕ) : Prop := w.IsChain (fun a b => b + 2 ≤ a)

noncomputable def gapSpan (ℓ N : ℕ) : Submodule ℂ (VirUniv leeYangC) :=
  Submodule.span ℂ
    {x | ∃ u : List ℕ, GapTwo u ∧ energyVir u = N ∧ u.length ≤ ℓ ∧ negWordVir leeYangC u = x}

theorem goodFil_leeYang_eq_gapSpan (p N : ℕ) :
    goodFil leeYangC p (N : ℤ) = (gapSpan p N).map (maxIdealVir leeYangC).mkQ := by

theorem finrank_pieceIn_le (n : ℤ) (p : ℕ) :
    Module.finrank ℂ ((virGradedFlag leeYangC).pieceIn p n) ≤
      ∑ j ∈ Finset.Icc 0 (2 * (p : ℤ)), leadingCoefficient n j := by
\end{leancode}

\begin{dictionary}
A word is a list $w=[i_1,\dots,i_q]$ standing for $L_{-(i_1+2)}\cdots L_{-(i_q+2)}\ltwentyeightvac$ (\lean{negWordVir}, Chapter~\ref{ch:L06}); \lean{GapTwo w} says
consecutive entries drop by at least $2$ (\lean{List.IsChain}). \lean{goodFil leeYangC p N} is the image of the length-$\le p$ words in $L_N$
(Definition~\ref{l19:def:lengthflag}); \lean{(maxIdealVir leeYangC).mkQ} the quotient map. Exactly the theorem.
\end{dictionary}

\paragraph{Proof on paper.} Straightening: modulo the maximal ideal, which contains the level-4 singular vector and hence the relation
$L_{-2}^2\equiv\tfrac35L_{-4}$ and all its descendants, any product $L_{-a}L_{-b}$ with $a-b\in\{0,1\}$ inside a word can be rewritten as a combination of words
that are smaller in a suitable order (the sum of squares of the parts strictly increases, \lean{sq_spread}), without increasing the length. Induction on that
order. The count of gap-two words of length $\le p$ is $\sum_{r\le p}G_r(N)$. \qed

\begin{physics}
The Lee--Yang null vector $(L_{-2}^2-\tfrac35L_{-4})\ltwentyeightvac$ is the ``exclusion rule'' of a quasi-particle description: two stress-tensor modes cannot
be closer than distance two.
\end{physics}

\begin{status}
Proved outright; axioms standard (re-checked: \lean{LeeYangCharacter.goodFil_leeYang_eq_gapSpan}, \lean{LeeYangCharacter.finrank_pieceIn_le}).
\end{status}

\begin{theorem}[Rogers--Ramanujan in Rocha-Caridi form]\label{l28:thm:andrews}
For every $N$, $\sum_{r\le N}G_r(N)=\mathrm{rc}(N)$. Hence $\dim L_N\le\mathrm{rc}(N)$.
\end{theorem}

\noindent File \leanpath{VOA/\allowbreak Examples/\allowbreak RRLowerReduction.lean}.
\begin{leancode}
theorem sum_gapTwoCoefficient_eq_rc (N : ℕ) :
    ((∑ r ∈ Finset.range (N + 1), gapTwoCoefficient (N : ℤ) r : ℕ) : ℤ) = rc N := by

theorem finrank_le_rc (N : ℕ) :
    (Module.finrank ℂ ((virasoroSimple leeYangC).grading (N : ℤ)) : ℤ) ≤ rc N := by
\end{leancode}

\begin{dictionary}
Exactly the theorem; the left side is cast to $\Z$. The second statement is the case $p=N$ of Theorem~\ref{l28:thm:gapspan} rewritten with the first.
\end{dictionary}

\paragraph{Proof on paper.} Andrews' polynomial identity: for each $n$, a finite sum of Gaussian binomials over $j$ (the gap-two side, truncated) equals an
alternating sum of Gaussian binomials (the Rocha-Caridi side, truncated). Let $n\to\infty$ coefficientwise: for $n=7N+7$ every Gaussian coefficient that contributes to
$q^N$ has stabilized, giving $\sum_jp_{\le j}(N-j(j+1))=\sum_k(-1)^kp(N-k(5k+3)/2)$, which is the claim. No Jacobi triple product and no infinite product are needed. \qed

\begin{physics}
The second Rogers--Ramanujan identity, in the form ``fermionic sum $=$ bosonic alternating sum'': the two standard ways of writing the Lee--Yang vacuum
character agree.
\end{physics}

\begin{status}
Proved outright; axioms standard (re-checked).
\end{status}

\begin{lemma}[Corank is at most the order of vanishing]\label{l28:lem:corank}
Let $P(c)$ be a square matrix of polynomials whose value at $c_0$ is symmetric. Then $(c-c_0)^{\dim\ker P(c_0)}$ divides $\det P(c)$.
\end{lemma}

\noindent File \leanpath{VOA/\allowbreak Examples/\allowbreak LYGramRank.lean}.
\begin{leancode}
theorem pow_dvd_det_of_ker {ρ : Type} [Fintype ρ] (P : Matrix ρ ρ ℂ[X]) (c0 : ℂ)
    (hsym : (P.map (Polynomial.eval c0)).transpose = P.map (Polynomial.eval c0)) :
    ∃ Q : ℂ[X], P.det =
      (X - Polynomial.C c0) ^ Module.finrank ℂ (LinearMap.ker (P.map (Polynomial.eval c0)).mulVecLin) * Q := by
\end{leancode}

\begin{dictionary}
\lean{ℂ[X]} is the polynomial ring in one variable (here the variable is called \lean{X}); \lean{P.map (Polynomial.eval c0)} is $P(c_0)$; \lean{.mulVecLin} the
associated linear map; \lean{∃ Q, ... = (X - C c0)^k * Q} is divisibility. Exactly the lemma.
\end{dictionary}

\paragraph{Proof on paper.} Choose a basis whose first $k$ vectors span $\ker P(c_0)$. In that basis (a constant change of basis, which changes $\det$ by a nonzero
constant) the first $k$ columns of $P(c)$ vanish at $c_0$, so each is divisible by $c-c_0$; pull out $(c-c_0)^k$. \qed

\begin{physics}
The Jantzen-type principle: along a one-parameter deformation of the central charge, null states can only appear where the Kac determinant vanishes, and at
least to the order of the number of null states.
\end{physics}

\begin{status}
Proved; axioms standard (re-checked as part of \lean{leeYangGramRankBound_holds}).
\end{status}

\begin{theorem}[The Lee--Yang vacuum character]\label{l28:thm:lychar}
$\mathrm{rc}(N)\le\operatorname{rank}G_N\le\dim L_N$ for every $N$; hence $\dim L_N=\mathrm{rc}(N)=\sum_rG_r(N)$, and the Rogers--Ramanujan lower bound holds.
\end{theorem}

\noindent Files \leanpath{VOA/\allowbreak Examples/\allowbreak RRLowerGram.lean}, \leanpath{LYGramRank.lean}.
\noindent File \leanpath{VOA/Examples/LYGramRankCharacter.lean}.
\begin{leancode}
theorem gramRank_le_finrank (N : ℕ) :
    (vermaGramMatrix leeYangC 0 N).rank ≤
      Module.finrank ℂ ((virasoroSimple leeYangC).grading (N : ℤ)) := by

theorem leeYangGramRankBound_holds : LeeYangGramRankBound := by

theorem leeYangRochaCaridiCharacter_holds : LeeYangRochaCaridiCharacter :=

theorem leeYangRRLowerBound_holds : LeeYangCharacter.LeeYangRRLowerBound :=
\end{leancode}

\begin{dictionary}
No hypotheses. \lean{gramRank_le_finrank}: the irreducible quotient $L(c,0)$ of the Verma module has dimension $\operatorname{rank}G_N$ at level $N$
(Theorem~\ref{l06:thm:dimrank}), and it is a composition factor of the simple VOA's weight space with multiplicity one, via the local composition series
(\lean{virLocalCompositionSeries}). The three \lean{_holds} theorems are the definitions of Definition~\ref{l28:def:rr} (and of \lean{LeeYangRRLowerBound}).
\end{dictionary}

\paragraph{Proof on paper.} Deform along the curve $c(t)=13-6t-6/t$, $h=0$, which passes through $c=-22/5$ at $t_0=2/5$; on the whole curve $L_{-1}\ltwentyeightvac$
is singular. Choose, at $t_0$, a complement $f$ of the descendants $L_{-\lambda}L_{-1}\ltwentyeightvac$ at level $1+m$, and let $D(c)$ be the Gram matrix of $f$. A
factorization identity of the Kac determinant along the curve (\lean{FFLower.curve_identity}) expresses $\det D(c(t))$ through Kac factors whose orders of vanishing
at $t_0$ are computed exactly (\leanpath{LYGramOrders.lean}: \lean{aLY m}, \lean{bLY m}). Lemma~\ref{l28:lem:corank} bounds $\dim\ker D(c_0)$ by that order, so
$\operatorname{rank}G_{1+m}\ge\operatorname{rank}D(c_0)\ge|f|-\mathrm{ord}=p(1+m)-p(m)-\mathrm{aLY}(m)+\mathrm{bLY}(m)$, and the exact count
\lean{rc_eq_orders} identifies the right side with $\mathrm{rc}(1+m)$. Combined with Theorem~\ref{l28:thm:andrews} this gives equality. \qed

\begin{physics}
This is Feigin--Fuchs/Rocha-Caridi for the simplest non-unitary minimal model, proved from the Kac determinant alone: the vacuum module of Lee--Yang has exactly the
null states that the Kac determinant predicts along a generic deformation, and no accidental ones.
\end{physics}

\begin{status}
Proved outright; axioms standard (re-checked: \lean{RRLower.leeYangGramRankBound_holds}, \lean{leeYangRochaCaridiCharacter_holds},
\lean{leeYangRRLowerBound_holds}). This is a result of independent interest: the vacuum character of $M(2,5)$.
\end{status}

\begin{theorem}[The Lee--Yang PBW character, filtration level by filtration level]\label{l28:thm:pbwchar}
For every weight $N$ and every $p$, $\dim F_pL_N=\sum_{s\le p}G_s(N)$; that is, the statement \lean{LeeYangPBWCharacter} of
Definition~\ref{l28:def:gaptwo} holds. It is equivalent to the Rogers--Ramanujan lower bound.
\end{theorem}

\noindent File \leanpath{VOA/\allowbreak Examples/\allowbreak LeeYangCharacterReduction.lean}.
\begin{leancode}
theorem leeYangPBWCharacter_iff_lowerBound : LeeYangPBWCharacter ↔ LeeYangRRLowerBound :=
\end{leancode}

\begin{dictionary}
Both sides are the \lean{Prop}s of Definitions~\ref{l28:def:gaptwo} and \ref{l28:def:rr}; the equivalence has no hypotheses. Since
\lean{leeYangRRLowerBound_holds} proves the right side (Theorem~\ref{l28:thm:lychar}),
\lean{leeYangPBWCharacter_iff_lowerBound.mpr leeYangRRLowerBound_holds} proves \lean{LeeYangPBWCharacter}; this is the input
\lean{hcharacter} of Theorem~\ref{l28:thm:classfamily} at the level of $L$, i.e.\ step (f) of Remark~\ref{l28:rem:arc}. The statements use
\lean{Module.finrank}, which is $0$ on an infinite-dimensional space; here every weight space of $L$ is finite-dimensional because the
finitely many gap-two words span it (Theorem~\ref{l28:thm:gapspan}), so \lean{finrank} is the true dimension.
\end{dictionary}

\paragraph{Proof on paper.} ($\Rightarrow$) Take $p=N$: $F_NL_N=L_N$, so $\dim L_N=\sum_{s\le N}G_s(N)$. ($\Leftarrow$) By
Theorem~\ref{l28:thm:gapspan} the gap-two words of weight $N$ span $L_N$, and there are $\sum_sG_s(N)$ of them; the lower bound says
$\dim L_N$ is at least that number, so they are linearly independent, hence a basis. Then the gap-two words of length $\le p$ are linearly
independent, and they span $F_pL_N$ (Theorem~\ref{l28:thm:gapspan} again), so $\dim F_pL_N$ is their number $\sum_{s\le p}G_s(N)$. \qed

\begin{physics}
The gap-two words are a quasi-particle basis of the Lee--Yang vacuum module, and the length of a word is its $R$-charge: the refined
character counts Lee--Yang states by $R$-charge.
\end{physics}

\begin{status}
Proved outright; axioms standard (re-checked: \lean{LeeYangCharacter.}\allowbreak\lean{leeYangPBWCharacter_iff_lowerBound}). With
Theorem~\ref{l28:thm:lychar}, \lean{LeeYangPBWCharacter} is proved, unconditionally.
\end{status}

\begin{warning}[title={Dead ends for the Lee--Yang character}]
(i) Maps \emph{out of} $V(c,0)$ (including into Fock spaces) only bound $\dim L$ from above. (ii) A single Jantzen deformation or the Kac determinant along a
generic line is too weak: a generic line overshoots by $\dim(M(7)+M(13))$; the special curve $h=h_{1,1}$, on which $L_{-1}\ltwentyeightvac$ stays singular, is
essential. (iii) A 5-adic leading-order argument fails at $N=10$, and a modular-ODE route was not completed. (Project records; not machine-checked.)
\end{warning}


%======================================================================
\section{The files of this chapter, described collectively}\label{l28:sec:infra}
%======================================================================

The table lists the file families of this chapter with their role. ``Infrastructure'' means: gradings on tensor products, mode
computations, bookkeeping of weights, restrictions of operators to subspaces, and the intermediate rungs of the ladder. Every family is in the
import closure of the umbrella \lean{VOA} (checked by parsing the import graph of the state this chapter was written against).
\begin{center}\small
\begin{longtable}{p{0.30\textwidth}p{0.06\textwidth}p{0.56\textwidth}}
\toprule
family & files & role\\
\midrule\endhead
\leanpath{Sl2SemiInfinite*} & 34 & ghosts, ghost currents at level $4$, BRST current and charge, $\ltwentyeightQ^2=0\iff k=-4$, relative subcomplex, gradings (Section~\ref{l28:sec:brst})\\
\leanpath{Sl2Coset*} & 22 & the three factors, diagonal action, coset subspace and VOA, coset-to-cohomology map, vacuum class (Sections~\ref{l28:sec:brst}, \ref{l28:sec:coset})\\
\leanpath{GradedUnitaryCohomology*} & 4 & abstract Kähler complex, harmonic bound, GU data, the quartet example (Section~\ref{l28:sec:hodge})\\
\leanpath{LeeYangCosetTransfer*} & 4 & Hodge isomorphism, dimension squeeze, filtered-isomorphism transfer (Sections~\ref{l28:sec:hodge}, \ref{l28:sec:squeeze})\\
\leanpath{LeeYangTransferAssembly*} & 9 & the concrete affine and auxiliary data, the relative weight complex, the superseded assembly (Sections~\ref{l28:sec:endpoints}, \ref{l28:sec:headline})\\
\leanpath{LeeYangShortcut*} & 12 & Virasoro image, uniqueness of the normalized form, class-family theorem, shortcut assembly (Sections~\ref{l28:sec:coset}, \ref{l28:sec:squeeze})\\
\leanpath{ConcreteGood*}, \leanpath{BRSTGoodness*} & 8+8 & interpreted $R$-grading of the actual complex and its goodness (Section~\ref{l28:sec:package})\\
\leanpath{ConcreteHodge*}, \leanpath{BRSTKahler*}, \leanpath{BGPackage*} & 6+5+6 & product form, Kähler identities, $\Lambda$, the Beem--Garner package and its earlier variants (Section~\ref{l28:sec:package})\\
\leanpath{GhostConj*}, \leanpath{BRSTConjIdentity*} & 5+11 & ghost conjugation and the two conjugation identities (Section~\ref{l28:sec:package})\\
\leanpath{HarmonicMetric*} & 6 & metric on harmonic and explicit representatives (Section~\ref{l28:sec:package})\\
\leanpath{LevelOneUnitary*}, \leanpath{LevelOneHeis*} & 4+3 & unitarity of $L_1$; Heisenberg separation of its charge sectors (Sections~\ref{l28:sec:package}, \ref{l28:sec:koszul})\\
\leanpath{LeadingKoszul*}, \leanpath{KoszulConcrete*}, \leanpath{KoszulResidual*} & 7+3+18 & abstract Koszul--CE collapse; block decomposition of the actual complex; Koszul exactness and symbol maps (Section~\ref{l28:sec:koszul})\\
\leanpath{InvariantSymbol*}, \leanpath{CyclicEval*} & 3+6 & cyclic-evaluation bound and its assembly with the level-one facts (Section~\ref{l28:sec:koszul})\\
\leanpath{SymbolModel*}, \leanpath{JetStrict*}, \leanpath{JetProduct*}, \leanpath{JetFermionic*}, \leanpath{Adamovic*} & 8+2+7+3+5 & jet ring, cyclicity, strictness squeeze, jet-image lower bound, fermionic identity, Adamović symbols and their descent (Section~\ref{l28:sec:jet})\\
\leanpath{LeeYangCharacter*}, \leanpath{RRLower*}, \leanpath{LYGram*} & 5+7+4 & gap-two spanning, Andrews' identity, Rocha-Caridi, Kac-determinant order count (Section~\ref{l28:sec:character})\\
\bottomrule
\end{longtable}
\end{center}
No file of this chapter contains \lean{sorry}, \lean{admit} or \lean{native_decide} (a text search finds these words only in two docstrings, as ordinary English).

%======================================================================
\section{Verification record}\label{l28:sec:verification}
%======================================================================

\textbf{Records of earlier builds.} When the headline first compiled, the library still had the project axiom \lean{VOA.pbw}, and
\texttt{\#print axioms} of the headline was \lean{[propext, Classical.choice, Quot.sound, VOA.pbw]}; \lean{VOA.pbw} entered through
\lean{Desc.wordSymbolDescends}. The aggregate \texttt{lake build} was green then and in later aggregate builds that include these files
(13\,543, 13\,672 and 13\,868 jobs); these are records of the project's builds, not re-run for this edition. The headline is in the
library's maintained axiom audit. Since \lean{VOA.pbw} became a theorem, the axiom list is the standard three.

\textbf{What was re-run for this edition.} One scratch file importing \lean{VOA.Examples.AdamovicDescCyclic},
\lean{LeeYangTransferAssembly}, \lean{LeeYangShortcutAssembly}, \lean{InvariantSymbolBridge}, \lean{LeeYangCharacterReduction},
\lean{LYGramRank}, \lean{Sl2SemiInfiniteAnomalyCriterion}, \lean{Sl2SemiInfiniteGhostAffine}, \lean{LeeYangCosetTransferFiltered},
\lean{AffineSl2BoundaryVacuumCharacter} and the Chapter~\ref{ch:L29} modules \lean{TwoSeriesLeeYang}, \lean{TwoSeriesKoszulLeeYang} and
\lean{TwoSeriesGenIncl}, run as \texttt{lake env lean -DmaxSynthPendingDepth=3} from the repository root against the build products already
in the checkout (no \texttt{lake build} was run). For each of the following 55 declarations it printed
\lean{depends on axioms: [propext, Classical.choice, Quot.sound]}, with no error and no \lean{sorryAx}:
\begin{plaincode}
VOA.AdamovicSymbol.leeYang_gradedUnitary_of_affine
VOA.LeeYangTransferAssembly.leeYang_gradedUnitary_of_affine
VOA.AdamovicSymbol.Desc.wordSymbolDescends
VOA.AdamovicSymbol.Desc.symbol_eq_zero_of_mem
VOA.JetFermionic.jetFermionicIdentity
VOA.RRLower.leeYangRRLowerBound_holds
VOA.RRLower.leeYangRochaCaridiCharacter_holds
VOA.RRLower.leeYangGramRankBound_holds
VOA.RRLower.sum_gapTwoCoefficient_eq_rc
VOA.RRLower.finrank_le_rc
VOA.LevelOneUnitary.bUnitary
VOA.LevelOneHeis.bConformal_eq
VOA.LeeYangTransferAssembly.auxiliaryGradedUnitary
VOA.AffineSl2BelowCritical.universal_gradedUnitary_belowCritical_unconditional
VOA.Sl2CosetConcrete.leeYangBRSTComplex
VOA.Sl2CosetConcrete.leeYangBRSTComplex_cohomology_vacuum_ne_zero
VOA.Sl2CosetConcrete.cosetVOA
VOA.Sl2CosetConcrete.cosetCentralCharge
VOA.Sl2CosetConcrete.relativeVacuum_survives
VOA.GradedUnitaryCohomology.KahlerComplex.mem_harm_iff
VOA.GradedUnitaryCohomology.KahlerComplex.harmToLead_injective
VOA.GradedUnitaryCohomology.KahlerComplex.finrank_cohomFilt_row_le_lead
VOA.GradedUnitaryCohomology.GUData.flag_nondegenerate
VOA.GradedUnitaryCohomology.GUData.layer_signed_pos
VOA.LeeYangCosetTransfer.maps_eq_of_finrank_squeeze
VOA.LeeYangShortcut.gradedUnitaryAt_of_class_family
VOA.LeeYangShortcut.virasoroToH0_ker_le_maxIdeal
VOA.LeeYangShortcut.weightForm_eq_Bsimp
VOA.LeeYangShortcut.leeYang_gradedUnitary_of_affine'
VOA.ConcreteHodge.relativeDifferentialGood_of_affine
VOA.ConcreteHodge.ghostConjRhoD_proved
VOA.ConcreteHodge.ghostConjBE_proved
VOA.ConcreteHodge.repVacuumNormalized_proved
VOA.ConcreteHodge.repCreationAdjoint_proved
VOA.ConcreteHodge.repWordsReal_proved
VOA.ConcreteHodge.leeYang_gradedUnitary_of_affine_levelOne
VOA.ConcreteHodge.endpointKahler_filtered
VOA.ConcreteHodge.concreteKahler_weakGU
VOA.KoszulResidual.kExact_endpoint
VOA.KoszulResidual.koszulResidual_endpoint
VOA.KoszulResidual.leeYang_gradedUnitary_of_affine_strict
VOA.SymbolModel.jetCyclicDual
VOA.SymbolModel.jet_cyclic
VOA.LeadingKoszul.LeadingBicomplex.equivKer
VOA.InvariantSymbol.finrank_le_gapTwoCoefficient_of_cyclicEvaluation
VOA.LeeYangCharacter.finrank_pieceIn_le
VOA.LeeYangCharacter.goodFil_leeYang_eq_gapSpan
VOA.LeeYangCharacter.leeYangPBWCharacter_iff_lowerBound
VOA.AffineSl2BoundaryCharacter.character_eq_product
VOA.Sl2SemiInfinite.semiInfiniteQ_sq_zero_iff
VOA.Sl2SemiInfinite.ghostAffineAction
VOA.InverseDS.GradedIsometricEmbedding.gradedUnitary_of_filteredIso
VOA.TwoSeries.leeYang_gradedUnitary_of_affine_viaTwoSeries
VOA.TwoSeries.leeYang_gradedUnitary_of_affine_viaInvariants
VOA.TwoSeries.JetPoincareTwo.twoSeries_gradedUnitary_of_affine_final
\end{plaincode}
(The namespace \lean{VOA.InverseDS.GradedIsometricEmbedding} of \lean{gradedUnitary_of_filteredIso} is a historical name; the file
\leanpath{LeeYangCosetTransferFiltered.lean} belongs to this chapter.) The same run printed the type of the headline as
\begin{plaincode}
VOA.AdamovicSymbol.leeYang_gradedUnitary_of_affine : VOA.GradedUnitary VOA.LeeYangTransferAssembly.affineForm
    VOA.LeeYangTransferAssembly.affineRealStructure VOA.LeeYangTransferAssembly.affineFlag →
  VOA.VirGradedUnitarity.VirGradedUnitary VOA.LeeYang.leeYangC VOA.LeeYangTransferAssembly.leeYangRealStructure
\end{plaincode}
so the hypothesis \lean{hA} is its only argument. The build products used are those already in the checkout; they were not rebuilt for this
edition, so the result rests on the project's last aggregate build.

\textbf{Quotations.} Every \texttt{leancode} block of this chapter was compared, piece by piece with comments and whitespace normalized,
with the current Lean sources; all occur verbatim.

%======================================================================
\section{Findings, oddities and Lean defects}\label{l28:sec:defects}
%======================================================================

No statement-level defect was found: the headline says what the mathematics says, about the intended concrete objects, with its one hypothesis
visible. The defects are documentation defects, collected here once. Severities: S1, a statement or definition is wrong or vacuous; S2, a definition is weaker or stronger than its name or the literature, or a hypothesis is redundant; S3, an advertised result is only conditional (a gap); S4, documentation (docstrings, notes) is wrong or stale.

\begin{warning}[title={Documentation defects (Lean docstrings and project notes)}]
\begin{enumerate}
\item \textbf{(Docstrings, S4) ``OPEN'' on proved propositions.} The docstrings of the following \lean{Prop}s or structures call them open, although each is
  proved for the objects used by the headline: \lean{LeeYangRRLowerBound} (\leanpath{LeeYangCharacterReduction.lean}, ``OPEN (classical, not formalized)'');
  \lean{LeeYangRochaCaridiLowerBound} (\leanpath{RRLowerReduction.lean}); \lean{LeeYangGramRankBound} and the module docstring of
  \leanpath{RRLowerGram.lean}; \lean{LeeYangPBWCharacter} (\leanpath{LeeYangTransferAssemblyCharacter.lean}); \lean{LeadingInvariantSymbolDimension} and
  several \lean{OpenInputs} fields (\leanpath{LeeYangTransferAssemblyInputs.lean}); \lean{RelativeDifferentialGood} (\leanpath{ConcreteGoodInterpreted.lean});
  \lean{GhostConjRhoD}, \lean{GhostConjBE} (\leanpath{GhostConjEndpoint.lean}); \lean{RepVacuumNormalized}, \lean{RepCreationAdjoint}, \lean{RepWordsReal}
  (\leanpath{HarmonicMetricBridge.lean}); \lean{HarmonicVacuumNormalized}, \lean{HarmonicCreationAdjoint}, \lean{HarmonicWordsReal}
  (\leanpath{LeeYangShortcutAssembly.lean}); \lean{JetProductBound} and the module docstring (\leanpath{JetStrictEndpoint.lean}); \lean{JetLeadingSymbol},
  \lean{JetCharacterBound} (\leanpath{JetStrictSqueeze.lean}); \lean{GrJetStrict} (\leanpath{SymbolModelJetCyclic.lean}); \lean{WordSymbolDescends} and the
  module docstring (\leanpath{AdamovicSymbolDescent.lean}); \lean{LevelOneBottomDetected}, \lean{LevelOneHeisenbergSingular}
  (\leanpath{CyclicEvalHeisenberg.lean}).
\item \textbf{(Docstrings, S4) rung docstrings list proved inputs as open.} \lean{..._wordSymbols_two} (\leanpath{AdamovicSymbolDegreeOne.lean}),
  \lean{..._wordSymbols} (\leanpath{AdamovicSymbolDescent.lean}), \lean{..._wordSymbols_final} (\leanpath{LYGramEndpoint.lean}) and \lean{..._wordSymbols_rc}
  (\leanpath{RRLowerEndpoint.lean}) mark \lean{desc}, \lean{chi}, \lean{lowerRC} ``OPEN''; all are proved. The module docstrings of
  \leanpath{LeadingKoszulKahler.lean} (``remain to be supplied by a concrete instantiation'', with a final \lean{OPEN} comment) and of
  \leanpath{InvariantSymbolCyclic.lean} (the physical inputs ``are not formalized here'') are superseded by \leanpath{KoszulConcrete*}, \leanpath{KoszulResidual*}
  and \leanpath{CyclicEval*}.
\item \textbf{(Name clash, S4)} Warning~\ref{l28:warn:clash}.
\item \textbf{(Project notes, S4)} Project records written before \lean{VOA.pbw} became a theorem give the headline's axioms as
  including \lean{VOA.pbw}; that was correct then and is obsolete now (Section~\ref{l28:sec:verification}).
\end{enumerate}
\end{warning}

%======================================================================
\section{Referee's checklist}\label{l28:sec:checklist}
%======================================================================

\begingroup\small
\begin{longtable}{p{0.30\textwidth}p{0.30\textwidth}p{0.22\textwidth}p{0.10\textwidth}}
\toprule
Lean declaration & Intended meaning & Hypotheses / status & Fidelity check\\
\midrule\endhead
\lean{AVOA}, \lean{BVOA}, \lean{UVOA}, \lean{ALevel} & $L_{-4/3}$, $L_1$, $V^{-11/3}$ & definitions & Def.~\ref{l28:def:factors}\\
\lean{correctedWordState}, \lean{correctedAffineRPiece}, \lean{affineFlag} & corrected flag & definitions & Def.~\ref{l28:def:corrected}\\
\lean{IsCompactRealStructure}, \lean{compactRealStructure}, \lean{affineRealStructure} & compact conjugation & constructed & Def.~\ref{l28:def:compact}\\
\lean{affineForm}, \lean{leeYangRealStructure} & the two endpoint data & constructed & Def.~\ref{l28:def:lydatum}\\
\lean{AdamovicSymbol.}\allowbreak\lean{leeYang_}\allowbreak\lean{gradedUnitary_}\allowbreak\lean{of_}\allowbreak\lean{affine} & $\ltwentyeightGU(L_{-4/3})\Rightarrow\ltwentyeightGU(\text{Lee--Yang})$ & \textbf{conditional on \lean{hA}} (open) & Thm.~\ref{l28:thm:headline}\\
\lean{LeeYangTransferAssembly.}\allowbreak\lean{leeYang_}\allowbreak\lean{gradedUnitary_}\allowbreak\lean{of_}\allowbreak\lean{affine} & superseded assembly & \lean{hA}, \lean{OpenInputs} (coset identification) & Warn.~\ref{l28:warn:clash}\\
\lean{AffineSl2Action}, \lean{TranslatedAffineSl2Action} & affine action of level $k$ & definitions & Def.~\ref{l28:def:action}\\
\lean{brstState}, \lean{semiInfiniteQ} & BRST current and charge & definitions & Def.~\ref{l28:def:charge}\\
\lean{ghostAffineAction}, \lean{semiInfiniteQ_}\allowbreak\lean{sq_}\allowbreak\lean{zero_}\allowbreak\lean{iff} & ghosts at level 4; $Q^2=0\iff k=-4$ & structural hypotheses; proved & Thm.~\ref{l28:thm:qsq}\\
\lean{relativeCarrier} & relative states & definition & Def.~\ref{l28:def:relative}\\
\lean{leeYangBRSTComplex} & the concrete complex & unconditional & Thm.~\ref{l28:thm:complex}\\
\lean{cosetSubspace}, \lean{cosetOmega} & commutant, coset stress tensor & definitions & Def.~\ref{l28:def:coset}\\
\lean{cosetVOA}, \lean{cosetCentralCharge} & coset VOA, $c=-22/5$ & unconditional & Thm.~\ref{l28:thm:cosetvoa}\\
\lean{cosetToComplex_}\allowbreak\lean{closed}, \lean{leeYangCosetToH0} & coset $\to H^0$ & unconditional & Thm.~\ref{l28:thm:cosetmap}\\
\lean{virasoroToH0_}\allowbreak\lean{ker_}\allowbreak\lean{le_}\allowbreak\lean{maxIdeal} & Virasoro image surjects onto $L$ & unconditional & Thm.~\ref{l28:thm:virimage}\\
\lean{KahlerComplex} & Beem--Garner datum & definition & Def.~\ref{l28:def:kahler}\\
\lean{mem_harm_iff} & harmonic $\iff$ four pieces vanish & unconditional & Thm.~\ref{l28:thm:harm}\\
\lean{harmToLead_}\allowbreak\lean{injective}, \lean{finrank_}\allowbreak\lean{cohomFilt_}\allowbreak\lean{row_}\allowbreak\lean{le_}\allowbreak\lean{lead} & harmonic bound & unconditional & Thm.~\ref{l28:thm:harmbound}\\
\lean{GUData} & form and conjugation on $C$ & definition & Def.~\ref{l28:def:gudata}\\
\lean{flag_nondegenerate}, \lean{layer_eq}, \lean{layer_signed_pos} & Beem--Garner Thm.~4.1 & unconditional & Thm.~\ref{l28:thm:bg}\\
\lean{maps_eq_of_finrank_squeeze} & squeeze & unconditional & Thm.~\ref{l28:thm:squeeze}\\
\lean{gradedUnitary_of_filteredIso} & transfer by filtered isometry & unconditional & Thm.~\ref{l28:thm:filterediso}\\
\lean{gradedUnitaryAt_}\allowbreak\lean{of_}\allowbreak\lean{class_}\allowbreak\lean{family} & GU at one weight from classes & explicit hypotheses & Thm.~\ref{l28:thm:classfamily}\\
\lean{Unitary} & unitarity & definition & Def.~\ref{l28:def:unitary}\\
\lean{LevelOneUnitary.bUnitary}, \lean{bConformal_eq} & $L_1$ unitary; $\omega=\frac14h_{-1}^2$ & \textbf{unconditional} & Thm.~\ref{l28:thm:bunitary}\\
\lean{universal_}\allowbreak\lean{gradedUnitary_}\allowbreak\lean{belowCritical_}\allowbreak\lean{unconditional}, \lean{auxiliaryGradedUnitary} & $\ltwentyeightGU(V^k)$, $k<-2$ & \textbf{unconditional} & Thm.~\ref{l28:thm:aux}\\
\lean{RelativeDifferentialGood}, \lean{relativeDifferentialGood_}\allowbreak\lean{of_}\allowbreak\lean{affine} & goodness & \lean{hA}, \lean{hB} & Def.~\ref{l28:def:good}, Thm.~\ref{l28:thm:good}\\
\lean{GhostConjRhoD}, \lean{GhostConjBE}, \lean{ghostConj*_proved} & ghost conjugation identities & unconditional & Def.~\ref{l28:def:ghostconj}, Thm.~\ref{l28:thm:ghostconj}\\
\lean{Rep*}, \lean{rep*_proved} & metric on representatives & unconditional & Thm.~\ref{l28:thm:repmetric}\\
\lean{leadCohomEquivInvariants}, \lean{finrank_harmDeg_le_invariants} & $H(Q^-)\cong$ invariants & bicomplex data & Thm.~\ref{l28:thm:koszul}\\
\lean{kExact_endpoint}, \lean{koszulResidual_endpoint} & Koszul step, actual complex & \lean{hA}; \lean{hp} (from \lean{hA}) & Thm.~\ref{l28:thm:kexact}\\
\lean{gapTwoCoefficient}, \lean{LeeYangPBWCharacter}, \lean{LeadingInvariantSymbol*} & $G_r(n)$ and the statements & definitions & Def.~\ref{l28:def:gaptwo}\\
\lean{leeYangPBWCharacter_iff_}\allowbreak\lean{lowerBound} & PBW character $\iff$ lower bound; so \lean{LeeYangPBWCharacter} holds & \textbf{unconditional} & Thm.~\ref{l28:thm:pbwchar}\\
\lean{endpointKahler_}\allowbreak\lean{filtered} & step (b): the image is filtered & \lean{hA} & Thm.~\ref{l28:thm:filtered}\\
\lean{finrank_le_gapTwoCoefficient_}\allowbreak\lean{of_}\allowbreak\lean{cyclicEvaluation} & cyclic-evaluation bound & unconditional & Thm.~\ref{l28:thm:cyclic}\\
\lean{jetImage}, \lean{jetDeg}, \lean{jetDer} & jet ring, $A_r$ & definitions & Def.~\ref{l28:def:jet}\\
\lean{wordSym}, \lean{WordSymbolDescends}, \lean{JetLeadingSymbol}, \lean{GrJetStrict} & symbols, descent, strictness & definitions & Def.~\ref{l28:def:symbols}\\
\lean{jet_cyclic} & cyclicity at $(0,1)$ & unconditional & Thm.~\ref{l28:thm:jetcyclic}\\
\lean{Desc.symbol_eq_zero_of_mem}, \lean{Desc.wordSymbolDescends} & descent, every degree & unconditional & Thm.~\ref{l28:thm:desc}\\
\lean{JetProductBound}, \lean{jetSeries}, \lean{JetFermionicIdentity} & jet bound, fermionic series & definitions & Def.~\ref{l28:def:jetbound}\\
\lean{jetFermionicIdentity}, \lean{coeff_jetSeries_le}, \lean{grJetStrict_of_}\allowbreak\lean{leadingSymbol_product} & fermionic identity; strictness & \textbf{unconditional}; (3) from inputs & Thm.~\ref{l28:thm:jetfermionic}\\
\lean{LeeYangRRLowerBound}, \lean{rc}, \lean{LeeYangGramRankBound} & character statements & definitions & Def.~\ref{l28:def:rr}\\
\lean{goodFil_leeYang_eq_gapSpan}, \lean{finrank_pieceIn_le} & gap-two spanning & unconditional & Thm.~\ref{l28:thm:gapspan}\\
\lean{sum_gapTwoCoefficient_eq_rc}, \lean{finrank_le_rc} & Rogers--Ramanujan (Andrews) & unconditional & Thm.~\ref{l28:thm:andrews}\\
\lean{pow_dvd_det_of_ker} & corank $\le$ order & unconditional & Lem.~\ref{l28:lem:corank}\\
\lean{leeYangGramRankBound_holds}, \lean{leeYangRochaCaridiCharacter_holds}, \lean{leeYangRRLowerBound_holds} & Lee--Yang vacuum character & \textbf{unconditional} & Thm.~\ref{l28:thm:lychar}\\
\bottomrule
\end{longtable}
\endgroup

%======================================================================
\section{Notes and references}\label{l28:sec:notes}
%======================================================================

\textbf{Graded unitarity.} A.~Arabi Ardehali, C.~Beem, M.~Lemos and L.~Rastelli, arXiv:2507.23781 (graded unitarity of chiral algebras of
four-dimensional $\mathcal N=2$ theories; the affine vacuum determinant used for the auxiliary factor is their equation (4.13)); C.~Beem and A.~Garner,
``Graded unitarity and relative semi-infinite cohomology'', arXiv:2509.10364 (Lemma~3.16, Proposition~3.24, Theorem~4.1, Remark~4.4 are formalized in
Section~\ref{l28:sec:hodge}); C.~Dong and X.~Lin, ``Unitary vertex operator algebras'', \emph{J.\ Algebra} 397 (2014), 252--277; H.~Kaufmann's conjecture
is discussed in Chapter~\ref{ch:L19}.

\textbf{Chiral algebras and Argyres--Douglas theories.} C.~Beem, M.~Lemos, P.~Liendo, W.~Peelaers, L.~Rastelli and B.~C.~van~Rees, ``Infinite chiral symmetry
in four dimensions'', \emph{Comm.\ Math.\ Phys.} 336 (2015), 1359--1433, arXiv:1312.5344 (the chiral-algebra map and gauging as BRST reduction); C.~Córdova
and S.-H.~Shao, ``Schur indices, BPS particles, and Argyres--Douglas theories'', \emph{JHEP} 2016(1), 040, arXiv:1506.00265; T.~Creutzig, ``W-algebras for
Argyres--Douglas theories'', arXiv:1701.05926 ($(A_1,D_{2n+1})$ and $L_{-4n/(2n+1)}(\ltwentyeightsl)$).

\textbf{Cosets and semi-infinite cohomology.} P.~Goddard, A.~Kent and D.~Olive, ``Virasoro algebras and coset space models'', \emph{Phys.\ Lett.\ B} 152
(1985), 88--92, and ``Unitary representations of the Virasoro and super-Virasoro algebras'', \emph{Comm.\ Math.\ Phys.} 103 (1986), 105--119;
T.~Arakawa, T.~Creutzig and A.~R.~Linshaw, ``W-algebras as coset vertex algebras'', \emph{Invent.\ Math.} 218 (2019), 145--195, arXiv:1801.03822 (the coset
identification that the Lean proof avoids); B.~L.~Feigin, ``The semi-infinite homology of Kac--Moody and Virasoro Lie algebras'', \emph{Russian Math.\ Surveys}
39 (1984), 155--156; I.~B.~Frenkel, H.~Garland and G.~J.~Zuckerman, ``Semi-infinite cohomology and string theory'', \emph{Proc.\ Natl.\ Acad.\ Sci.\ USA} 83
(1986), 8442--8446.

\textbf{The affine and Virasoro inputs.} D.~Adamović, ``A realization of certain modules for the $N=4$ superconformal algebra and the affine Lie algebra
$A_1^{(1)}$'', \emph{Transform.\ Groups} 21 (2016), 299--327, arXiv:1407.1527 (the free-field realization behind the symbols $e\mapsto u^2$, \ldots); T.~Arakawa and
J.~Dai, the boundary-level character and classical freeness (Chapter~\ref{ch:L20}); A.~Rocha-Caridi, ``Vacuum vector representations of the Virasoro algebra'', in
\emph{Vertex Operators in Mathematics and Physics}, MSRI Publ.\ 3 (1985), 451--473; B.~L.~Feigin and D.~B.~Fuchs, ``Verma modules over the Virasoro algebra'',
Lecture Notes in Math.\ 1060 (1984), 230--245; V.~G.~Kac, ``Contravariant form for infinite-dimensional Lie algebras and superalgebras'', Lecture Notes in Phys.\ 94
(1979), 441--445 (the Kac determinant); G.~E.~Andrews, ``A polynomial identity which implies the Rogers--Ramanujan identities'', \emph{Scripta Math.} 28 (1970),
297--305; S.~Chern, D.~Tran and A.~Wakhare, arXiv:2605.02251, Theorem~1.2 (the fermionic identity, at $k=1$, $t=q$, $z=1$, combined with Jacobi's cube identity).

\endgroup
